\documentclass[11pt,a4paper]{report}

% ==========================================
% NEW ELEGANT PREAMBLE
% ==========================================
\usepackage[utf8]{inputenc}
\usepackage[T1]{fontenc}
\usepackage{amsmath, amssymb, amsthm, mathtools, environ}

% --- FONT CONFIGURATION ---
% Elegant Palatino font (Currently active)
\usepackage{mathpazo} 

% Alternative Font Options (Uncomment one to try):
% \usepackage{lmodern} % Latin Modern: The clean, standard LaTeX look.
% \usepackage{libertine} \usepackage{libertinust1math} % Linux Libertine: Very elegant, classic book feel.
% \usepackage{fourier} % Utopia: Robust, slightly heavier strokes, fantastic for complex math.
% \usepackage{stix2} % STIX2: Times-based, incredibly comprehensive math symbol support.

% --- EMPHASIS STYLE OVERRIDE ---
\DeclareTextFontCommand{\emph}{\bfseries}

\usepackage[margin=1in, headheight=35pt]{geometry}
\usepackage[dvipsnames]{xcolor}
\usepackage{hyperref}

% --- HYPERREF AESTHETICS ---
\hypersetup{
    colorlinks=true,
    linkcolor=MidnightBlue,
    filecolor=Magenta,      
    urlcolor=MidnightBlue,
    citecolor=MidnightBlue,
    bookmarksopen=true
}

\usepackage{tcolorbox}
\tcbuselibrary{skins, breakable, theorems}

% --- TYPOGRAPHY & LAYOUT ENHANCEMENTS ---
\usepackage{parskip} % Modern paragraph spacing (no indents, clean vertical space)
\linespread{1.05}    % Palatino benefits immensely from slightly increased line spacing

% --- LIST SPACING FIX ---
\usepackage{enumitem}
\setlist{
    topsep=0.5ex,
    itemsep=0.5ex,
    parsep=\parskip,
    partopsep=0pt
}

% Make description labels flow like normal inline text with stretchable spaces
\setlist[description]{
    style=unboxed,
    labelsep=0.333em plus 0.167em minus 0.111em
}

\usepackage{fancyhdr}
\usepackage{extramarks}
\pagestyle{fancy}
\fancyhf{} % clear all header and footer fields
\fancyhead[R]{\sffamily\color{Gray}\thepage}
\fancyhead[L]{\sffamily\color{Gray}\leftmark}
\fancyfoot[L]{\sffamily\color{Gray}\rightmark}
\fancyfoot[R]{\sffamily\color{Gray}\lastxmark}
\renewcommand{\headrule}{\hbox to\headwidth{\color{Gray!50}\leaders\hrule height 0.5pt\hfill}}
\renewcommand{\footrule}{\hbox to\headwidth{\color{Gray!50}\leaders\hrule height 0.5pt\hfill}}

% Redefine marks to capture both section and subsection
\renewcommand{\sectionmark}[1]{%
    \markright{\MakeUppercase{\thesection.\ #1}}%
    \extramarks{}{}% Clear the subsection mark when a new section starts
}
\renewcommand{\subsectionmark}[1]{%
    \extramarks{\MakeUppercase{\thesubsection.\ #1}}{\MakeUppercase{\thesubsection.\ #1}}%
}

% Fail-Safe TikZ Configuration
\usepackage{tikz}
\usetikzlibrary{calc, arrows.meta, positioning, 3d, shapes.geometric, decorations.pathreplacing, patterns}
\usepackage{pgfplots}
\pgfplotsset{compat=1.18}

% ==========================================
% CUSTOM CHAPTER & SECTION FORMATTING
% ==========================================
\usepackage{titlesec}

\definecolor{bglight}{HTML}{F8F9FA}    

% Math operators
\DeclareMathOperator{\id}{id}
\DeclareMathOperator{\Vol}{Vol}

% Custom Quotes Macro
\newcommand{\qt}[1]{\textit{``#1''}}

% Custom Inline Meta-Note for stage directions (V1.16.2+)
\newcommand{\inlinemetanote}[1]{\textit{[\textcolor{Gray}{#1}]}}

% Standard Math Environments
\newtheorem{theorem}{Theorem}[chapter]
\newtheorem{lemma}[theorem]{Lemma}
\theoremstyle{definition}
\newtheorem{definition}[theorem]{Definition}

% --- NUMBEREd ENVIRONMENTS ---
%\newtheorem{lemma}[theorem]{Lemma}
\newtheorem{corollary}[theorem]{Corollary}
%\newtheorem{definition}[theorem]{Definition}
\newtheorem{proposition}[theorem]{Proposition}
\newtheorem{example}[theorem]{Example}
\newtheorem{exercise}[theorem]{Exercise}

% --- UNNUMBERED ENVIRONMENTS ---
% The asterisk (*) prevents them from being numbered!
\newtheorem*{theorem*}{Theorem}
\newtheorem*{lemma*}{Lemma}
\newtheorem*{proposition*}{Proposition}
\newtheorem*{definition*}{Definition}
\newtheorem*{corollary*}{Corollary}
\newtheorem*{claim*}{Claim}
\newtheorem*{example*}{Example}
\newtheorem*{exercise*}{Exercise}

\newtheorem*{notation}{Notation}
\newtheorem*{remark}{Remark}
\newtheorem*{summary}{Summary}
\newtheorem*{warmup}{Warm up}
\newtheorem*{question}{Question}
\newtheorem*{answer}{Answer}
\newtheorem*{importantremark}{Important remark}
\newtheorem*{goals}{Goals}

% ==========================================
% REDESIGNED META-ENVIRONMENTS
% ==========================================

% 1. Clean Spoken Protocol
\newtcolorbox{spoken-clean}[1][]{
    enhanced,
    frame hidden,
    colback=MidnightBlue!5!white,
    borderline west={3pt}{0pt}{MidnightBlue},
    title={\sffamily\bfseries\textcolor{MidnightBlue}{Spoken Protocol [#1]}},
    attach title to upper=\par\vspace{1ex},
    breakable,
    %fontupper=\itshape
}


% 2. AI Stroke-by-Stroke Note
\newtcolorbox{nice-box}[1][]{
    enhanced,
    colback=bglight,
    colframe=gray!40,
    boxrule=0.5pt,
    drop fuzzy shadow,
    title={\sffamily\bfseries #1},
    coltitle=black,
    attach boxed title to top left={xshift=10pt, yshift=-10pt},
    boxed title style={colback=gray!20, colframe=white, boxrule=0pt},
    top=14pt,
    breakable
}

% 3. Best Teacher / Didactic Insight
\newtcolorbox{didactic-insight}[1][Pedagogical Concept]{
    enhanced,
    arc=3mm,
    colback=ForestGreen!5!white,
    colframe=ForestGreen!80!black,
    title={\sffamily\bfseries Insight: #1},
    frame hidden,
    colback=ForestGreen!3!white,
    borderline west={2pt}{0pt}{ForestGreen!80!black},
    title={\sffamily\bfseries\textcolor{ForestGreen!80!black}{\textit{Insight:} #1}},
    attach title to upper=\par\vspace{0.5ex},
    breakable
}

% 4a. Modular Visual Blackboard Replication Box (V1.17)
\NewTColorBox{color-box}{ m O{} }{
    enhanced,
    colframe=#1,
    colback=white,
    colupper=black,
    fonttitle=\bfseries\sffamily,
    coltitle=black,
    title={#2},
    boxrule=1.5pt,
    arc=4pt,
    drop fuzzy shadow=#1!30,
    attach boxed title to top left={xshift=10pt, yshift=-10pt},
    boxed title style={colback=white, colframe=#1, boxrule=0.5pt, top=2pt, bottom=2pt, left=4pt, right=4pt},
    % Hide title area completely if the optional argument is not provided
    IfValueF={#2}{notitle}
}

% 4b. Legacy Visual Blackboard Replication Boxes (Backward Compatibility)
\newcommand{\NewChalkBox}[3]{%
    \newtcolorbox{#1}[1][Highlighted Board Formula]{%
        enhanced,
        colback=white,
        colframe=#2,
        fonttitle=\bfseries\sffamily,
        coltitle=black,
        title=##1,
        boxrule=1.5pt,
        arc=4pt,
        drop fuzzy shadow=#3,
        attach boxed title to top left={xshift=10pt, yshift=-10pt},
        boxed title style={colback=white, colframe=#2, boxrule=0.5pt, top=2pt, bottom=2pt, left=4pt, right=4pt}
    }%
}

\NewChalkBox{orange-box}{BurntOrange}{BurntOrange!30}
\NewChalkBox{blue-box}{MidnightBlue}{MidnightBlue!30}
\NewChalkBox{dark-green-box}{ForestGreen}{ForestGreen!30}
\NewChalkBox{apple-green-box}{LimeGreen}{LimeGreen!30}
\NewChalkBox{yellow-box}{Goldenrod!90!black}{Goldenrod!30}
\NewChalkBox{violet-box}{Violet}{Violet!30}
\NewChalkBox{red-box}{BrickRed}{BrickRed!30}
\NewChalkBox{green-box}{Green}{Green!30}
\NewChalkBox{purple-box}{Purple}{Purple!30}
\NewChalkBox{gray-box}{Gray}{Gray!30}
\NewChalkBox{apricot-box}{Apricot}{Apricot!30}
\NewChalkBox{bittersweet-box}{Bittersweet}{Bittersweet!30}
\NewChalkBox{teal-blue-box}{TealBlue}{TealBlue!30}
\NewChalkBox{cyan-box}{Cyan}{Cyan!30}
\NewChalkBox{magenta-box}{Magenta}{Magenta!30}
\NewChalkBox{brown-box}{Brown}{Brown!30}
\NewChalkBox{pink-box}{Pink}{Pink!30}

% ==========================================
% LEGACY FORMULA BOXES (Backward Compatibility)
% ==========================================
\NewChalkBox{orange-formula}{BurntOrange}{BurntOrange!30}
\NewChalkBox{blue-formula}{MidnightBlue}{MidnightBlue!30}
\NewChalkBox{dark-green-formula}{ForestGreen}{ForestGreen!30}
\NewChalkBox{apple-green-formula}{LimeGreen}{LimeGreen!30}
\NewChalkBox{yellow-formula}{Goldenrod!90!black}{Goldenrod!30}
\NewChalkBox{violet-formula}{Violet}{Violet!30}
\NewChalkBox{red-formula}{BrickRed}{BrickRed!30}

% 5. Mathematical Setup (Stroke-by-Stroke)
\newtcolorbox{math-stroke}[1][Mathematical Setup (Stroke-by-Stroke)]{
    enhanced,
    frame hidden,
    colback=white,
        opacityback=0,
    borderline west={2pt}{0pt}{MidnightBlue!30},
    title={\sffamily\bfseries\textcolor{MidnightBlue!80!black}{#1}},
    attach title to upper=\par\vspace{1ex},
    breakable
}

% 6. Redundant Explanation
\newtcolorbox{redundant-explanation}[1][Redundant Explanation of Step]{
    enhanced,
    arc=2mm,
    colback=Purple!5!white,
    colframe=Purple!60!black,
    title={\sffamily\bfseries #1},
    breakable
}

% 7. Meta-Note
\newtcolorbox{meta-note}[1][Administrative Meta-Note]{
    enhanced,
    frame hidden,
    colback=Gray!10,
    borderline={0.5pt}{0pt}{Gray, dashed},
    title={\sffamily\bfseries\textcolor{Gray!90!black}{#1}},
    attach title to upper=\par\vspace{1ex},
    breakable
}


% 8. AI Observational Note
\newtcolorbox{ai-note}[1][AI Note]{
    enhanced,
    colback=Goldenrod!5!white,
    colframe=Goldenrod!90!black,
    fonttitle=\bfseries\sffamily,
    coltitle=black,
    title=#1,
    boxrule=1pt,
    arc=3pt,
    breakable
}

% 9. Student Question, Student Interaction, or Common Confusion.
\newtcolorbox{student-interaction}[1][Student Interaction]{
    enhanced,
    frame hidden,
    colback=Violet!5!white,
    borderline west={3pt}{0pt}{Violet},
    title={\sffamily\bfseries\textcolor{Violet!90!black}{#1}},
    attach title to upper=\par\vspace{1ex},
    fontupper=\itshape,
    left skip=20pt,
    right skip=20pt,
    breakable
}


% 10. Explanation of Steps
\newenvironment{explanation-of-steps}{
    \par\vspace{1ex}\noindent\textbf{Explanation of Steps:}
}{
    \par\vspace{1ex}
}

% 10.5 Short Inline Proof (Safe to nest inside math-stroke)
\newenvironment{short-proof}[1][Proof]{
    \par\vspace{1ex}\noindent\textbf{\sffamily #1:} \itshape
}{
    \hfill$\blacksquare$\par\vspace{1ex}
}

% 11. Proof Wrapper (Overrides default amsthm proof)
\let\proof\relax
\let\endproof\relax
\newtcolorbox{proof}[1][Proof]{
    enhanced,
    breakable,
    frame hidden,
    colback=white,
        opacityback=0,
    borderline west={3pt}{0pt}{Gray!80!black},
    % Creative Title Badge to indicate the start of the proof
    title={\tikz[baseline=(A.base)]\node[fill=Gray!80!black, text=white, rounded corners=2pt, inner xsep=6pt, inner ysep=3pt, font=\sffamily\bfseries\large] (A) {#1};},
    attach title to upper=\par\vspace{1.5ex},
    % Diagonal "PROOF" watermark overlaid across the nested contents
    finish unbroken={%
        \node[rotate=35, scale=5, text=Gray, opacity=0.06, font=\sffamily\bfseries] at (frame.center) {PROOF};
        \node[anchor=south east, xshift=-4pt, yshift=4pt] at (frame.south east) {\tikz\node[draw=Gray!80!black, fill=white, line width=1pt, rounded corners=2pt, inner xsep=4pt, inner ysep=2pt, font=\sffamily\bfseries\small, text=Gray!80!black] {Q.E.D.};};
    },
    finish first={%
        \node[rotate=35, scale=5, text=Gray, opacity=0.06, font=\sffamily\bfseries] at (frame.center) {PROOF};
    },
    finish middle={%
        \node[rotate=35, scale=5, text=Gray, opacity=0.06, font=\sffamily\bfseries] at (frame.center) {PROOF};
    },
    finish last={%
        \node[rotate=35, scale=5, text=Gray, opacity=0.06, font=\sffamily\bfseries] at (frame.center) {PROOF};
        \node[anchor=south east, xshift=-4pt, yshift=4pt] at (frame.south east) {\tikz\node[draw=Gray!80!black, fill=white, line width=1pt, rounded corners=2pt, inner xsep=4pt, inner ysep=2pt, font=\sffamily\bfseries\small, text=Gray!80!black] {Q.E.D.};};
    }
}

% 12. Lecture Break
\newtcolorbox{lecture-break}[1][Lecture Break]{
    enhanced,
    frame hidden,
    colback=white,
    borderline top={0.5pt}{0pt}{Gray, dashed},
    borderline bottom={0.5pt}{0pt}{Gray, dashed},
    title={\sffamily\bfseries\centering\textcolor{Gray!90!black}{#1}},
    attach title to upper=\par\vspace{1ex},
    fontupper=\centering\itshape\color{Gray!90!black},
    breakable
}

% ==========================================
% INVISIBLE LAYER & FALLBACK ENVIRONMENTS (V1.17)
% ==========================================
% These environments act as AI scratchpads and are typically stripped by the C# backend.
% The environ package completely hides them from the compiled PDF if they leak through.
\NewEnviron{ai-tikz-planner-invisible-content}{}
\NewEnviron{ai-type-check-invisible-content}{}
\NewEnviron{ai-proof-skeleton-invisible-content}{}
\NewEnviron{ai-example-invisible-content}[1][]{}
\NewEnviron{ai-discard-invisible-content}[1][]{}
\NewEnviron{ai-retroactive-patch}[1][]{}
\NewEnviron{ai-global-state-checkpoint-invisible-content}{}

% Data Integrity Fallbacks (Rendered as visible warning boxes in case they appear in the source)
\newtcolorbox{ai-logical-gap}[1][Logical Gap]{
    enhanced,
    colback=BrickRed!5,
    colframe=BrickRed,
    title={\sffamily\bfseries\textcolor{white}{#1}},
    fonttitle=\sffamily\bfseries,
    coltitle=BrickRed!75!black,
    title={#1},
    breakable
}
\newtcolorbox{ai-generation-loop-fallback}[1][Generation Loop Fallback]{
    enhanced,
    colback=BrickRed!5,
    colframe=BrickRed,
    title={\sffamily\bfseries\textcolor{white}{#1}},
    fonttitle=\sffamily\bfseries,
    coltitle=BrickRed!75!black,
    title={#1},
    breakable
}
\newtcolorbox{ai-raw-ocr-fallback}[1][Raw OCR Fallback]{
    enhanced,
    colback=BrickRed!5,
    colframe=BrickRed,
    title={\sffamily\bfseries\textcolor{white}{#1}},
    fonttitle=\sffamily\bfseries,
    coltitle=BrickRed!75!black,
    title={#1},
    breakable
}
\newtcolorbox{ai-phonetic-dump}[1][Phonetic Dump]{
    enhanced,
    colback=BrickRed!5,
    colframe=BrickRed,
    title={\sffamily\bfseries\textcolor{white}{#1}},
    fonttitle=\sffamily\bfseries,
    coltitle=BrickRed!75!black,
    title={#1},
    breakable
}
\newtcolorbox{ai-modality-conflict}[1][Modality Conflict]{
    enhanced,
    colback=BrickRed!5,
    colframe=BrickRed,
    title={\sffamily\bfseries\textcolor{white}{#1}},
    fonttitle=\sffamily\bfseries,
    coltitle=BrickRed!75!black,
    title={#1},
    breakable
}
\newtcolorbox{ai-off-camera-state}[1][Off-Camera State]{
    enhanced,
    colback=BrickRed!5,
    colframe=BrickRed,
    title={\sffamily\bfseries\textcolor{white}{#1}},
    fonttitle=\sffamily\bfseries,
    coltitle=BrickRed!75!black,
    title={#1},
    breakable
}
\newtcolorbox{ai-async-board-update}[1][Async Board Update]{
    enhanced,
    colback=BrickRed!5,
    colframe=BrickRed,
    title={\sffamily\bfseries\textcolor{white}{#1}},
    fonttitle=\sffamily\bfseries,
    coltitle=BrickRed!75!black,
    title={#1},
    breakable
}
\newtcolorbox{ai-kinetic-emphasis}[1][Kinetic Emphasis]{
    enhanced,
    colback=BrickRed!5,
    colframe=BrickRed,
    title={\sffamily\bfseries\textcolor{white}{#1}},
    fonttitle=\sffamily\bfseries,
    coltitle=BrickRed!75!black,
    title={#1},
    breakable
}

% Custom Lecture Chapter Command
\newcommand{\lecturechapter}[4]{%
  \chapter[#1, #2: #4]{{\Large\color{Gray}#1: #3} \texorpdfstring{\\[1ex]}{-- } {\LARGE #4}}%
}

% --- Chapter Formatting ---
% Removes the massive default whitespace at the top of the page
\titlespacing*{\chapter}{0pt}{-20pt}{40pt}

% Redesigns the chapter header (Elegant, right-aligned, with stable rules)
\titleformat{\chapter}[display]
  {\normalfont\sffamily\Large\bfseries\color{MidnightBlue!90!black}}
  {\raggedleft\large\color{Gray}\MakeUppercase{\chaptertitlename}\ \thechapter}
  {1ex}
  {\rule{\textwidth}{1.5pt}\vspace{2ex}\raggedleft \break}
  [\vspace{1ex}\rule{\textwidth}{1.5pt}]

% --- Section Formatting ---
\titleformat{\section}
  {\normalfont\sffamily\Large\bfseries\color{MidnightBlue}}
  {\thesection}{1em}{}

\titleformat{\subsection}
  {\normalfont\sffamily\large\bfseries\color{MidnightBlue}}
  {\thesubsection}{1em}{}

\begin{document}

% ==========================================
% AutoExtraction Source: 02-16-monday.mp4
% Model: gemini-3-flash-preview
%
% MERGED AND CLEANED by Gemini Code Assist on 2026-05-01
% This file combines TEIL 1 and TEIL 2 to create a coherent transcript,
% removing redundant overlapping sections and TEIL 3.
% ==========================================

\chapter{Monday: February 16th 2026 \& Real Analysis }


\begin{nice-box}[Context]
This lecture marks the beginning of \emph{Analysis II: mehrere Variablen} (Multivariable Calculus) for the Spring 2026 semester at ETH Zürich. The lecturer, Joaquim Serra, introduces the course structure, logistics, and provides a high-level overview comparing the fundamental goals of Analysis I (single-variable calculus) with the upcoming objectives of Analysis II (vector calculus and higher-dimensional integration).
\end{nice-box}

% ==========================================
% SECTION 1: INTRODUCTION AND LOGISTICS
% ==========================================
\section{Introduction and Administrative Information}

\begin{spoken-clean}[00:00:00 - 00:01:20]
Good morning. Good morning. Thank you. Please, please, please. Uh, good morning and welcome to Analysis 2. Uh, I have here a projection with my name, Joaquim Serra, the name of the coordinator, Enric Florit-Simón. Now, this is my introduction. Um, and then you know the times of the lectures, obviously, because you are here. Um, and I wanted to share before starting with the mathematical content some practical informations that are very important because we are a large group and the — the communication and everything needs to be done in a more or less structured way. Otherwise, otherwise we collapse.
\end{spoken-clean}

\begin{meta-note}[Projected Content]
The lecturer displays the course website for \qt{Analysis II: mehrere Variablen Spring 2026}. Key details shown:
\begin{description}
    \item[Lecturer:] Joaquim Serra
    \item[Coordinator:] Enric Florit-Simón
    \item[Lectures:] Mon 08:15--10:00 (ETA F 5), Wed 08:15--10:00 (HG F 7 / HG F 5), Thu 16:15--18:00 (ETA F 5).
\end{description}
\end{meta-note}

\begin{spoken-clean}[00:01:20 - 00:03:30]
So, as I guess you did — so as a general rule, uh, everything that applied to Analysis 1 applies to Analysis 2, unless otherwise stated. We will try to follow the same rules and then you know what to expect. In particular, you will have this, um, metadata page that I can see that from — from up there is a bit difficult to see, right? But I will try to put it bigger, yeah. Um, so you know that you need to register for the exercise classes. Uh, for the communication, please try to use whenever possible the forum. We will look at the forum and — and reply, and also your — your peers can also reply to questions because many — so then we can use the collective knowledge, right? Um, there will be recordings of the lecture, but not until the later on, okay? Because, uh, is important for the lecture and for you, but mainly, I mean, for you is your interest, uh, if you want to come or not. But for the lecture is very good that you come to the lecture, you ask questions, you stop me when I'm going too fast, these kind of things. You know, it's very important. So it's a — and also is where you meet with your peers and you can communicate, ask questions to your peers. It's — it's — it's better for the learning.
\end{spoken-clean}

\begin{spoken-clean}[00:03:30 - 00:05:07]
Um, we will tell you more about the exam later on, right? Today I will tell you a bit more what is important to be successful in the course, what to expect about this course, and, uh, the details and the rules of the exam will come later on. You don't need to know now. The bonus, there is a $0.25$ bonus, will be in questions that you will answer in some exercise classes. We will tell you in advance which ones. And, uh, essentially you need to be in the — in the exercise class and answer it. And, uh, do not worry because, uh, the — the rule will be thought — in any case it's just $0.25$ points — uh, but the rules will be thought so that if you miss one day or casuistics like this, it's not that you get zero points, right? Okay. Um, lecture notes. There will be lecture notes similar to those in Analysis 1, and we will put them in this page, hopefully by the end of each week, because I — I update them after the lecture, mainly depending on your questions, depending on your feedback. And that's it, uh, I think that the rules are mainly those and the — the organization. Is there some question, uh, about organization? So essentially, as I said, uh, if you have doubts, uh, is like in Analysis 1, unless I say something that is surprising to you, as in Analysis 1. Is it okay, this? Yeah. Okay, now I will switch — I have some other things I want to project later, but now, because we cannot record at the same time blackboard and projection, I will switch one moment to the blackboard. \inlinemetanote{The lecturer turns off the projector and prepares the blackboard} So I can mute, I think. Yeah. Okay. Now we need the light. \inlinemetanote{Adjusts classroom lighting} Okay. Ah, I think that there is already on. Okay.
\end{spoken-clean}

% ==========================================
% SECTION 2: OVERVIEW - ANALYSIS I VS II
% ==========================================
\section{Overview: Analysis I vs. Analysis II}

\begin{spoken-clean}[00:05:07 - 00:07:20]
Okay, now, just to start, um, we will give a very quick, so quick that is almost a cartoon, but a very quick overview of what is Analysis 1 and Analysis 2, right? It's very, very quick. \inlinemetanote{Starts writing on the left blackboard} So, Analysis 1. Um, is — sorry, but for me is very difficult to write big and then it's — so — so if — if you don't see something, I will try to say at loud everything I write, but if — if something — I mean, Analysis is... so, um, Analysis 1, calculus, sorry, it's the first day, right? Bear with me. Calculus is written — is spelled correctly? Yeah. Calculus of one variable, right? And then, if I ne — if — if we had to choose a representative theorem of Analysis 1 that a bit represents the course, uh, what would you choose? So, can you tell me — I know it's the first day, but let's start a bit of communication. Can you tell me at least two answers — so two people can tell me what is the representative — one representative theorem of Analysis 1? Yeah?
\end{spoken-clean}

\begin{student-question}
Fundamental Theorem of Calculus.
\end{student-question}

\begin{spoken-clean}[continued]
Yeah. I agree. Some other representative theorem? He said the Fundamental Theorem of Calculus. Makes sense, right? Okay, after — is there some other guess? Is — yeah, there?
\end{spoken-clean}

\begin{student-question}
The Intermediate Value Theorem.
\end{student-question}

\begin{spoken-clean}[continued]
The Intermediate Value Theorem. Okay. I think we need to use this for the recording, otherwise I — I will repeat. \inlinemetanote{Referring to the catch-box microphone} Later on when I get — I get better with the throwing the — we will use that. So, okay, so fundamental — so we said two theorems that are — that are indeed important. Perhaps the most representative, I would agree and — and I saw many people saying yes, so is the Fundamental Theorem of Calculus, right? And is this... \inlinemetanote{Starts writing the formula on the board}
\end{spoken-clean}

\begin{nice-box}[Board Action: The Fundamental Theorem of Calculus]
The lecturer writes the Fundamental Theorem of Calculus (FTC) on the board. He initially makes a common sign error in the subtraction, which is subsequently corrected.
\end{nice-box}

\begin{math-stroke}[The Fundamental Theorem of Calculus (Recall)]
\begin{ai-type-check-invisible-content}
% f: [a, b] \to \mathbb{R} is differentiable.
% f' is Riemann integrable.
\end{ai-type-check-invisible-content}
In Analysis I, the central result is the Fundamental Theorem of Calculus (FTC), which relates the integral of a derivative to the values of the function at the boundaries:
\[ \text{FTC} \qquad \int_{a}^{b} f'(x) \, dx = \textcolor{red}{f(a) - f(b)} \quad \text{\textcolor{red}{(Incorrect)}} \]
\end{math-stroke}

\begin{spoken-clean}[00:08:35 - 00:08:50]
Is it correct? Yeah, looks like... \inlinemetanote{Pauses and looks at the board} Yeah, there is some mistake? \inlinemetanote{Corrects the board} Good, well spotted.
\end{spoken-clean}

\begin{math-stroke}[The Fundamental Theorem of Calculus (Corrected)]
The corrected statement of the theorem:
\[ \text{FTC} \qquad \int_{a}^{b} f'(x) \, dx = f(b) - f(a) \]
\end{math-stroke}

\begin{spoken-clean}[00:08:50 - 00:10:32]
Now, to — to unfold and to write and to understand every detail about this equation over here, you needed to introduce a lot of things. You needed to define what is that (i.e., the derivative $f'$), right? What is this (i.e., the Riemann integral $\int_{a}^{b}$)? And even — and you spent a lot of time, right? — what is that (i.e., the real numbers $\mathbb{R}$)? It was not — it was not easy. Yeah. Okay, so this is a fairly, fairly compact theorem, yet the ingredients you had to learn, uh, to unfold all of this and to understand the meaning of every symbol here were a bit of work. Now, how will Analysis 2 look like? \inlinemetanote{Moves to   the next section of the board} So this will be calculus, calculus, still calculus in, uh, more than one variable. Well, since you already know one, let's say at least two variables, because otherwise is already covered by Analysis 1. And some representative theorem that would be the equivalent — so the counterpart of this one (i.e., the FTC) — is the Gauss — it has many names, and is always the same theorem. So but — so is the Gauss, Stokes, uh, what else? Divergence, sometimes is called the Divergence Theorem. Uh, and — and that's it, I think. Green, sometimes Green Theorem.
\end{spoken-clean}

\begin{math-stroke}[Analysis II: The Goal]
Analysis II extends these concepts to higher dimensions:
\begin{description}
    \item[Analysis 1:] Calculus of one variable.
    \item[Analysis 2:] Calculus in $\ge 2$ variables.
\end{description}
The representative results in Analysis II are a family of related theorems:
\begin{itemize}
    \item Gauss Theorem
    \item Stokes Theorem
    \item Divergence Theorem
    \item Green Theorem
\end{itemize}
\end{math-stroke}

\begin{spoken-clean}[00:10:32 - 00:11:00]
And now I will write the theorem, and you will not understand exactly what each thing means. I will tell you intuitively, and then the goal of the lecture will be at the end of the lecture you will be able to understand in full precision everything that I wrote here. Okay? So the theorem goes as follows. So we have — we have the space, the space, let's say I will put it in the space $\mathbb{R}^3$. Three coordinates of space, right? This thing... \inlinemetanote{Draws a 3D coordinate system on the board}
\end{spoken-clean}

\begin{math-stroke}[Visualizing the Higher-Dimensional Space]
\begin{center}
\begin{tikzpicture}[scale=1.2, node distance=0.5cm]
% \begin{ai-tikz-planner-invisible-content}
% 1. Background: 3D Axes (x, y, z).
% 2. Midground: A generic 3D volume Omega (the "potato").
% 3. Foreground: Labels for the axes and the volume.
% \end{ai-tikz-planner-invisible-content}
    % 3D Axes
    \draw[->, thick] (0,0,0) -- (2,0,0) node[right] {$x$ (or $x_1$)};
    \draw[->, thick] (0,0,0) -- (0,2,0) node[above] {$z$ (or $x_3$)};
    \draw[->, thick] (0,0,0) -- (0,0,2) node[below left] {$y$ (or $x_2$)};
    
    \node[below right] at (2.5, 2.5, 0) {$\mathbb{R}^3$};
\end{tikzpicture}
\end{center}
\end{math-stroke}

% [SYSTEM] Segment complete. Please prompt "Continue" for the remainder of the segment.
\begin{spoken-clean}[00:11:00 - 00:12:28]
Right? And then in the — in this space (i.e., $\mathbb{R}^3$) we are given a domain, we are given some — or some surface that encloses some interior — so this — this is the interior of the surface and the exterior. Exterior. And this is a surface, right? The surface I call it $S$. What is a surface? It could be a sphere, it could be a sort of — the skin of a \qt{potato}. It's a surface. Two-dimensional object. Right? It has and — and what I'm — this kind of surface will have a interior and a exterior. And the interior of the surface — so the — if the surface is the skin of the potato, the filling of the potato, the white part, right? Is $\Omega$. I call it $\Omega$.
\end{spoken-clean}

\begin{math-stroke}[Visualizing the Domain and Surface]
\begin{ai-tikz-planner-invisible-content}
% 1. Background: 3D Axes (x, y, z).
% 2. Midground: The 3D "potato" volume Omega, drawn with a semi-transparent fill.
% 3. Foreground: The surface S (boundary of Omega) and labels.
\end{ai-tikz-planner-invisible-content}
\begin{center}
\begin{tikzpicture}[scale=1.5, node distance=0.5cm]
    % 3D Axes
    \draw[->, thick, gray!50] (0,0,0) -- (2.5,0,0) node[right] {$x_1$};
    \draw[->, thick, gray!50] (0,0,0) -- (0,2.5,0) node[above] {$x_3$};
    \draw[->, thick, gray!50] (0,0,0) -- (0,0,2.5) node[below left] {$x_2$};

    % The Potato (Omega)
    \draw[thick, MidnightBlue, fill=MidnightBlue!10, opacity=0.6] 
        (1,1,0) to[out=90,in=180] (1.5,2,0) 
        to[out=0,in=90] (2.5,1,0) 
        to[out=-90,in=0] (1.5,0.2,0) 
        to[out=180,in=-90] cycle;
    
    % Labels
    \node[MidnightBlue] at (1.7, 1.1, 0) {\Large $\Omega$};
    \node[MidnightBlue, right] at (2.5, 1.5, 0) {$S$ (Surface)};
    \node[gray] at (3, 0.5, 0) {Exterior};
    \node[gray] at (1.7, 0.7, 0) {\footnotesize interior};
\end{tikzpicture}
\end{center}
\begin{explanation-of-steps}
The lecturer uses the \qt{potato} analogy to describe a bounded volume $\Omega \subset \mathbb{R}^3$ enclosed by a two-dimensional surface $S = \partial \Omega$. This geometric setup is foundational for the higher-dimensional versions of the Fundamental Theorem of Calculus.
\end{explanation-of-steps}
\end{math-stroke}

\begin{spoken-clean}[00:12:28 - 00:13:32]
And then the equivalent of this theorem (i.e., the FTC) in this setup is this. Given a vector field — now you don't know what it is a vector field, let me tell you. $f$ equal $f_1, f_2, f_3$. These are three — three functions of the space, right? At each point of space you have three numbers. Here you have three numbers, here you have three numbers. Yeah. This is a vector field and how we — do we draw vector fields? Um, we draw it like this, right? It's like a — a field of some cereal, right? But the — in every point you have an arrow instead of a — a plant.
\end{spoken-clean}

\begin{math-stroke}[Defining the Vector Field]
A vector field $f$ assigns a vector to every point in space:
\[ f = (f_1, f_2, f_3) \]
where each $f_i: \mathbb{R}^3 \to \mathbb{R}$ is a scalar function.

\begin{center}
\begin{tikzpicture}[scale=1.5]
% \begin{ai-tikz-planner-invisible-content}
% 1. Background: The potato volume Omega.
% 2. Foreground: A non-uniform vector field F drawn inside the volume.
% \end{ai-tikz-planner-invisible-content}
    % The Potato (Omega)
    \draw[thick, MidnightBlue, fill=MidnightBlue!5] 
        (1,1) to[out=90,in=180] (1.5,2) 
        to[out=0,in=90] (2.5,1) 
        to[out=-90,in=0] (1.5,0.2) 
        to[out=180,in=-90] cycle;
    
    % Vector Field F (Arrows inside)
    \foreach \x in {1.3, 1.6, 1.9, 2.2}
        \foreach \y in {0.6, 1.0, 1.4, 1.8}
            {
                \pgfmathsetmacro{\u}{0.2*cos(\x*50 + \y*30)}
                \pgfmathsetmacro{\v}{0.2*sin(\x*40 - \y*20)}
                \draw[->, BrickRed, thin] (\x,\y) -- (\x+\u, \y+\v);
            }
    
    \node[BrickRed] at (2.8, 0.8) {$f$ (Vector Field)};
\end{tikzpicture}
\end{center}
\end{math-stroke}

\begin{spoken-clean}[00:13:32 - 00:15:15]
Okay, so this is the vector field $f$. And the theorem says that the integral in $\Omega$ of the — the derivative of the field, the $i$ component with respect to the $x_i$ variable, and you sum over $i$. This is equal to the flux of the field across the boundary, across the $S$. And it's called like this. Is $f \cdot \nu$ differential of surface. This is the same theorem you have got here in one dimension (Recall: FTC) in several variables.
\end{spoken-clean}

\begin{math-stroke}[The Divergence Theorem (Gauss's Theorem)]
The \qt{final boss} of Analysis II relates the volume integral of the divergence to the surface integral of the flux:
\begin{color-box}{Black}[Divergence Theorem]
\[ \int_{\Omega} \left( \sum_{i} \frac{\partial f_i}{\partial x_i} \right) \, dx = \int_{S} f \cdot \nu \, dS \]
\end{color-box}

\begin{explanation-of-steps}
This equation is the higher-dimensional analogue of the Fundamental Theorem of Calculus. It states that the total \qt{source} or \qt{sink} of a vector field within a volume $\Omega$ (the left side) is exactly equal to the net flow or \qt{flux} of that field through the enclosing surface $S$ (the right side).
\end{explanation-of-steps}
\end{math-stroke}

\begin{spoken-clean}[00:15:15 - 00:17:38]
Now the small problem is that maybe some of you in physics or — or in previous courses have seen some of these notations, probably not. I don't expect you to — we — and we will need to introduce each of these notations. These partial derivatives, what is a domain, what is a surface? What is this surface mathematically? How do you model a surface? How do you compute a integral over a surface? You know how to compute a integral over the real line, right? Over a segment, but over a surface? This will need some definition and some work. And these objects that are the partial derivatives. So here in — in one equation you have a bit the goal and the syllabus of the course. How will — how will we proceed? So first, we will start looking at the space. In Analysis 1, you did everything in $\mathbb{R}$. Now we will need to introduce $\mathbb{R}^3$, but we are mathematicians so we will — well, some of us. Some of you are physicists or are proto-physicists, would like to become physicists. So we will use this because we don't want to spend time proving the same theorem in — in the plane $\mathbb{R}^2$, then in $\mathbb{R}^3$ in the space, and then when you go to do general relativity you need $\mathbb{R}^4$, right? We do everything in $\mathbb{R}^n$. $n$ is the dimension. So what is $\mathbb{R}^n$? Is the set of $n$ — $n$-tuples of real numbers. This is a definition. Okay? And what is — this is the same as the set of points — we will denote a point $x$, but now $x$ will carry $n$ coordinates. $x_1, x_2, \dots, x_n$. Now this is a point. And these $x_i$ are real numbers. This you know probably because you — you are doing or you have done linear algebra, right? So this should be quite familiar. But still is good that we recall.
\end{spoken-clean}

\begin{math-stroke}[Definition of \texorpdfstring{$\mathbb{R}^n$}{Rn}]
\begin{definition}[The Space $\mathbb{R}^n$]
The space $\mathbb{R}^n$ is defined as the set of all $n$-tuples of real numbers:
\[ \mathbb{R}^n = \{ n\text{-tuples of real numbers} \} \]
Formally, a point $x \in \mathbb{R}^n$ is represented by its coordinates:
\[ x = (x_1, x_2, \dots, x_n) \quad \text{where } x_i \in \mathbb{R} \text{ for all } i \in \{1, \dots, n\} \]
\end{definition}
\end{math-stroke}

\begin{spoken-clean}[00:17:38 - 00:21:13]
Okay. Now, the first part of the course we will be studying $\mathbb{R}^n$ and the properties of $\mathbb{R}^n$. The same a bit as in $\mathbb{R}$ you studied the — the axiom of the supremum, right? Or the — the fact that you can take, uh, limits of Cauchy sequences, these kind of facts, right? The same we will need to study in $\mathbb{R}^n$. Then there will be some new things or — or things that you could also study in $\mathbb{R}$ but become very trivial and now will be much more interesting is the topology of $\mathbb{R}^n$. So first, the $\mathbb{R}^n$ as a space, as Euclidean space. So this is a bit the syllabus. 2: the topology of $\mathbb{R}^n$. You will see what this means indeed, the topology. Then 3: this object over here \inlinemetanote{points at the partial derivatives in the Divergence Theorem} (i.e., the partial derivatives). Partial derivatives. So when you have — when you have several variables you can differentiate a quantity with respect to each of them, right? And differential and — so the differentiation of functions of several variables. Then 4: integration. Integration in $\mathbb{R}^n$, here in a domain of $\mathbb{R}^n$, and on surfaces. But because we want to learn also how to integrate on surfaces, we will do a bit of introduction to differential geometry. Introduction to, uh, surfaces and \qt{submanifolds}. Now I know it's a lot of words, but we will introduce one by one. So surface is a usual two-dimensional object, right? Typically smooth. When it is a three-dimensional object in $\mathbb{R}^4$, for instance, we call this a submanifold. You can also call it surface but — but it's more precise the word submanifold. Or — or — or some two-dimensional object in a five-dimensional space. The — any — any combination of dimensions. And then when you know this, when you know, uh, yes, you will need to know how to integrate. Integration over $\mathbb{R}^n$ and surfaces.
\end{spoken-clean}

\begin{math-stroke}[Course Syllabus]
The roadmap for Analysis II consists of five major pillars:
\begin{enumerate}
    \setcounter{enumi}{0} \item $\mathbb{R}^n$ as Euclidean space.
    \setcounter{enumi}{1} \item Topology of $\mathbb{R}^n$.
    \setcounter{enumi}{2} \item Partial derivatives and differentiation of multivariable functions.
    \setcounter{enumi}{3} \item Introduction to surfaces and \qt{submanifolds} (Differential Geometry).
    \setcounter{enumi}{4} \item Integration over $\mathbb{R}^n$ and surfaces.
\end{enumerate}
\end{math-stroke}

\begin{spoken-clean}[00:21:13 - 00:22:00]
And hopefully then you will know all of these elements. This element, uh, this is the normal vector. This I can tell you now what is this. This is just the vector, the unit vector, the vector of length one that is normal to the surface, right? If you have a surface, this — this $\nu$ is just the — the normal to the surface. So this is easy. Um, but still we will — we will define this in detail. And after we know all of these, we will be able to prove that equation over here (i.e., the Divergence Theorem) and work with it and see why it is useful. Right? For many things. Okay. Uh, so this is a rough plan of the course.
\end{spoken-clean}

\begin{math-stroke}[The Normal Vector]
\begin{center}
\begin{tikzpicture}[scale=1.5]
    % Surface S
    \draw[thick, MidnightBlue] (0,0) to[bend left=30] (2,0.5);
    \node[MidnightBlue, below] at (1,0.2) {$S$};
    
    % Normal Vector nu
    \draw[->, thick, ForestGreen] (1, 0.35) -- (1, 1.1) node[above] {$\nu$};
    \node[ForestGreen, right] at (1, 0.7) {\footnotesize Unit Normal Vector};
\end{tikzpicture}
\end{center}
\end{math-stroke}

% [SYSTEM] Segment complete. Please prompt "Continue" for the remainder of the segment.
% ==========================================
% SECTION 3: THE ROLE OF PROOFS IN ANALYSIS
% ==========================================
\section{The Role of Proofs in Analysis}

\begin{spoken-clean}[00:22:00 - 00:24:15]
Now I wanted to — to do a very small survey about your attitude, um, attitude towards proofs. Okay? This is important to mention that. \inlinemetanote{Moves to the right chalkboard} So I will put two columns. Column left is \qt{I find proofs} — so proofs I mean rigorous, long, very careful mathematical proofs with all details. Okay? This is mean proofs. So here I will put: I find proof, um, essential, illuminating, useful to clarify, the best part of any course. And on this column over here I will put: I find the proof innecessary, confusing, mysterious, I don't know when I have finished the proof or not, right? Because it looks like — it sounds convincing but — mysterious, no clear — no clear rules of the game.
\end{spoken-clean}

\begin{math-stroke}[Survey: Attitude Towards Proofs]
The lecturer creates a comparative table on the board to gauge student sentiment regarding formal mathematical proofs.
\begin{center}
\begin{tabular}{l | l}
\textbf{Positive / Essential} & \textbf{Negative / Challenging} \\ \hline
Essential & Innecessary \\
Illuminating & Confusing \\
Useful to clarify & Mysterious \\
The best part & No clear rules of the game
\end{tabular}
\end{center}
\end{math-stroke}

\begin{spoken-clean}[00:24:15 - 00:26:15]
And now, okay, this is extreme — I mean, I'm — I'm making all this on purpose. This is extreme, right? This is extreme. I know that any — most of you would be maybe in the middle. But I would like to know which of you lean more towards this opinion or more towards this opinion. Okay? So let's start with this one. So who is here? Who loves proofs and... okay. \inlinemetanote{Students raise hands} Okay. And here... \inlinemetanote{Students raise hands for the other column} Okay, it's interesting. So it's roughly — there are a bit more people that love proofs, but — but okay. So there is no correct answer, right? Uh, so I would be — well, depending on where — in between.
\end{spoken-clean}

\begin{spoken-clean}[00:26:15 - 00:28:30]
So why — why am I asking this? So does it coin — is it — does it have some correlation with physicists and mathematicians or not? Sometimes. What do you think? It — it could make sense, not always. There are physicists that... but — it could make sense. So I think that, um, about this discussion of proofs, I wanted to — to use ten minutes today to tell you two things. First, why I believe — many people can have many answers — but why I believe that the proofs are important and what will we test in this course, right? Because it's — it's a bit correlated of course.

So the first thing is that in — in algebra many times you have a theorem that is very beautiful, very clean, and you can apply in many situations without modifying it, right? In analysis, we will see that, sometimes the theorem that you learned doesn't exactly apply to the situation you — you need. Then you need to know how to modify a little bit the theorem. And the only way to do that is to know the proof and then modify a — some step here and some step there.
\end{spoken-clean}

\begin{didactic-insight}[Analysis as a Malleable Tool]
The lecturer highlights a fundamental difference between Algebra and Analysis: in Analysis, theorems are often not \qt{black boxes} but templates. Understanding the proof is what grants the mathematician the flexibility to adapt a result to non-standard domains or boundary conditions.
\end{didactic-insight}

\begin{spoken-clean}[00:28:30 - 00:31:17]
Also, many times the proof contains the ideas on why the theorem is true, right? The theorem is a bunch of symbols as we wrote there sometimes or statements, but it's not very illuminating. It's okay to know that a fact is true, but when you know a bit why, it's — it's — it's more useful many times, right? Being able to export and to use in situations that are new and also, um, understanding better, meaning being able also sometimes to remember the — the theorem. Because if you — if you get quick at remembering proofs and arguments, sometimes you don't remember all the assumptions of a theorem and you can figure — you can figure them out just by — by repeating the proof in — in your mind, right?

This doesn't mean that you need to learn the proof in all the details all the time. It's very important — so that the proof is useful and illuminating and so on — it's very important that when you have this super long proof, you extract the main ideas of the proof. A very short — so you need to have several levels in your head. So a proof that is that long, you can remember three sentences like hints that are the key ingredients of the proof for you. And it's very good exercise to do that. Then you expand. You learn how to expand these — these summaries. Is it clear that? And how you learn how to expand? You learn how to expand, um, hints by doing exercises. That's why the exercises are important. If you become very good in solving exercises, you'll be able to compress the information very well because then you don't need to remember all the details. You just remember the highlights of the proof and then you fill the details because you are expert in doing that, right? So becoming expert in exercises is — is good to find the proof also more interesting and more useful and not that, um, much of an effort.

Now, some people that are here that find the proof confusing, mysterious, innecessary, is because sometimes the — the rules of the game and why the game is interesting is not that clear. So why when you finish a proof, why can you say that the proof is finished, right? Because the professor says so that this the proof QED? Is common sense? Um, what is — most people see that, oh, this looks like a argument that is conclusive. Most people. But not every time. And there are very subtle arguments. So that's why we need to spend sometimes a lot of time with things that seem obvious because I tell you: a continuous function crosses the zero axis at some point (Recall: Intermediate Value Theorem). Okay, of course it crosses. But — this intermediate value theorem, right? Continuous function crosses the — of course, but you need to prove it, right? And the proof cannot be just the intuition of I make a drawing... \inlinemetanote{Draws a line crossing an axis} I have the zero line, I have two points, one here, one here, I have a continuous function, so if I cannot hold the — the pen from the paper, I need to cross the line at some point. This is not a mathematical proof. But it's a correct intuition.
\end{spoken-clean}

\begin{math-stroke}[Intuition vs. Rigor]
\begin{center}
\begin{tikzpicture}[scale=1.0]
    % Axes
    \draw[->] (-1,0) -- (3,0) node[right] {$x$};
    \draw[->] (0,-1) -- (0,2) node[above] {$y$};
    
    % Points
    \fill (0.5, -0.5) circle (2pt) node[below] {$f(a) < 0$};
    \fill (2.5, 1.5) circle (2pt) node[above] {$f(b) > 0$};
    
    % Continuous curve
    \draw[thick, MidnightBlue] (0.5, -0.5) .. controls (1.5, -0.2) and (1.0, 1.2) .. (2.5, 1.5);
    
    % Intersection
    \fill[BrickRed] (1.35, 0) circle (2pt);
    \node[BrickRed, below right] at (1.35, 0) {Root?};
\end{tikzpicture}
\end{center}
\begin{explanation-of-steps}
The lecturer uses the Intermediate Value Theorem to illustrate that while geometric intuition (the \qt{pen on paper} argument) is powerful, it does not constitute a formal proof. Rigorous analysis requires proving these \qt{obvious} facts from foundational axioms.
\end{explanation-of-steps}
\end{math-stroke}

\begin{spoken-clean}[00:31:17 - 00:31:17]
Now, so what can help with this? I think that something that can help with this, with the rules of what is a proof and what is not a proof... \inlinemetanote{Turns on the projector} Now we will need to turn on the projector again. Okay. Let me see. \inlinemetanote{Projector starts up} Hmm. Okay. So if someone is interested, you can go — you can Google, let's see, Google... you can Google, uh, Natural Number Game. \inlinemetanote{Lecturer searches for the game on Google} This is a serious stuff. I mean, I'm — I'm not — this is a very serious stuff. It's difficult. I mean, I know that there are some cartoons, but this is — is a serious difficult game. So — so it's — so this is a — so one — one argument and — and now seriously, in the light of artificial intelligence and — and all what is going on. So one argument that is very interesting about proofs is that — maybe some of you didn't realize — but a proof is computable. So if a proof is correct or not can be computed. That's why it's not an opinion, right? You can — there are several computation languages for... \textit{[Audio cuts off abruptly]}
\end{spoken-clean}

\begin{meta-note}[Projected Content: The Natural Number Game]
The lecturer shows the \qt{Natural Number Game} website, an interactive tool for learning formal proof verification using the Lean proof assistant. The interface shows a progression of \qt{worlds} representing different mathematical concepts.
\end{meta-note}

\begin{math-stroke}[The Structure of Formal Proofs (Lean)]
\begin{center}
\begin{tikzpicture}[node distance=1.5cm, every node/.style={draw, rounded corners, fill=gray!10, font=\footnotesize}]
    % Nodes
    \node (Tutorial) {Tutorial World};
    \node (Addition) [below=of Tutorial] {Addition World};
    \node (Multiplication) [below left=of Addition] {Multiplication World};
    \node (Power) [below right=of Addition] {Power World};
    \node (Function) [below=of Multiplication] {Function World};
    \node (Proposition) [below=of Power] {Proposition World};
    \node (AdvAdd) [below=of Function] {Advanced Addition World};
    \node (AdvMult) [below=of Proposition] {Advanced Multiplication World};
    \node (Inequality) [below=of AdvAdd] {Inequality World};

    % Connections
    \draw[->] (Tutorial) -- (Addition);
    \draw[->] (Addition) -- (Multiplication);
    \draw[->] (Addition) -- (Power);
    \draw[->] (Multiplication) -- (Function);
    \draw[->] (Power) -- (Proposition);
    \draw[->] (Function) -- (AdvAdd);
    \draw[->] (Proposition) -- (AdvMult);
    \draw[->] (AdvAdd) -- (Inequality);
\end{tikzpicture}
\end{center}
\end{math-stroke}

% [SYSTEM] Video complete.

% --- TEIL 2 ---
\begin{spoken-clean}[00:00:00 - 00:01:56]
...people that find the proof confusing, mysterious, innecessary, is because sometimes the --- the rules of the game and why the game is interesting is not that clear. So, why when you finish a proof, why can you say that the proof is finished, right? Because the professor says so that this is the proof Q.E.D.? Is common sense? Um, what is --- most people see that, oh, this looks like a argument that is conclusive. Most people. But not every time. And there are very subtle arguments. So, that's why we need to spend sometimes a lot of time with things that seem obvious because I tell you: a continuous function crosses the zero axis at some point (Recall: Intermediate Value Theorem). Okay, of course it crosses. But --- this intermediate value theorem, right? Continuous function crosses the --- of course, but you need to prove it, right? And the proof cannot be just the intuition of I make a drawing --- I have the zero line, I have two points, one here, one here, I have a continuous function, so if I cannot hold the --- the pen from the paper, I need to cross the line at some point. This is not a mathematical proof. But it's a correct intuition. Now, so what can help with this? I think that something that can help with this, with the rules of what is a proof and what is not a proof... \inlinemetanote{The lecturer turns on the projector and searches on Google} Now we will need to turn on the projector again. Okay. Let me see. Um. Okay. So, if someone is interested, you can go --- you can Google, let's see, Google... you can Google, uh, Natural Number Game. \inlinemetanote{The lecturer navigates to the Natural Number Game website}
\end{spoken-clean}

\begin{meta-note}[Projected Content: Lean Game Server]
The lecturer displays the \qt{Lean Game Server} website, specifically the \qt{Natural Number Game}. The site is described as a repository of learning games for the proof assistant Lean (Lean 4) and the mathematical library \qt{mathlib}. Other games shown include \qt{Relativity}, \qt{Field Theory}, \qt{Thermodynamics}, and \qt{Real Analysis, The Game}.
\end{meta-note}

\begin{spoken-clean}[00:01:56 - 00:04:08]
So, if someone is interested, you can go --- you can Google, let's see, Google... you can Google, uh, Natural Number Game. Is --- this is a serious stuff. I mean, I'm --- I'm not --- this is a very serious stuff. It's difficult. I mean, I know that there are some cartoons, but this is --- is a serious difficult game. So --- so it's --- so this is a --- so one --- one argument and --- and now seriously, in the light of artificial intelligence and --- and all what is going on. So, one argument that is very interesting about proofs is that --- maybe some of you didn't realize --- but a proof is computable. So, if a proof is correct or not can be computed. That's why it's not an opinion, right? You can --- there are several computation languages for proofs, one of the most famous is this Lean, where you can write the --- the proof in the language. It's very tedious to write the proof in the language, but now there are tools that help you. And now the --- it's like a computer program. If the program compiles, the proof is correct up to the axioms. You provide the axioms, you provide the reasoning, and if the program compiles, it's --- proof is correct. And if it doesn't compile, it doesn't mean that it's not correct. It means that you need to put more details. Maybe it's not correct, maybe yes, but it's not detailed enough. Right? So, first argument is that, um, so --- or some important argument is that proof are computable, this means that this ends all the discussion. If I wrote a proof, it's true, that's it. End of the discussion. Okay? Which is interesting. So, it --- it doesn't matter if Donald Trump says that the proof --- the fundamental theorem --- I mean, no, he can say whatever he wants, but the theorem will be true, okay? That nowadays is --- this is an interesting aspect.
\end{spoken-clean}

\begin{didactic-insight}[AI and the Verifiability of Proofs]
The lecturer uses the emergence of proof assistants like Lean to ground the concept of mathematical truth. By framing a proof as a \qt{computable} object rather than a matter of opinion, he highlights the absolute rigor required in Analysis. This serves as a meta-commentary on the reliability of AI models (like ChatGPT or Gemini), which may \qt{hallucinate} plausible-sounding but logically flawed arguments, whereas a compiled Lean proof is objectively verified.
\end{didactic-insight}

\begin{spoken-clean}[00:04:08 - 00:08:20]
This also an interesting aspect for these guys, right? \inlinemetanote{Points at the ChatGPT and Gemini tabs on his browser} Right? So, I --- I had on purpose the --- the ChatGPT or the Gemini or whatever. So, you are using those, right? You are using these --- I know that some people tell you \qt{don't use that} --- well, I know you are using those. Um, and I think it's --- I think you are --- you --- you know much better how to use than --- than people think. Um, so --- but you see, these guys still, especially with more advanced proofs, hallucinate. And say things sometimes that are wrong. And then knowing how to spot this is another motivation. Knowing how to spot wrong arguments and --- and being careful with correct and incorrect arguments --- um, is a good motivation to learn how to do correct proofs. Okay? And you can train by doing exercises, and if you are not sure, if you are --- now I --- I put it down. But if you are not sure what are the rules of the game, what is a proof, what is not a proof, try to --- I know it's difficult, try to go to the Natural Number Game and play a bit, and you will see what is a fully rigorous proof up to --- up to the end. And why we don't do this all the time. Because now --- the next thing I wanted to tell you is in Analysis 2, the pace will be quicker than in Analysis 1. We'll start roughly at the same pace, do --- doing many details, and then we will go a bit faster. What does it mean faster? That I will not be as detailed. You see? Some things is \qt{because this and that, then this follows}. And then you --- if you don't see that this follows, if many don't see, you ask question: \qt{Wow, uh, why this follows? I don't see --- we don't see}. Then we can provide details. But we will not be able to stop all the time to provide details because some of the arguments you already know how to do them. The Analysis 1 arguments, I will say \qt{okay, this is like in Analysis 1}, right? So, the pace is quicker and quicker. And the more you advance in your studies, it will be quicker. This doesn't mean that you can be okay with not --- uh, understanding, right? You need to --- if some statement you don't know how to prove, you can try by yourself, you can ask your peers, you can unfold everything. And this is a very good exercise. So, you need to do your exercises, you need to try to understand all the steps of the proof, meaning that ideally you should be able with time and --- and effort to --- to go back to the axioms, right? You don't need to do every time. You --- you just need to --- to be --- to know that you are able to. Okay? And then you are okay. So, if you can do this with --- with most --- proofs, um, it --- it means that you are following very well the course. Okay?
\end{spoken-clean}

\begin{spoken-clean}[00:08:20 - 00:10:30]
Is there something else I wanted to tell you before we start with the lecture? I don't think so. So, we have roughly ten minutes before the pause. I will start with the proper --- lecture --- the second half. And then --- uh, let me ask you, do you have some questions about what I said in this first part? A bit the plan of the course, how to be successful in the course, what do we expect from you? I --- now I have the impression that I said many things and maybe --- uh, I will need to insist --- uh, later on and come back to the objectives. Let me summarize if --- if there are no questions let me summarize in three points what are the --- a bit the main learning goals. Okay? So, maybe point number one that is very relevant to both physics and --- physicists and mathematicians is being able to do computations. Computations. Computations means --- um, advanced computations. Compute the --- the solution of some O.D.E., right? Several variable stuff. A computation can be compute some integral. Why the --- the area of the sphere is --- whatever, $4\pi$ or of the radius one, $4\pi r^2$? You need to be able to compute that integrals and --- with a numeric --- so with a --- um, --- with a concrete output like a numeric value. This --- this could be a concrete goal. Then you need to be able to follow proofs. Follow the --- the logic. And this you train with exercises, that some of them are computation --- also the computations, but some of them are more computation, some of them are more theoretical. And also you can train with the lecture notes. Whenever you don't understand a passage of the lecture notes, as I said before, you train --- train to fill out details. If you cannot, no problem, you ask, you ask your peers, then you ask me, and is a good exercise. Right? So, you have plenty of material to --- of material to train.
\end{spoken-clean}

% [SYSTEM] Segment complete. Please prompt "Continue" for the remainder of the segment.
\begin{spoken-clean}[00:10:30 - 00:11:36]
And essentially these are the two --- the two main legs: so computation and proofs. And then the third leg that I will introduce a bit is --- a bit assessing if a proof is correct. This is something we didn't test in the past, but now it's becoming more and more important. Okay? So, some type of exercise that I will try to introduce at some point is giving you some proof with a flaw argument. Sometimes I make mistakes, so then it's good because you can practice this and then you see... but I will try to make less mistakes. And then if I would then do some proofs with incorrect arguments on purpose and then you can spot \qt{ah, this subtlety you --- we didn't get it right} or here --- here the ChatGPT made a hallucination, okay? It's getting more and more harder and harder to get these models to hallucinate, but still in Analysis 2 it's --- is harder now. But --- but then we will put trickier questions. Um. Okay. So, it's five minutes. Do you have any question, someone? Ah, there? Is there a question? Yeah. Okay. Thank you. Sorry, I cannot see very well so...
\end{spoken-clean}

\begin{student-question}[Rigor in Tests]
I wanted to ask how rigorous do the proofs have to be in the test? Do we have to do them like the Analysis 1 test where we have --- we don't have to write down the axioms every time, or we just mention the most important theorem, or do we have to do very rigorous proofs?
\end{student-question}

\begin{spoken-clean}[continued]
Let me --- let me put the question about --- about --- let me put the question a bit different. The proof are always very rigorous. The only thing that they don't need to be as detailed as. You see? You can do some claims and you don't need to justify to the smallest detail every single claim. In Analysis 1, you see the difference of the theorem, it's clear.
\end{spoken-clean}

\begin{meta-note}[Blackboard Transition]
The lecturer erases the left and center chalkboards to summarize the transition from Analysis I to Analysis II.
\end{meta-note}

\begin{spoken-clean}[00:12:33 - 00:13:17]
So, what would be a --- a proof of this? So, how would you summarize in a sentence a proof of --- of this equation over here? Let's try this. \inlinemetanote{Points at the FTC on the board} Yeah. Fundamental Theorem of Calculus. \inlinemetanote{Starts writing on the board}
\end{spoken-clean}

\begin{math-stroke}[The Fundamental Theorem of Calculus vs. Divergence Theorem]
The lecturer restates the central comparison between the two courses:
\begin{itemize}
    \item \textbf{Analysis 1:} Calculus of one variable.
    \[ \text{FTC} \qquad \int_{a}^{b} f'(x) \, dx = f(b) - f(a) \]
    \item \textbf{Analysis 2:} Calculus in $\ge 2$ variables.
    \[ \text{Gauss/Divergence} \qquad \int_{\Omega} \sum_{i=1}^n \frac{\partial f_i}{\partial x_i} \, dx = \int_{S} f \cdot \nu \, dS \]
\end{itemize}
\end{math-stroke}

\begin{spoken-clean}[00:13:17 - 00:15:18]
So, anyone --- someone want to tell me one sentence that kind of explains the idea of the proof of that? \inlinemetanote{Points at the FTC}
\end{spoken-clean}

\begin{student-question}[Student Answer]
We can show that the derivative of the primitive converges to the actual value of the function at the --- I think it was $a$ or no, $b$.
\end{student-question}

\begin{spoken-clean}[continued]
The derivative of the --- yeah. So, the derivative with respect to this --- this value, right? Yeah, this is a good --- a good answer. So, if --- so everyone has --- this is what I was telling you about before. So, when you see a --- something like this, the first time you see, you see all the proof, and then there --- there needs to be something that you remember. That is some summary of the proof that for each person might be a bit different. There are also several ways to do the proof. And then with this, you are able to reconstruct the details. You want to know what --- what I remember when I see this? I see --- I always remember the \qt{telescopic sum}. Do you know what is a telescopic sum, right? So, when I see this, I see --- I see that. Right? I see... let me put it there.
\end{spoken-clean}

\begin{didactic-insight}[The Telescopic Sum Analogy]
The lecturer provides a powerful mental anchor for the Fundamental Theorem of Calculus: the telescopic sum. By viewing the integral as a sum of infinitesimal increments where internal terms cancel out, leaving only the boundary values, the student gains a structural intuition that carries over to the Divergence Theorem in higher dimensions.
\end{didactic-insight}

\begin{math-stroke}[Intuition: The Telescopic Sum]
The proof of the FTC can be visualized by partitioning the interval $[a, b]$ into $N$ small sub-intervals:
\begin{center}
\begin{tikzpicture}
    \draw[thick] (0,0) -- (5,0);
    \foreach \x in {0, 0.5, 1, 1.5, 2, 2.5, 3, 3.5, 4, 4.5, 5}
        \draw (\x, 0.1) -- (\x, -0.1);
    \node[below] at (0, -0.1) {$a = x_0$};
    \node[below] at (5, -0.1) {$b = x_N$};
    \node[below] at (2.5, -0.1) {$x_i$};
\end{tikzpicture}
\end{center}
The difference $f(b) - f(a)$ is expanded as a sum of increments:
\begin{align*}
f(b) - f(a) &= f(x_N) - f(x_0) \\
&= \underbrace{(f(x_N) - f(x_{N-1})) + (f(x_{N-1}) - f(x_{N-2})) + \dots + (f(x_1) - f(x_0))}_{\text{Telescopic Sum}} \\
&= \sum_{i=1}^N \big( f(x_i) - f(x_{i-1}) \big)
\end{align*}
\begin{explanation-of-steps}
By multiplying and dividing each term by the step size $h = x_i - x_{i-1}$, each increment $\frac{f(x_i) - f(x_{i-1})}{h}$ approximates the derivative $f'(x_i)$. In the limit as $h \to 0$, this sum becomes the Riemann integral $\int_a^b f'(x) \, dx$.
\end{explanation-of-steps}
\end{math-stroke}

\begin{spoken-clean}[00:15:18 - 00:17:15]
So, you are unfolding this sum. Now, if you call the distance between these two consecutive points $h$, you multiply and divide by $h$, right? And each of these pieces is approximately the derivative. End of the proof. Okay? I will not go at this speed. But this is a way to remember the proof. And you need to find your own way in each proof. And then it's not that hard to remember them. Okay, so it's nine. We do the 15 minutes pose and then we start with the proper lecture. \inlinemetanote{The lecturer takes a break. The video resumes with the lecturer at the board} So, I will --- I will do my best, I will do my best to respect the schedule and the times, but please --- uh, I would ask you to do also your best to respect the 15 minutes break and be here 30 seconds before the bell --- rings and so on. Because otherwise we lost two minutes that of course are not subtracted from the exam, it just that I rush some proof, right? So, it's --- is in your interest to be on time.
\end{spoken-clean}

% ==========================================
% SECTION 4: THE SPACE R^n
% ==========================================
\section{The Space \texorpdfstring{$\mathbb{R}^n$}{Rn}}

\begin{spoken-clean}[00:17:15 - 00:19:02]
So, $\mathbb{R}^n$ we said before is the set of $n$-tuples of real numbers, right? \inlinemetanote{Writes on the board} And we will be concerned with $\mathbb{R}^n$, $n$ in $\mathbb{N}$ is the dimension of the space. Can you see from up there this? Yeah? You have very good eyes because I could not see anything. Okay. Now, $\mathbb{R}^n$ by itself --- well, --- just vectors of --- of numbers of real numbers is not --- that interesting, we will need to add some extra structure. So, what is the Euclidean --- Euclidean structure? You already know, I think, but we will say again.
\end{spoken-clean}

\begin{math-stroke}[Definition of \texorpdfstring{$\mathbb{R}^n$}{Rn}]
\begin{definition}[The Space $\mathbb{R}^n$]
The space $\mathbb{R}^n$ is the set of all $n$-tuples of real numbers:
\[ \mathbb{R}^n = \{ x = (x_1, \dots, x_n) \mid x_i \in \mathbb{R} \} \]
The natural number $n \in \mathbb{N}$ denotes the \emph{dimension} of the space.
\end{definition}
\end{math-stroke}

\begin{spoken-clean}[00:19:02 - 00:21:18]
So, in $\mathbb{R}^n$ what can we do? We can measure --- we can, well, sum vectors, right? \inlinemetanote{Writes the sum formula} We can multiply --- this is like linear algebra. If $\lambda$ is a real, we can multiply by a --- a scalar. This is a scalar, a real number. So, $\lambda x_1, \dots, \lambda x_n$. You know that from linear algebra very well. And then we can also measure the Euclidean structure means that we can also measure distances between points, right? So, we have the norm, the Euclidean norm. \inlinemetanote{Writes the norm formula}
\end{spoken-clean}

\begin{math-stroke}[Vector Space Structure of \texorpdfstring{$\mathbb{R}^n$}{Rn}]
The space $\mathbb{R}^n$ is equipped with the standard operations of a vector space:
\begin{description}
    \item[Vector Addition:] For $x, y \in \mathbb{R}^n$:
    \[ x + y = (x_1 + y_1, \dots, x_n + y_n) \]
    \item[Scalar Multiplication:] For $x \in \mathbb{R}^n$ and $\lambda \in \mathbb{R}$:
    \[ \lambda x = (\lambda x_1, \dots, \lambda x_n) \]
\end{description}
\end{math-stroke}

% [SYSTEM] Segment complete. Please prompt "Continue" for the remainder of the segment.
\begin{spoken-clean}[00:21:18 - 00:22:37]
So, first the norm, the Euclidean norm of a vector $x$ is the root, the square root of the square --- the sum of the squares of the components. $x_1^2$ to --- $x_n^2$. Right? \inlinemetanote{Writes the formula on the board} This is the length of the vector. So, in the plane, let's say... \inlinemetanote{Draws a 2D coordinate system} Let's say this is $x_1$ and this is $x_2$, this is the point $x$. The length of --- of this vector $x$, the length, is $x_1^2 + x_2^2$ square root. Pythagoras. Right? Pythagoras theorem. And the same applies to every dimension.
\end{spoken-clean}

\begin{math-stroke}[The Euclidean Norm]
The Euclidean norm (or length) of a vector $x \in \mathbb{R}^n$ is defined as:
\[ \|x\| = \sqrt{\sum_{i=1}^n x_i^2} = \sqrt{x_1^2 + \dots + x_n^2} \]

\begin{center}
\begin{tikzpicture}[scale=1.5]
% \begin{ai-tikz-planner-invisible-content}
% 1. Background: 2D Axes.
% 2. Midground: Dashed lines to components x1, x2.
% 3. Foreground: The vector x and its label.
% \end{ai-tikz-planner-invisible-content}
    % Axes
    \draw[->, thick] (-0.5,0) -- (3,0) node[right] {$x_1$};
    \draw[->, thick] (0,-0.5) -- (0,2.5) node[above] {$x_2$};

    % Components
    \draw[dashed] (2,0) -- (2,1.5) -- (0,1.5);
    \node[below] at (2,0) {$x_1$};
    \node[left] at (0,1.5) {$x_2$};

    % Vector
    \draw[->, very thick, MidnightBlue] (0,0) -- (2,1.5) node[midway, above left] {$\|x\|$};
    \fill (2,1.5) circle (2pt) node[above right] {$x = (x_1, x_2)$};
\end{tikzpicture}
\end{center}

\begin{explanation-of-steps}
The Euclidean norm is a direct generalization of the Pythagorean theorem to $n$ dimensions, representing the straight-line distance from the origin to the point $x$.
\end{explanation-of-steps}
\end{math-stroke}

\begin{spoken-clean}[00:22:37 - 00:24:42]
Okay. And we can also measure scalar products or --- inner products that have, you know, maybe that they have to do with angles, right? The scalar products. So, we will use these notations: $x \cdot y$ or sometimes this also, $\langle x, y \rangle$. Both are fine. \inlinemetanote{Writes the summation formula} This should not be too surprising. Scalar product. And then I said the norm of a vector, that is the length of the vector, and then the distance between points. What is the distance between points? If I take two points here... \inlinemetanote{Draws two points $x$ and $y$ in space} One point $x$ and one point $y$ over here, the distance between the two points is the length of the vector that joins them, no? The distance between $x$ and $y$ is the length of the vector that joins them. This a definition. This is the Euclidean distance.
\end{spoken-clean}

\begin{math-stroke}[Scalar Product and Distance]
The space is further structured by the \emph{scalar product} (or inner product):
\begin{color-box}{BrickRed}[Scalar Product]
\[ x \cdot y = \langle x, y \rangle = \sum_{i=1}^n x_i y_i \]
\end{color-box}

The \emph{Euclidean distance} between two points $x, y \in \mathbb{R}^n$ is defined as the norm of their difference:
\begin{equation} \label{eq:distance_def}
d(x, y) = \|y - x\| = \sqrt{\sum_{i=1}^n (y_i - x_i)^2}
\end{equation}

\begin{center}
\begin{tikzpicture}[scale=1.2]
    \coordinate (X) at (0.5, 0.5);
    \coordinate (Y) at (3, 2);
    
    \draw[->, thick, gray!50] (0,0) -- (4,0);
    \draw[->, thick, gray!50] (0,0) -- (0,3);
    
    \fill (X) circle (2pt) node[below left] {$x$};
    \fill (Y) circle (2pt) node[above right] {$y$};
    
    \draw[thick, BurntOrange, <->] (X) -- (Y) node[midway, above left] {$d(x, y)$};
\end{tikzpicture}
\end{center}
\end{math-stroke}

\begin{spoken-clean}[00:24:42 - 00:26:37]
Now, the first lemma I will prove about Euclidean distance is the triangle inequality. Triangle inequality. Do you know the --- some of you know the triangle inequality? Have you seen that? Who --- how many people have not seen the triangle inequality? Ah, you all know the triangle inequality. Okay, then I will be quick on that. So, the triangle inequality says this: for all $x, y, z$ in $\mathbb{R}^n$, for all --- triples of vectors, if I do $z - x$, the length of this is less than the length of going from $x$ to $y$ plus the length of going from $y$ to $z$. Right? \inlinemetanote{Writes the inequality on the board} Maybe you have seen the triangle inequality with the absolute value? Have --- but you have also seen that one, with the modulus, with the length. Right? So, what this says is that whenever you have three points $x, y, z$ in the space in any position... \inlinemetanote{Draws a triangle with vertices $x, y, z$} The length --- so the --- the distance between $x$ and $z$ is always longer than this --- the sum of these two distances. Right? Or equal. And actually it is equal when the point is aligned with the two, right? But I'm not saying this yet.
\end{spoken-clean}

\begin{math-stroke}[Lemma: The Triangle Inequality]
\begin{color-box}{black}[Triangle Inequality]
\begin{lemma}[Triangle Inequality] \label{lem:triangle_inequality}
For all $x, y, z \in \mathbb{R}^n$, the following inequality holds:
\begin{equation} \label{eq:triangle_xyz}
\|z - x\| \le \|y - x\| + \|z - y\|
\end{equation}
\end{lemma}
\end{color-box}

\begin{center}
\begin{tikzpicture}[scale=1.5]
    \coordinate (X) at (0,0);
    \coordinate (Y) at (1, 1.5);
    \coordinate (Z) at (3, 0.5);
    
    \draw[thick, MidnightBlue] (X) -- (Z) node[midway, below] {$\|z - x\|$};
    \draw[thick, BrickRed] (X) -- (Y) node[midway, above left] {$\|y - x\|$};
    \draw[thick, BrickRed] (Y) -- (Z) node[midway, above right] {$\|z - y\|$};
    
    \fill (X) circle (1.5pt) node[below left] {$x$};
    \fill (Y) circle (1.5pt) node[above] {$y$};
    \fill (Z) circle (1.5pt) node[below right] {$z$};
\end{tikzpicture}
\end{center}

\begin{explanation-of-steps}
Geometrically, the triangle inequality states that the shortest path between two points is a straight line. Any detour through a third point $y$ must result in a total distance that is greater than or equal to the direct distance.
\end{explanation-of-steps}
\end{math-stroke}

\begin{spoken-clean}[00:26:37 - 00:29:18]
So, let's prove that. So, to prove that equivalently I will prove --- equivalently, I'm saying that $\|x + y\|$ is less than or equal to $\|x\| + \|y\|$. \inlinemetanote{Writes the equivalent form on the board and labels it with a red star} So, why is this equivalent to this one? Well, equivalent is means that if you know how to prove this, easily you prove that, right? Why is this equivalent to this one? Well, you see the difference of the theorem, it's clear. So, what would be a proof of this? So, you need to do this substitution. You can call $a$ the vector $y - x$, say, then $b$ is equal to $z - y$, and then automatically if you sum $a + b$ is --- is $z - x$. Right? \inlinemetanote{Writes the substitution on the board} This is like putting --- so one --- so if you want to go from here to here, you just put as he said $z$ equal to --- so you need to put here is $x + y$. Well, you need to do this substitution, this kind of substitution, and it works both ways, right? Because if you do this, you will get $a + b$ less than or equal to to that and --- and the opposite. If you --- use this inequality with this substitution, you will get that. Now, this is as detailed as I will be, again. A bit the difference between Analysis 1 and 2. If this equivalence poses some trouble, it shouldn't pose too much, but you can work a bit by yourself, is a nice exercise. Okay? And if you cannot, um, ask later.
\end{spoken-clean}

\begin{math-stroke}[Equivalence of Triangle Inequality Forms]
The statement in Equation \ref{eq:triangle_xyz} is equivalent to the following simpler form:
\begin{color-box}{black}[Equivalent Form]
\begin{equation} \label{eq:triangle_star}
\|x + y\| \le \|x\| + \|y\| \tag{\textcolor{BrickRed}{$\ast$}}
\end{equation}
\end{color-box}

\begin{proof}[Proof of Equivalence]
To see the equivalence, we perform a change of variables. Let:
\begin{align*}
a &= y - x \\
b &= z - y
\end{align*}
Summing these vectors yields:
\[ a + b = (y - x) + (z - y) = z - x \]
Substituting these into Equation \ref{eq:triangle_xyz}, we obtain:
\[ \|a + b\| \le \|a\| + \|b\| \]
which is exactly the form of Equation \ref{eq:triangle_star}. Since the mapping $(x, y, z) \mapsto (a, b)$ is surjective, the two statements are logically equivalent.
\end{proof}
\end{math-stroke}

\begin{spoken-clean}[00:29:18 - 00:31:17]
Okay, so this is the --- since we convinced ourselves that the two statements are equivalent and this is shorter, I will prove this one. So, now the goal is to prove that inequality. Okay? And now I will do it maybe in the next blackboard. \inlinemetanote{Moves to the next chalkboard} So, I start writing... let me call this inequality star. \inlinemetanote{Writes \qt{*} next to the formula} So, this what I want to show is equivalent after squaring, right? I square and I said $\|x + y\|^2$ is equal --- is less than or equal to $\|x\|^2 + \|y\|^2 + 2\|x\|\|y\|$. Right? \inlinemetanote{Writes the squared inequality} That inequality is equivalent to this one because is an inequality between positive numbers and the square --- squaring is a monotone function. This is the part that I said in words but I will not write because we are in Analysis 2. Okay? This you learned in Analysis 1. Now, once I wrote it in this form, I can develop this. What is the modulus of squared? This is equal by definition to the sum of $(x_i + y_i)^2$ between $i$ equal 1 and $n$. Right? \inlinemetanote{Writes the summation} And this is --- and this I can go --- this is equal, develop the --- now these are real numbers, so I can develop the squares. This is the sum... \textit{[Audio cuts off abruptly]}
\end{spoken-clean}

\begin{math-stroke}[Proof Strategy for the Triangle Inequality]
\begin{ai-proof-skeleton-invisible-content}
% 1. Square both sides of the inequality (*).
% 2. Expand the left side using the definition of the norm and summation.
% 3. Use the Cauchy-Schwarz inequality to bound the cross terms.
\end{ai-proof-skeleton-invisible-content}
We aim to prove Equation \ref{eq:triangle_star}. Since both sides are non-negative, it is equivalent to prove the squared version:
\[ \|x + y\|^2 \le (\|x\| + \|y\|)^2 = \|x\|^2 + \|y\|^2 + 2\|x\|\|y\| \]
Expanding the left-hand side using the definition of the Euclidean norm:
\begin{align*}
\|x + y\|^2 &= \sum_{i=1}^n (x_i + y_i)^2 \\
&= \sum_{i=1}^n (x_i^2 + 2x_i y_i + y_i^2) \\
&= \sum_{i=1}^n x_i^2 + 2\sum_{i=1}^n x_i y_i + \sum_{i=1}^n y_i^2 \\
&= \|x\|^2 + 2(x \cdot y) + \|y\|^2
\end{align*}
\textit{[Audio cuts off abruptly]}
\end{math-stroke}

% [SYSTEM] Video complete.

% --- TEIL 3 ---
% ==========================================
% SECTION 4: THE SPACE R^n (CONTINUED)
% ==========================================

\begin{spoken-clean}[00:00:00 - 00:01:24]
Right? Because if you do this (i.e., the substitution $a = y - x$ and $b = z - y$), you will get $a + b$ less than or equal to to that and --- and the opposite. If you, uh, use this inequality with this substitution, you will get that. Now, this is as detailed as I will be, again. A bit the difference between Analysis 1 and 2. If this equivalence poses some trouble, it shouldn't pose too much, but you can work a bit by yourself, is a nice exercise. Okay? And if you cannot, um, ask later.

Okay, so this is the --- since we convinced ourselves that the two statements are equivalent and this is shorter, I will prove this one (i.e., Equation \ref{eq:triangle_star}). So, now the goal is to prove that inequality. Okay? And now I will do it maybe in the next blackboard. \inlinemetanote{Moves to the center-left chalkboard} So, I start writing... let me call this inequality star. \inlinemetanote{Writes \qt{*} next to the formula} So, this what I want to show is equivalent after squaring, right? I square and I said...
\end{spoken-clean}

\begin{math-stroke}[Proof of the Triangle Inequality (Initial Attempt)]
\begin{ai-proof-skeleton-invisible-content}
% 1. Goal: Prove ||x + y|| <= ||x|| + ||y||.
% 2. Method: Square both sides and expand.
% 3. Note: The lecturer initially writes x - y on the board.
\end{ai-proof-skeleton-invisible-content}
To prove the triangle inequality, we consider the squared norm of the sum of two vectors. \inlinemetanote{The lecturer initially writes the formula for the difference $x - y$ on the board}
\begin{equation} \label{eq:triangle_sq_wrong}
\|x - y\|^2 \le \|x\|^2 + \|y\|^2 + 2\|x\|\|y\| \quad \textcolor{red}{\text{(Incorrect notation)}}
\end{equation}
Expanding the left-hand side by the definition of the Euclidean norm:
\[
\sum_{i=1}^n (x_i - y_i)^2 = \dots
\]
\end{math-stroke}

\begin{spoken-clean}[00:01:24 - 00:02:15]
...I square and I said $x$ minus $y$ squared (i.e., actually $x + y$ as intended for the proof) is equal --- is less than or equal to $x$ squared plus $y$ squared plus two times the product of the modulus. Right? That inequality is equivalent to this one because is an inequality between positive numbers and the square --- squaring is a monotone function. This is the part that I said in words but I will not write because we are in Analysis 2. Okay? This you learned in Analysis 1. Now, once I wrote it in this form, I can develop this. What is the modulus of squared? This is equal by definition to the sum of $x_i$ minus $y_i$ squared between $i$ equal 1 and $n$. Right?
\end{spoken-clean}

\begin{student-question}[Student Spotting the Sign Error]
Um, shouldn't it be $x + y$ at the beginning of the proof? There \inlinemetanote{points at the board}.
\end{student-question}

\begin{spoken-clean}[continued]
Ah, yes. Actually is the same but, yeah, now that you say it, is better. So, you can put $x + y$ is --- is equivalent to put $x + y$ and $x - y$ but --- but to match this inequality (i.e., Equation \ref{eq:triangle_star}) otherwise it would be with a minus. But all the proof you can do with the plus and with the minus but --- okay, then here will be a plus. Thank you. Thank you for the --- nice. \inlinemetanote{Corrects the board}
\end{spoken-clean}

\begin{math-stroke}[Proof of the Triangle Inequality (Corrected)]
\begin{proof}
We aim to prove Equation \ref{eq:triangle_star}: $\|x + y\| \le \|x\| + \|y\|$. Since both sides are non-negative, this is equivalent to proving:
\begin{equation} \label{eq:triangle_sq_correct}
\|x + y\|^2 \le (\|x\| + \|y\|)^2 = \|x\|^2 + \|y\|^2 + 2\|x\|\|y\|
\end{equation}
Expanding the left-hand side using the definition of the Euclidean norm and the properties of real numbers:
\begin{align*}
\|x + y\|^2 &= \sum_{i=1}^n (x_i + y_i)^2 \\
&= \sum_{i=1}^n (x_i^2 + y_i^2 + 2x_i y_i) \\
&= \underbrace{\sum_{i=1}^n x_i^2}_{\|x\|^2} + \underbrace{\sum_{i=1}^n y_i^2}_{\|y\|^2} + 2\underbrace{\sum_{i=1}^n x_i y_i}_{x \cdot y} \\
&= \|x\|^2 + \|y\|^2 + 2(x \cdot y)
\end{align*}
Substituting this back into Equation \ref{eq:triangle_sq_correct}, we see that the inequality holds if and only if:
\[
\|x\|^2 + \|y\|^2 + 2(x \cdot y) \le \|x\|^2 + \|y\|^2 + 2\|x\|\|y\|
\]
Canceling the common terms $\|x\|^2$ and $\|y\|^2$ and dividing by 2, we find that the triangle inequality is equivalent to the following condition:
\begin{equation} \label{eq:cs_goal}
x \cdot y \le \|x\|\|y\|
\end{equation}
\end{proof}
\end{math-stroke}

\begin{spoken-clean}[00:03:28 - 00:05:50]
Right? Which is equal to the modulus of $x$ squared plus the modulus of $y$ squared. This is the modulus of $x$ squared when you sum over $n$, over $i$. This is the modulus of $y$ squared and this is two times the scalar product. Right? So, I want to prove --- this is a bit like a snake, but these are equalities. This is equal to this, is equal to this. So, what I'm left with --- I'm left with proving that this thing is less than or equal to this thing. And you see here, this --- this I can simplify, right? I can simplify because it's both sides. So, what I want to show is that two $x$ dot $y$ is less than or equal than the product of the modulus. Right? So, after this calculation, what I want to prove is equivalent to this. Let me simplify also the two. Is equivalent to this (i.e., Equation \ref{eq:cs_goal}).

And this is an inequality that is famous. Now I will prove it in case you don't know. But this is the Cauchy-Schwarz inequality. Okay? And this is true, this inequality is true because is the Cauchy-Schwarz. \inlinemetanote{Writes \qt{true Cauchy-Schwarz inequality} on the board} And now, in case just to warm up and in case you didn't remember or you never saw the proof of the Cauchy-Schwarz inequality, I will prove the Cauchy-Schwarz inequality, right? So, if I can prove this inequality, you see that the proof of the triangle inequality is done, is concluded. I'm just left with proving this inequality that I told you is called the Cauchy-Schwarz inequality. So, let me do this in a --- so I will put a square that means end of the proof but up to this ingredient that I will prove next.
\end{spoken-clean}

\begin{nice-box}[Lemma: Cauchy-Schwarz Inequality]
\begin{lemma}[Cauchy-Schwarz Inequality] \label{lem:cauchy_schwarz}
For all $x, y \in \mathbb{R}^n$, the following inequality holds:
\begin{equation} \label{eq:cauchy_schwarz}
x \cdot y \le \|x\|\|y\|
\end{equation}
\end{lemma}
\end{nice-box}

\begin{math-stroke}[Proof of the Cauchy-Schwarz Inequality]
\begin{proof}
The first observation is that if either $x = 0$ or $y = 0$, the inequality is trivially satisfied as $0 \le 0$. Therefore, we assume without loss of generality (WLOG) that both $x \neq 0$ and $y \neq 0$.

Recall the definition of the scalar product:
\[ x \cdot y = \sum_{i=1}^n x_i y_i \]
We introduce a parameter $\lambda > 0$. For any such $\lambda$, we can multiply and divide each term by $\lambda$:
\[ 2(x \cdot y) = \sum_{i=1}^n 2(x_i \lambda) \left( \frac{y_i}{\lambda} \right) \]
Using the elementary inequality for real numbers $2ab \le a^2 + b^2$ (which follows from $(a-b)^2 \ge 0$), we apply this to each term in the sum:
\begin{align*}
2(x \cdot y) &\le \sum_{i=1}^n \left( (x_i \lambda)^2 + \left( \frac{y_i}{\lambda} \right)^2 \right) \\
&= \lambda^2 \sum_{i=1}^n x_i^2 + \frac{1}{\lambda^2} \sum_{i=1}^n y_i^2 \\
&= \lambda^2 \|x\|^2 + \frac{1}{\lambda^2} \|y\|^2
\end{align*}
This inequality holds for any $\lambda > 0$. To obtain the tightest bound, we choose $\lambda$ such that the two terms on the right are equal. Setting $\lambda^2 \|x\|^2 = \frac{1}{\lambda^2} \|y\|^2$ yields:
\[ \lambda^2 = \frac{\|y\|}{\|x\|} \]
Substituting this optimal value of $\lambda^2$ back into the inequality:
\begin{align*}
2(x \cdot y) &\le \left( \frac{\|y\|}{\|x\|} \right) \|x\|^2 + \left( \frac{\|x\|}{\|y\|} \right) \|y\|^2 \\
&= \|y\|\|x\| + \|x\|\|y\| \\
&= 2\|x\|\|y\|
\end{align*}
Dividing by 2, we arrive at the Cauchy-Schwarz inequality:
\[ x \cdot y \le \|x\|\|y\| \]
\end{proof}
\end{math-stroke}

\begin{spoken-clean}[00:05:50 - 00:08:20]
So, now I say Lemma, Cauchy-Schwarz inequality. Well, I don't copy. Lemma, this holds. Right? Cauchy-Schwarz inequality. In the notes will be as a separate lemma. But now I don't make me copy again because I'm going to write again exactly the same. So, this inequality holds for all $x$ and $y$ in $\mathbb{R}^n$. Okay? Now, proof. In the notes I have put two proofs. There are many ways to prove this. I have put two proofs of this inequality. I will with give one, okay?

The first observation is that if either $x$ equal zero or $y$ equal zero, the inequality is --- zero less than or equal to zero, right? Because if --- if the vector is --- if the vector one of the two vectors is zero, this product is zero and the modulus of one of the two will be zero so is zero less than or equal to zero. Which is trivially true. By this observation, I can assume without loss of generality that both vectors are not zero. So, assume without loss of generality (WLOG) both $x$ and $y$ are different from zero. This we have seen many times now you are in the second semester, this \qt{assume without loss of generality} phrasing.
\end{spoken-clean}

\begin{spoken-clean}[00:08:20 - 00:10:30]
Okay, so both vectors are non-zero. Now what I will do, um, I will take --- want to prove --- I will write what is the scalar product $x$ dot $y$. This is the sum between $i$ equal 1 and $n$, I copy again, $x_i y_i$. Right? And now I will do a trick. Um, I will take for every --- for every number $\lambda$ positive, I have the right to divide and multiply by $\lambda$. I will put $\lambda$ here and $\lambda$ over here. Right? I did nothing because I multiplied and divided by $\lambda$. And now I --- I will use an $n$ times an inequality of real numbers. You know that --- or let me --- let me multiply by two here and here. And now I will use an inequality of real numbers: $2ab$ is less than or equal to $a^2 + b^2$. Which is the same as --- essentially that $(a-b)^2$ is non-negative. So, this inequality of real numbers, I will apply $n$ times to each of these two numbers. This will be the $a$ and this will be the $b$ in each case. And what do I get? And with the two, right? I get this is less than or equal than the sum $i$ equal 1 to $n$. Now, the $\lambda$, I get $\lambda^2 x_i^2$ plus $y_i^2$ over $\lambda^2$. Okay? Which is the same as $\lambda^2$ mod $x$ squared plus 1 over $\lambda^2$ mod $y$ squared. And this holds for every $\lambda$. Right? This inequality holds for every $\lambda$ positive. I didn't use anything. Now is where I use that both are --- are not zero, so I can take $\lambda$ equal --- or $\lambda^2$ equal what? $y$ divided by $x$. We take that $\lambda$ and you get $\lambda^2$ two times, you see this --- when you get the 1 over $\lambda$ exactly simplifies what you want and you get two times the product. Okay? This is the proof of the Cauchy-Schwarz and this ends the proof of the triangle inequality. Is there any question about this proof?
\end{spoken-clean}

\begin{meta-note}[Segment Transition]
The lecturer has finished the proofs of the triangle inequality and the Cauchy-Schwarz inequality. He erases the center and right chalkboards to transition to the next major topic: Metric Spaces.
\end{meta-note}

% ==========================================
% SECTION 5: METRIC SPACES
% ==========================================
\section{Metric Spaces}

\begin{spoken-clean}[00:10:30 - 00:12:50]
This is the structure $\mathbb{R}^n$ with this structure of the scalar product, the length, is the main structure we will use in the proof to do analysis. But there are --- now we need to understand convergence sequences in $\mathbb{R}^n$, we need to understand Cauchy sequences and all these things that you understood in $\mathbb{R}$. Everything --- a bit the summary --- everything that you learned how to do in $\mathbb{R}$, now we need to learn how to do in $\mathbb{R}^n$. But --- to make it more efficient, we will learn all these concepts in a more general, um, --- from a more general viewpoint, that is from the viewpoint of metric spaces. Right? So that later when you modify --- so when you go to other --- courses or in this course later on we will have different metric spaces, I don't need to redefine all the time what convergence means and so on.

Now, I'm --- I need to warn you, especially for the --- well, I don't know if especially for the physicists, but this --- this --- now will get a bit more abstract. But bear with --- then we will come --- this will be for a --- for a couple of weeks or so. Later we will come back to $\mathbb{R}^n$, everything $\mathbb{R}^n$. But this is metric spaces topology is more axiomatic, more abstract. It's good to practice proofs because here you need to use axioms. You cannot use your intuitions all the time. You need to go back to the axioms of what is a metric space that I will tell you now and do the proofs. Sometimes working like mathematicians in more abstract viewpoint with few --- fewer axioms, um, is also easier. Well, depends, right? But it's easier because if you have less axioms --- now in $\mathbb{R}^n$ we have many axioms, we have the axiom of the supremum for each coordinate, you have many things. Now in the metric space, I will introduce three axioms. Just three. And then all the proofs need to use these three axioms. So, you don't need to think too much what are --- it's like --- like playing chess, I say sometimes, like playing chess only with the rooks. So, if you play chess only and there are only the rooks, it's not that difficult than playing with many pieces, right? So, but then you need a bit more of abstraction because it's not what you are used to. So, we are going to play chess with the rooks and we will define a metric space.
\end{spoken-clean}

\begin{didactic-insight}[The \qt{Chess with Rooks} Analogy]
The lecturer uses a playful analogy to explain the pedagogical value of abstraction. By stripping away the specific properties of $\mathbb{R}^n$ and focusing on the minimal axioms of a metric space, the \qt{rules of the game} become simpler but require more rigorous adherence. This prevents students from relying on geometric intuition that might fail in more exotic spaces (like function spaces).
\end{didactic-insight}

\begin{nice-box}[Definition: Metric Space]
\begin{definition}[Metric Space] \label{def:metric_space}
A \emph{metric space} is a pair $(X, d)$, where $X$ is a non-empty set and $d$ is a mapping (called the \emph{distance} or \emph{metric}):
\[ d: X \times X \to [0, \infty) \]
that satisfies the following three axioms for all $x, y, z \in X$:
\begin{enumerate}
    \setcounter{enumi}{0} \item \emph{Definiteness:} $d(x, y) = 0 \iff x = y$.
    \setcounter{enumi}{1} \item \emph{Symmetry:} $d(x, y) = d(y, x)$.
    \setcounter{enumi}{2} \item \emph{Triangle Inequality:} $d(x, z) \le d(x, y) + d(y, z)$.
\end{enumerate}
\end{definition}

\begin{explanation-of-steps}
The lecturer notes that while he formally defines $X$ as a set, he will not be overly pedantic about the case where $X$ might be empty, as it is trivial and lacks mathematical interest for the course.
\end{explanation-of-steps}
\end{nice-box}

\begin{spoken-clean}[00:12:50 - 00:15:15]
So, we will define a metric space is a pair $(X, d)$ where $X$ is a non-empty set. Well, let's say --- this also the kind of things that I will not be that careful about. So, if the set $X$ is empty or non-empty --- I mean, this is the part that I'm sorry to say that, but who cares? Uh, if the set is --- so if someone wants to say \qt{no, because with the empty set}, okay, in the notes probably everything is okay and I say that the set of is non-empty. In the lecture I will not focus on the case that the set can be empty or not, unless it is important and --- and hides some important subtlety. So, $X$ is a set that is non-empty, but let's not put it. And $d$ --- this is the definition, sorry, this is the definition of metric space. And $d$ is a map from the product $X$ times $X$ to the non-negative real numbers that will be called the distance. This is called the distance. And this distance satisfies three axioms.

Axiom number one is that is called is definite. That is for every pair of points $x$ and $y$ in $X$, in my space, so the distance between the two points is zero if and only if the points are the same. The distance between a point and itself is zero, of course, because what sort of distance would be non-zero from a point to itself? And if the points are different, the distance will not be zero because otherwise this distance would not be able to distinguish points, right? Axiom number two: the distance is symmetric. The distance between $x$ and $y$ is the same as the distance between $y$ and $x$. Okay? This axiom is very also very reasonable for a distance. In applications, uh, this is not always true if you think of a distance like the distance of flying --- uh, with a plane between two cities. You know that when you fly to the west is usually take less time than fly to the east? No, flying to the west takes less time than flying to the east? Due to the --- to the atmospheric drift. So, it's not always true in every situation that the distance is symmetric, but in many situations it is true. And the third axiom is the triangle inequality.
\end{spoken-clean}

\begin{spoken-clean}[00:15:15 - 00:17:15]
These are the axioms of a metric space. First example --- now I will give some examples, no? First example of a metric space: the Euclidean space $\mathbb{R}^n$ with the distance the length of a vector. Is very easy that it satisfies the --- the three axioms. So, examples. $\mathbb{R}^n$ with the Euclidean distance $d_{\text{Euclidean}}$. Why do we specify both the set and the distance? Because sometimes the same set has two or more interesting distances, right? For instance, in $\mathbb{R}^2$ I can put the Euclidean distance or I can put a different distance that you can check also satisfies that is called the --- in $\mathbb{R}^2$ there is this distance, uh, the New York distance (i.e., the Manhattan distance). That is mod $x_1$ minus $y_1$ plus mod $x_2$ minus $y_2$. Right? This is called the Manhattan distance. Why? Because this is the sort of distance when you live in one of these modern American cities or --- where the streets go north to south and east to west and you want to go from point $x$ to $y$, right? Then you cannot go diagonal because you hit the buildings and then you need to go like that. And then the distance between two points is --- is this thing. Okay? And you can check very easily that this still satisfies all the axioms of the distance. Right?
\end{spoken-clean}

\begin{math-stroke}[Examples of Metrics on \texorpdfstring{$\mathbb{R}^n$}{Rn}]
\begin{enumerate}
    \setcounter{enumi}{0} \item \emph{Euclidean Metric:} The standard distance in $\mathbb{R}^n$:
    \[ d_2(x, y) = \|x - y\| = \sqrt{\sum_{i=1}^n (x_i - y_i)^2} \]
    \setcounter{enumi}{1} \item \emph{Manhattan Metric (Taxicab Metric):} Common in grid-based layouts:
    \[ d_1(x, y) = |x_1 - y_1| + |x_2 - y_2| \]
\end{enumerate}

\begin{center}
\begin{tikzpicture}[scale=0.8]
% \begin{ai-tikz-planner-invisible-content}
% 1. Background: A grid representing city blocks.
% 2. Midground: Points x and y.
% 3. Foreground: The L1 path (Manhattan) and the L2 path (Euclidean).
% \end{ai-tikz-planner-invisible-content}
    % Grid
    \draw[step=0.5, gray!30, thin] (0,0) grid (4,3);
    
    % Points
    \coordinate (X) at (0.5, 0.5);
    \coordinate (Y) at (3.5, 2.5);
    \fill (X) circle (2pt) node[below left] {$x$};
    \fill (Y) circle (2pt) node[above right] {$y$};
    
    % Euclidean Distance
    \draw[thick, MidnightBlue, dashed] (X) -- (Y) node[midway, above left] {$d_2$};
    
    % Manhattan Distance
    \draw[very thick, BrickRed] (0.5, 0.5) -- (3.5, 0.5) -- (3.5, 2.5);
    \node[BrickRed, below] at (2, 0.5) {$d_1$};
\end{tikzpicture}
\end{center}
\end{math-stroke}

\begin{spoken-clean}[00:17:15 - 00:19:15]
So, when you have a bit more abstract example, when you have any metric space and you have a subset of it, you can just restrict this distance to the --- to $Y$. So, $d$ restricted to $Y$ cross $Y$ with $Y$, this is another metric space. Right? If you take a subset of your given metric space and restrict the distance, is --- will be always a metric space. So, example, more --- this is a bit always true, a more concrete example would be if this is $\mathbb{R}^3$ with the Euclidean distance, right? And I take a subset that is the two-dimensional sphere. Let me draw it like this. $S^2$ is the set of points $x$ in $\mathbb{R}^3$ such that the length of $x$ is one. Right? This $S^2$, the sphere. Then I can see the sphere as a metric space with the Euclidean distance. So, the distance between these two points of the sphere, this one and this one say, is just you take the straight segment and measure the length. And this satisfies the axioms of the distance. Is a restriction of the Euclidean distance. You can say \qt{okay, but this distance maybe is not the best}. Because as I was saying before, if you travel by plane, this is not the distance between points. The distance between points --- a first approximation of the distance between points when you travel by plane is the geodesic distance. That is the length of the arc, the shortest arc on the surface of the earth. This is another distance. Right? Same set can have multiple meaningful distances depending on the context.
\end{spoken-clean}

\begin{math-stroke}[Subspaces and the Sphere]
\begin{ai-type-check-invisible-content}
% (X, d) is a metric space. Y \subset X.
% d|_Y is the restriction of d to Y x Y.
\end{ai-type-check-invisible-content}
Any subset $Y \subset X$ of a metric space $(X, d)$ inherits a metric structure by restricting the distance function to $Y \times Y$.

\emph{Example: The Unit Sphere $S^2 \subset \mathbb{R}^3$:}
\[ S^2 = \{ x \in \mathbb{R}^3 \mid \|x\| = 1 \} \]
The sphere can be equipped with two distinct metrics:
\begin{description}
    \item[Inherited Euclidean Metric:] The straight-line distance through the interior of the sphere.
    \item[Geodesic Metric:] The distance measured along the surface of the sphere (great-circle distance).
\end{description}

\begin{center}
\begin{tikzpicture}[scale=1.5]
    % Sphere
    \draw[thick, gray!50] (0,0) circle (1);
    \draw[dashed, gray!50] (0,0) ellipse (1 and 0.3);
    
    % Points
    \coordinate (A) at (-0.7, 0.4);
    \coordinate (B) at (0.6, 0.2);
    \fill (A) circle (1.5pt);
    \fill (B) circle (1.5pt);
    
    % Euclidean distance (chord)
    \draw[thick, MidnightBlue] (A) -- (B);
    
    % Geodesic distance (arc)
    \draw[thick, BrickRed, bend left=20] (A) to (B);
    
    \node[MidnightBlue, below] at (0, 0.2) {\footnotesize Chord};
    \node[BrickRed, above] at (0, 0.5) {\footnotesize Arc};
\end{tikzpicture}
\end{center}
\end{math-stroke}

\begin{spoken-clean}[00:19:15 - 00:21:15]
And now maybe before --- let me do a couple more of examples. I will put it here. Maybe I need to erase something. So, one thing that can be useful when you try to do a proof in this setup that is a bit abstract of metric spaces, remember you have very few axioms you can use, but then you have less intuition. So, how you --- how can you overcome your lack of intuition? What you do is you --- you think of a specific example. Don't think of a general metric space. A general metric space means a billion of things at the same time, right? Is an abstract model of many things at the same time. Pick a concrete example. Pick the Euclidean space. Do your proof in the Euclidean space using your intuitions that you have from the space. And then try to see which axioms are you using. If you can use --- maybe you use the wrong proof and you are not using just those axioms. But then try to modify your proof using your intuition until you can do with those axioms. Okay? This can be useful. So, whenever things are abstract, don't think always in the abstract, just go to concrete examples. And this applies always because I have seen many times in more advanced courses doing examinations that you ask --- so people know the abstract theory very well and then you ask in a concrete case. \qt{So, give me the --- no, no, don't tell me the Ricci curvature on a --- no, what is in the Euclidean space?} And then people sometimes don't know how to answer. And this is a problem. So, the abstract is good but you need to know very well and think --- spend time thinking about concrete examples.

So, let me give you so that you can see also that things start to be interesting very soon. Let me give you some more interesting examples of metric space, right? Now we saw spaces of points. We can also consider spaces of functions. Now, you learned what is a continuous function. Let's say $a, b$ are two real numbers, $a < b$. And you can take $X$ the space of continuous functions from the closed interval $[a, b]$ to $\mathbb{R}$. $f$ continuous. This is a space of functions. So, one point in this space would look something like this... \inlinemetanote{Draws a curve on a coordinate system} This is one point in this space. How can we define a distance between functions? Well, in many ways. Something that is used often is the supremum distance or the maximum distance. Is the distance between $f$ and $g$ is the maximum in $[a, b]$ of the absolute value of the difference. So, you take two curves here, different, and then you look at the maximal oscillation --- so the maximum difference between the two functions. This is a distance between functions. If this is zero, the two functions are equal? Yes. Only if? Yes, is a bit more you need to reason but they are continuous so you can do that. Exercise. Does it satisfy the triangle inequality? Prove it. The symmetry is very easy that you replace $f$ and $g$ is the same because here you have the absolute value. But the triangle inequality you can prove it using the triangle inequality for the absolute value and playing with the maximum.

Another distance between functions... \inlinemetanote{Writes the integral formula} You measure the distance, you square, you integrate and you take the square root. This is another distance. Very important, the $L^2$ distance. We will not use this in a course, but you will see that you will use this later on many, many times. So, we will continue next time from --- from here.
\end{spoken-clean}

\begin{math-stroke}[Metrics on Function Spaces]
Let $X = C([a, b], \mathbb{R})$ be the set of all continuous functions $f: [a, b] \to \mathbb{R}$. We can define several metrics on this set:
\begin{enumerate}
    \setcounter{enumi}{0} \item \emph{Supremum Metric (\texorpdfstring{$L^\infty$}{L-infinity} Metric):}
    \[ d_\infty(f, g) = \max_{x \in [a, b]} |f(x) - g(x)| \]
    \setcounter{enumi}{1} \item \emph{\texorpdfstring{$L^2$}{L2} Metric:}
    \[ d_2(f, g) = \left( \int_a^b (f(x) - g(x))^2 \, dx \right)^{1/2} \]
\end{enumerate}

\begin{center}
\begin{tikzpicture}[scale=1.2]
% \begin{ai-tikz-planner-invisible-content}
% 1. Background: Axes.
% 2. Midground: Two continuous functions f and g.
% 3. Foreground: A vertical bar showing the maximum difference for d_infinity.
% \end{ai-tikz-planner-invisible-content}
    % Axes
    \draw[->] (0,0) -- (4,0) node[right] {$x$};
    \draw[->] (0,0) -- (0,2.5) node[above] {$y$};
    \node[below] at (0.5, 0) {$a$};
    \node[below] at (3.5, 0) {$b$};

    % Functions
    \draw[thick, MidnightBlue] (0.5, 1) .. controls (1.5, 2) and (2.5, 0.5) .. (3.5, 1.5) node[right] {$f$};
    \draw[thick, BrickRed] (0.5, 0.5) .. controls (1.5, 1) and (2.5, 1.5) .. (3.5, 0.8) node[right] {$g$};
    
    % Max difference
    \draw[<->, thick, ForestGreen] (1.2, 1.55) -- (1.2, 0.85) node[midway, left] {\footnotesize $d_\infty$};
\end{tikzpicture}
\end{center}

\begin{explanation-of-steps}
The $L^\infty$ metric measures the \qt{worst-case} difference between two functions, while the $L^2$ metric provides an \qt{average} measure of their difference over the entire interval. Both satisfy the metric axioms, demonstrating that the concept of distance extends far beyond simple points in space.
\end{explanation-of-steps}
\end{math-stroke}

\include{content/02-18-wednesday-copy-1}

\lecturechapter{Friday}{Feb 20th}{February 20th 2026}{Topology of Metric Spaces: Open and Closed Sets}

\begin{nice-box}[Context]
This lecture marks the transition from the basic definitions of metric spaces to the study of their \emph{topology}. The lecturer, Joaquim Serra, reviews the concepts of convergence and completeness before introducing the fundamental building blocks of topological spaces: open and closed sets. The discussion emphasizes the use of axioms to prove properties that will later be applied to multivariable differentiation and integration.
\end{nice-box}

% ==========================================
% SECTION 1: REVIEW AND INTRODUCTION
% ==========================================
\section{Review of Previous Lectures}

\begin{spoken-clean}[00:00:00 - 00:01:51]
We can continue. Uh, good afternoon. We can --- we can continue where we left it last time. So, I remind --- so just a quick summary of what we saw the last lectures. First lecture was mostly introduction, but then we saw metric spaces. And I told you --- because people already asked me, even though I thought the first day --- why we studied metric spaces. So, the model --- whenever it gets too abstract, the model you need to have in mind all the time for this $(X, d)$ metric space is $\mathbb{R}^n$ with the Euclidean distance, okay? So, whenever you have doubts, just plug $\mathbb{R}^n$ where there is an $X$ (a capital $X$) and $d$ is the usual --- the modulus of the difference of the points. And, um, I said that I wanted to do this in this more general framework because later, in the last part of the course, we will use other spaces. Okay?
\end{spoken-clean}

\begin{math-stroke}[Summary of Metric Space Foundations]
\begin{ai-type-check-invisible-content}
% (X, d) is a metric space.
% x_n is a sequence in X.
\end{ai-type-check-invisible-content}
The core structure studied thus far is the metric space $(X, d)$.
\begin{itemize}
    \item \textbf{Standard Model:} $(\mathbb{R}^n, d_{\text{Eucl}})$, where $d(x, y) = \|x - y\|$.
    \item \textbf{Sequences:} $(x_n)_{n \ge 0} \subset X$ converges to $x \in X$ ($x_n \to x$) if $d(x_n, x) \to 0$.
    \item \textbf{Completeness:} A space is \emph{complete} if every Cauchy sequence converges. We proved that $\mathbb{R}^n$ is a complete metric space.
\end{itemize}
\end{math-stroke}

\begin{spoken-clean}[00:01:51 - 00:02:22]
And for good reason. It's not --- I mean, we will use to prove some theorems in the --- in the course, the space of continuous functions. Uh, but --- if you want for the --- this is not something that belongs to mathematicians, right? Let me say very quickly for the physicists: so in classical mechanics, the points in the space is a vector in $\mathbb{R}^n$, right? But when you --- you will go to quantum mechanics, what is a point? What is the position of a particle in space? Is the wave function. This is a function. So, the function --- the wave functions are naturally points in your space, right? So, this is what we are saying here. We are trying to say things more general to be able to understand, um, the things that come naturally in the physics, for instance.
\end{spoken-clean}

\begin{didactic-insight}[The Physical Motivation for Abstraction]
The lecturer provides a profound justification for studying abstract metric spaces. In classical mechanics, a "point" is a coordinate in $\mathbb{R}^n$. In quantum mechanics, a "state" is a wave function. By treating functions as points in a metric space, the tools of Analysis (limits, continuity, topology) become directly applicable to the most advanced models of physics.
\end{didactic-insight}

\begin{spoken-clean}[00:02:22 - 00:03:13]
Okay. So, now, in this framework of metric spaces, we saw what it means a sequence, what it means that a sequence converges, what is a Cauchy sequence, what is a complete metric space, and we proved that $\mathbb{R}^n$ is always a complete metric space using that $\mathbb{R}$ is complete. This last time. And today we will see, uh, what is the definition of open and closed sets in a metric space and the notion of continuous functions in the setup of metric spaces. So, let me start with the definition of open set. So, let me start maybe with the definition of open ball. Open ball.
\end{spoken-clean}

% ==========================================
% SECTION 2: TOPOLOGY OF METRIC SPACES
% ==========================================
\setcounter{section}{1}
\section{Topology of Metric Spaces}
\subsection{Open and Closed Sets}

\begin{math-stroke}[Definition: Open Ball]
\setcounter{theorem}{37}
\begin{definition}[Open Ball]\label{def:open-ball}
We define the \emph{open ball} of radius $r > 0$ centered at a point $x \in X$ in a metric space $(X, d)$ as:
\[ B(x, r) = \{ y \in X \mid d(y, x) < r \} \]
\end{definition}

\begin{explanation-of-steps}
The open ball is the set of all points whose distance to the center $x$ is strictly less than the radius $r$. The strict inequality ($<$) is what makes the ball \qt{open}—it does not include its boundary.
\end{explanation-of-steps}

\begin{nice-box}[Notation Note]
The lecturer notes that $B_r(x)$ is an equivalent and common notation for $B(x, r)$ found in many textbooks.
\end{nice-box}
\end{math-stroke}

\begin{spoken-clean}[00:03:13 - 00:05:12]
So, we --- we define --- define the open ball, open ball of radius $r$ positive centered at a given point $x$ in my metric space as --- and we will denote with this: is ball centered at $x$ with radius $r$. Sometimes you will see equivalent notation is the ball of radius $r$ centered at $x$. This is also in many notes and books. Is the same. Okay? This is by definition the set of points in my space such that the distance --- sorry, the points will be $y$ --- such that the distance between $y$ and $x$ is less than $r$, right?

So, if you want a drawing of that, in the Euclidean space --- in a --- in a general metric space I cannot draw anything because it's too general --- but in the Euclidean space, let's say in $\mathbb{R}^2$, the metric ball would be the ball. So, the --- the open ball would be this, right? \inlinemetanote{Draws a circle on the board} What is inside of the --- of the ball. And it's open because you don't include --- if you want I can --- I can put the boundary of the ball, that we will introduce in a --- today also in a --- in a general topological space what the boundary means. This surface of the ball is not included, because there is a strict inequality, okay? The points that are at exactly distance $r$, I don't take them. This is the open ball.
\end{spoken-clean}

\begin{math-stroke}[Visualizing the Open Ball in \texorpdfstring{$\mathbb{R}^2$}{R2}]
\begin{center}
\begin{tikzpicture}[scale=1.5]
% \begin{ai-tikz-planner-invisible-content}
% 1. Background: 2D Axes.
% 2. Midground: A dashed circle representing the boundary.
% 3. Foreground: Shaded interior and center point x.
% \end{ai-tikz-planner-invisible-content}
    % Axes
    \draw[->, gray!50] (-1.5,0) -- (1.5,0) node[right] {$x_1$};
    \draw[->, gray!50] (0,-1.5) -- (0,1.5) node[above] {$x_2$};

    % The Ball
    \draw[dashed, thick, MidnightBlue, fill=MidnightBlue!10] (0,0) circle (1cm);
    \fill (0,0) circle (1.5pt) node[below right] {$x$};
    \draw[->, thick, ForestGreen] (0,0) -- (0.707, 0.707) node[midway, above left] {$r$};
    
    \node[MidnightBlue] at (1.2, -0.8) {$B(x, r)$};
    \node[gray, font=\footnotesize] at (1.2, 1.0) {Boundary (excluded)};
\end{tikzpicture}
\end{center}
\end{math-stroke}

\begin{spoken-clean}[00:05:12 - 00:06:15]
Now, next definition is open and closed sets. So, we say that a set $U$ of my metric space $X$ is open if for every $x$ in $U$, exists some positive radius $r$ (depending on $x$, possibly) such that the ball of radius $r$ centered at $x$ is contained in $U$. So, I will draw some in $\mathbb{R}^2$ some open set. Is essentially any --- any set --- this is a bit --- heuristically, this is the idea, right? Then we will see concrete open sets. But is a set that does not include its contour. So, its boundary. Again, I --- I didn't define the boundary yet, but now you know --- you --- you have some common sense knowledge of what is a boundary of this set, that is this part --- is the skin of the set, right? If the set is the inside without the skin, this is open.
\end{spoken-clean}

\begin{math-stroke}[Definition: Open and Closed Sets]
\setcounter{theorem}{38}
\begin{definition}[Open Set]\label{def:open-set}
A subset $U \subseteq X$ of a metric space $(X, d)$ is \emph{open} if every point in $U$ is the center of some open ball entirely contained within $U$:
\[ \forall x \in U, \exists r > 0 \text{ s.t. } B(x, r) \subseteq U \]
\end{definition}

\begin{center}
\begin{tikzpicture}[scale=1.2]
% \begin{ai-tikz-planner-invisible-content}
% 1. Background: A generic "potato" shape with a dashed boundary.
% 2. Midground: A point x inside.
% 3. Foreground: A small solid circle around x representing B(x, r).
% \end{ai-tikz-planner-invisible-content}
    % The Open Set U
    \draw[dashed, thick, fill=MidnightBlue!5] 
        (0,0) to[out=90,in=180] (1.5,1.5) 
        to[out=0,in=90] (3,0.5) 
        to[out=-90,in=0] (1.5,-1) 
        to[out=180,in=-90] cycle;
    \node at (2.5, 1.2) {\Large $U$};

    % Point x and its ball
    \coordinate (X) at (1, 0.2);
    \fill (X) circle (1.5pt) node[below] {$x$};
    \draw[thick, ForestGreen] (X) circle (0.4cm);
    \node[ForestGreen, right] at (1.4, 0.2) {\footnotesize $B(x, r) \subset U$};
\end{tikzpicture}
\end{center}

\begin{explanation-of-steps}
The intuition behind an open set is that no point in the set is on the \qt{edge}. For any point $x$, you can always move a tiny distance in any direction and still remain inside the set.
\end{explanation-of-steps}
\end{math-stroke}

% \begin{ai-global-state-checkpoint-invisible-content}
% timestamp: 00:06:00
% topic: Definitions of open balls and open sets.
% board_state: def:open-ball, def:open-set
% next_goal: Provide a concrete example of an open set (the square).
% open_loops: none
% \end{ai-global-state-checkpoint-invisible-content}

\begin{spoken-clean}[00:06:15 - 00:08:37]
Why? Because any point here, you can take a ball --- depending on the distance to the boundary, right? Depending on the distance to the skin. You make the ball super small and it will be contained in the set. Is --- as I said, any point in the interior will be at positive distance from the skin.

So, if you want a more concrete example, take the set $(0, 1) \times (0, 1)$, the square, open square, subset $\mathbb{R}^2$. This --- prove that it is open. Exercise. So, you take this square... \inlinemetanote{Draws a dashed square on the board} You take a point here. A point in the square will have two coordinates. And now, what is the definition of open? If you take the point --- this point is $(x_1, x_2)$ --- you need to find a --- you need to construct a radius $r$ positive such that the whole ball centered at this given point belongs to the --- to the square. In the case of the square, what is the formula for the radius? Well, you can take $r$ equal, I don't know, something like the minimum between this distance, this distance, this distance, and this distance. So, it will be something like the minimum between $x_1, x_2, 1 - x_1$ and $1 - x_2$. You can check that this radius works. The minimum of these four numbers. Right?
\end{spoken-clean}

\begin{math-stroke}[Example: The Open Unit Square]
Consider the open unit square in $\mathbb{R}^2$:
\[ U = (0, 1) \times (0, 1) \subset \mathbb{R}^2 \]
\begin{center}
\begin{tikzpicture}[scale=2]
% \begin{ai-tikz-planner-invisible-content}
% 1. Background: Dashed square boundary.
% 2. Midground: Point x = (x1, x2).
% 3. Foreground: Arrows showing distances to the four boundaries.
% \end{ai-tikz-planner-invisible-content}
    % The Square
    \draw[dashed, thick] (0,0) rectangle (1.5, 1.5);
    \node[below left] at (0,0) {$0$};
    \node[below] at (1.5, 0) {$1$};
    \node[left] at (0, 1.5) {$1$};

    % Point x
    \coordinate (X) at (0.4, 0.6);
    \fill (X) circle (1pt) node[above right] {$(x_1, x_2)$};

    % Distances to boundaries
    \draw[<->, BrickRed, thin] (0.4, 0) -- (0.4, 0.6) node[midway, left] {$x_2$};
    \draw[<->, BrickRed, thin] (0, 0.6) -- (0.4, 0.6) node[midway, above] {$x_1$};
    \draw[<->, BrickRed, thin] (0.4, 0.6) -- (0.4, 1.5) node[midway, left] {$1-x_2$};
    \draw[<->, BrickRed, thin] (0.4, 0.6) -- (1.5, 0.6) node[midway, above] {$1-x_1$};
    
    % The Ball
    \draw[ForestGreen, thick] (X) circle (0.25cm);
    \node[ForestGreen, below] at (0.4, 0.35) {\footnotesize $B(x, r)$};
\end{tikzpicture}
\end{center}
\begin{explanation-of-steps}
To prove $U$ is open, for any $(x_1, x_2) \in U$, we must find $r > 0$ such that $B(x, r) \subset U$. The distance to the nearest boundary is given by:
\[ r = \min(x_1, x_2, 1 - x_1, 1 - x_2) \]
Since $0 < x_1, x_2 < 1$, all four terms are strictly positive, so $r > 0$. Any point $y$ within distance $r$ of $x$ will necessarily have coordinates between $0$ and $1$, remaining inside the square.
\end{explanation-of-steps}
\end{math-stroke}

\begin{spoken-clean}[00:08:37 - 00:10:16]
Okay. Second point in the definition is what is a closed set. So, $A$ subset $X$ is called closed if the complement is open. $X - A$ is open. This set $(0, 1) \times (0, 1)$ is not closed. You can show. Instead, this set --- for instance, $[0, 1] \times [0, 1]$ --- this set is closed. Now it needs to contains the skin to be closed. Because then, when you pass to the open --- to the complement that is by definition open, it need not to contain the boundary, right? So, the closed sets, we draw them with their skin. This is closed. $A$.
\end{spoken-clean}

\begin{math-stroke}[Definition: Closed Set]
\begin{definition}[Closed Set]\label{def:closed-set}
A subset $A \subseteq X$ is \emph{closed} if its complement in $X$ is an open set:
\[ A \text{ is closed} \iff X \setminus A \text{ is open} \]
\end{definition}

\begin{center}
\begin{tikzpicture}[scale=1.2]
% \begin{ai-tikz-planner-invisible-content}
% 1. Background: A solid bounding box representing X.
% 2. Midground: A solid "potato" shape representing A.
% 3. Foreground: Labels for A and X \ A.
% \end{ai-tikz-planner-invisible-content}
    % Ambient Space X
    \draw[thick] (-1,-1.5) rectangle (4, 2);
    \node at (3.5, 1.7) {$X$};

    % Closed Set A
    \draw[thick, fill=gray!20] 
        (0.5,0) to[out=90,in=180] (1.5,1) 
        to[out=0,in=90] (2.5,0) 
        to[out=-90,in=0] (1.5,-0.8) 
        to[out=180,in=-90] cycle;
    \node at (1.5, 0.1) {\Large $A$};
    
    \node at (3, -1) {$X \setminus A$ (Open)};
\end{tikzpicture}
\end{center}

\begin{explanation-of-steps}
A closed set \qt{contains its boundary}. For example, the closed unit square $[0, 1] \times [0, 1]$ is closed because its complement (everything outside the square) is an open set.
\end{explanation-of-steps}
\end{math-stroke}

\begin{spoken-clean}[00:10:16 - 00:12:03]
Okay. And maybe the last part of this definition that you will not use --- we will not use actually --- but --- but is something that is a --- a common term that you will find in books and it's good that you know, is what is the topology. So, you will see in the notes you will see the word topology many times. You are not yet studying topology but --- just for information, the word topology --- topology generated --- so the --- the topology associated to the distance $d$ is the --- we can put a $\mathcal{T}$ like this, a sort of calligraphic $\mathcal{T}$, is the collection of all the open sets. This is the topology. So, when we --- when we say that something is a topological notion, it has to do with this collection of open sets.
\end{spoken-clean}

\begin{math-stroke}[Definition: Topology]
\begin{definition}[Topology]\label{def:topology}
The \emph{topology} $\mathcal{T}$ on a metric space $(X, d)$ is the collection of all open subsets of $X$:
\[ \mathcal{T} = \{ U \subseteq X \mid U \text{ is open} \} \]
\end{definition}
\end{math-stroke}

% \begin{ai-global-state-checkpoint-invisible-content}
% timestamp: 00:12:00
% topic: Definitions of closed sets and topology.
% board_state: def:closed-set, def:topology
% next_goal: State and prove the lemma on unions and intersections of open sets.
% open_loops: none
% \end{ai-global-state-checkpoint-invisible-content}

\begin{spoken-clean}[00:12:03 - 00:13:41]
Now, to see a bit more what are --- because this is still a bit abstract --- to see a bit more practically how we identify open and closed sets, uh, we are going to do it through sequences. Okay? So, this is the... well, let me first do a lemma, maybe, that it we will use many times. Lemma: unions and intersections of open sets. This is one of the basic properties. It's --- it's very simple, but since we will use many, many times, I prove it now. So, it says that --- I will put it with words and then we will prove with symbols. So, arbitrary unions of open sets are open.
\end{spoken-clean}

\begin{nice-box}[Proposition: Unions and Intersections of Open Sets]
\setcounter{theorem}{42}
\begin{proposition}\label{prop:open-sets-unions-intersections}
In any metric space $(X, d)$:
\begin{enumerate}
    \setcounter{enumi}{0} \item The union of an \emph{arbitrary} collection of open sets is open.
    \setcounter{enumi}{1} \item The intersection of a \emph{finite} collection of open sets is open.
\end{enumerate}
\end{proposition}
\end{nice-box}

\begin{spoken-clean}[00:13:41 - 00:15:50]
So, this means that if you have some collection $\{U_i\}_{i \in I}$ of $U_i$ subset $X$ open, then the union of this collection (that can be --- the union of the collection is also open). Right? And also finite intersections of open sets are open. Right? So, this means that if $I$ is finite (is a finite set), the intersection --- let's say, let's put it like this: the intersection of $U_i$ for $i$ between $1$ and capital $N$ (a bounded number) --- this is open set.
\end{spoken-clean}

\begin{proof}[Proof of Proposition \ref{prop:open-sets-unions-intersections}]
\begin{spoken-clean}[continued]
The proof of this is just using the definition. So, if I want to check that a union is open, I take a point $x$ in this set that is the union. Let me call this in the union. And then this in particular $x$ belongs to one of them, because it's a union. I call $i_0$ is one of the $i$s. I any of them. And then since this --- this guy is open, exists some radius such that the ball of radius $r$ centered at $x$ is contained in $U_{i_0}$. But this is contained in the union. Which proves that this union is open, right?
\end{spoken-clean}

\begin{math-stroke}[Proof: Arbitrary Unions]
Let $\{U_i\}_{i \in I}$ be a collection of open sets. Let $x \in \bigcup_{i \in I} U_i$.
\begin{description}
    \item[Step 1:] By the definition of union, there exists some index $i_0 \in I$ such that $x \in U_{i_0}$.
    \item[Step 2:] Since $U_{i_0}$ is open, there exists $r > 0$ such that $B(x, r) \subseteq U_{i_0}$.
    \item[Step 3:] Since $U_{i_0} \subseteq \bigcup_{i \in I} U_i$, it follows that $B(x, r) \subseteq \bigcup_{i \in I} U_i$.
\end{description}
Thus, the union is open.
\end{math-stroke}

\begin{spoken-clean}[continued]
With the finite intersection, you can do as an exercise. But I will tell you how to do, okay? Now explain the proof with, yeah. \inlinemetanote{The lecturer retrieves the catch-box microphone}
\end{spoken-clean}

\begin{student-interaction}[Student Question]
Except for the empty set and the set $X$ itself, can you give us some examples of both closed and open sets?
\end{student-interaction}

\begin{spoken-clean}[continued]
Um, we will discuss this when we --- when we discuss connected spaces. Okay? This is a good question. So, the --- the whole space and the empty set are always open and closed, and I will emphasize this when we see connected spaces. So, you need to wait a bit more. I could give an example, but it's not very interesting. It's just a disconnected space like the --- the union, the disjoint union of two intervals or something like that, then each interval is open and closed. But this will not be very interesting. You will see the interest of this question when we see connected spaces. Okay? But --- is a fair question.
\end{spoken-clean}

\begin{math-stroke}[Proof: Finite Intersections]
\begin{ai-proof-skeleton-invisible-content}
% 1. Take x in the intersection of U_1, ..., U_N.
% 2. x is in each U_i, so there exists r_i such that B(x, r_i) is in U_i.
% 3. Take r = min(r_1, ..., r_N).
% 4. Then B(x, r) is in every U_i, so it's in the intersection.
\end{ai-proof-skeleton-invisible-content}
Let $U_1, \dots, U_N$ be open sets. Let $x \in \bigcap_{i=1}^N U_i$.
\begin{description}
    \item[Step 1:] Since $x$ is in the intersection, $x \in U_i$ for all $i \in \{1, \dots, N\}$.
    \item[Step 2:] Since each $U_i$ is open, for each $i$ there exists $r_i > 0$ such that $B(x, r_i) \subseteq U_i$.
    \item[Step 3:] Let $r = \min(r_1, r_2, \dots, r_N)$. Since we are taking the minimum of a \emph{finite} number of positive values, $r > 0$.
    \item[Step 4:] For this $r$, $B(x, r) \subseteq B(x, r_i) \subseteq U_i$ for all $i$. Therefore, $B(x, r) \subseteq \bigcap_{i=1}^N U_i$.
\end{description}
Thus, the finite intersection is open.
\end{math-stroke}
\end{proof}

% \begin{ai-global-state-checkpoint-invisible-content}
% timestamp: 00:18:00
% topic: Unions and intersections of open sets.
% board_state: prop:open-sets-unions-intersections, proof:unions-open
% next_goal: State and prove the lemma for closed sets using De Morgan's laws.
% open_loops: none
% \end{ai-global-state-checkpoint-invisible-content}

\begin{spoken-clean}[00:18:05 - 00:19:05]
So, I am saying that when --- when you take open sets, when you take finitely many intersections of open sets, uh, you get something open. Why? I will say with words, and then you try to write down, okay? So, you --- you have finitely many open sets. And now you take a point that is in the intersection of these finitely many. Right? This means that the point is in all of them. And now, for each of the sets, you pick a radius that is different in every time, but there are finitely many radii that you can put your ball in the open set. And if you take the minimum of all these radii, it will be good. With the minimum of the finitely many radii, you will fall inside of the intersection. Okay?
\end{spoken-clean}

\begin{spoken-clean}[00:19:05 - 00:21:42]
You don't need to understand what I said, please try to do the exercise. If there are problems doing this, uh, you have it in the notes. But I think this is a --- is very good to try in these first proofs of the --- of the theory that are very short one-line proofs, is very good opportunity to do as an exercise to see if we understood the definitions. Because is one line. So, it's not --- you don't need to think too much.

Okay. Now, analogous lemma but for closed sets. Unions and intersections of closed sets. And it is the same but --- the same but the opposite. So, it says that finite unions of closed sets are closed, and arbitrary intersections of closed sets are closed. And the proof is essentially --- well, you can repeat the same proof or you can reason a bit more --- or a variant of the same proof --- or you can reason passing to the complement. So, you can use that when you have a finite union --- so remember that being closed is defined as the complement being open. So, if you pass to the complement, the complementation converts closed sets in open and open in closed, right? So, you can always pass statements about open sets to the complement and you get a statement about closed sets. If you have some property of open sets and you pass to the complements (all the statement), you get a property of closed sets. And this is what we have here. Why? Because $X$ minus a union of sets is the same as the intersection of the complements, right?
\end{spoken-clean}

\begin{nice-box}[Proposition: Unions and Intersections of Closed Sets]
\setcounter{theorem}{43}
\begin{proposition}\label{prop:closed-sets-unions-intersections}
In any metric space $(X, d)$:
\begin{enumerate}
    \setcounter{enumi}{0} \item The intersection of an \emph{arbitrary} collection of closed sets is closed.
    \setcounter{enumi}{1} \item The union of a \emph{finite} collection of closed sets is closed.
\end{enumerate}
\end{proposition}
\end{nice-box}

\begin{math-stroke}[Proof via De Morgan's Laws]
\begin{proof}
The properties of closed sets follow directly from the properties of open sets using De Morgan's Laws:
\begin{align}
X \setminus \left( \bigcup_{i \in I} A_i \right) &= \bigcap_{i \in I} (X \setminus A_i) \label{eq:demorgan-union} \\
X \setminus \left( \bigcap_{i \in I} A_i \right) &= \bigcup_{i \in I} (X \setminus A_i) \label{eq:demorgan-intersection}
\end{align}
\begin{description}
    \item[Arbitrary Intersections:] Let $\{A_i\}_{i \in I}$ be closed. Then each $X \setminus A_i$ is open. By Proposition \ref{prop:open-sets-unions-intersections}, their arbitrary union $\bigcup (X \setminus A_i)$ is open. By Equation \ref{eq:demorgan-intersection}, this union is the complement of the intersection $\bigcap A_i$. Since the complement is open, the intersection is closed.
    \item[Finite Unions:] Let $A_1, \dots, A_N$ be closed. Then each $X \setminus A_i$ is open. By Proposition \ref{prop:open-sets-unions-intersections}, their finite intersection $\bigcap (X \setminus A_i)$ is open. By Equation \ref{eq:demorgan-union}, this intersection is the complement of the union $\bigcup A_i$. Since the complement is open, the union is closed.
\end{description}
\end{proof}
\end{math-stroke}

\begin{spoken-clean}[00:21:42 - 00:24:23]
So, when you do the complement of a union, you get the intersection of the complements. And the same: the complement of a intersection is equal to the union of the complements. And then you apply the previous result. Because if these sets are closed, these sets will be open. And the intersection --- so the finite intersection of open is open. Okay?

Let me give some example very --- very easy that this property --- when we said finite intersections or finite unions --- is important that it is finite. That when you get infinite intersections of open sets, you can get a closed set that is not open. Example: finite is important. So, if you take the simplest possible --- well, the simplest possible no, but one of the spaces that you are most familiar with that is $\mathbb{R}$ with the distance of the absolute value (with the Euclidean distance), and you take the sets $U_k = (-1/k, 1/k)$. And you take its intersection, then the intersection for $k$ equals $1$ to infinity of $U_k$ --- this gives you the point $\{0\}$. So, these are open intervals, and it's very easy because the intervals do not contain the end points, each of them is open. But when you do their intersection, you get just the point zero, and the point zero is not open because is --- you cannot put any ball in that contains anything, right? The point zero is closed but is not open. So, these are open and this is not open! That's why it's important that the intersection is finitely many.
\end{spoken-clean}

\begin{math-stroke}[Example: Infinite Intersection of Open Sets]
\setcounter{theorem}{44}
\begin{example}\label{ex:infinite-intersection-open}
The finiteness condition in Proposition \ref{prop:open-sets-unions-intersections} is necessary. Consider the collection of open intervals in $\mathbb{R}$:
\[ U_k = \left( -\frac{1}{k}, \frac{1}{k} \right) \quad \text{for } k \in \mathbb{N} \]
The intersection of this infinite collection is:
\[ \bigcap_{k=1}^\infty U_k = \{ 0 \} \]
\begin{center}
\begin{tikzpicture}[scale=1.5]
% \begin{ai-tikz-planner-invisible-content}
% 1. Background: Horizontal axis.
% 2. Midground: Nested intervals (-1, 1), (-1/2, 1/2), etc.
% 3. Foreground: The point {0} at the center.
% \end{ai-tikz-planner-invisible-content}
    \draw[->, thick] (-1.5,0) -- (1.5,0) node[right] {$\mathbb{R}$};
    
    % Intervals
    \draw[ultra thick, MidnightBlue!30] (-1, 0.1) -- (1, 0.1);
    \node[MidnightBlue!30, above] at (1, 0.1) {\footnotesize $U_1$};
    
    \draw[ultra thick, MidnightBlue!60] (-0.5, 0.2) -- (0.5, 0.2);
    \node[MidnightBlue!60, above] at (0.5, 0.2) {\footnotesize $U_2$};
    
    \draw[ultra thick, MidnightBlue] (-0.25, 0.3) -- (0.25, 0.3);
    \node[MidnightBlue, above] at (0.25, 0.3) {\footnotesize $U_3$};

    % The limit point
    \fill[BrickRed] (0,0) circle (2pt) node[below] {$\{0\}$};
\end{tikzpicture}
\end{center}
\begin{explanation-of-steps}
While each $U_k$ is an open set, their infinite intersection is the singleton set $\{0\}$, which is \emph{closed} but \emph{not open} in $\mathbb{R}$. This demonstrates that the property of openness is not necessarily preserved under infinite intersections.
\end{explanation-of-steps}
\end{example}
\end{math-stroke}

% \begin{ai-global-state-checkpoint-invisible-content}
% timestamp: 00:24:00
% topic: Counter-example for infinite intersections of open sets.
% board_state: ex:infinite-intersection-open
% next_goal: Define interior and closure.
% open_loops: none
% \end{ai-global-state-checkpoint-invisible-content}

\begin{spoken-clean}[00:24:23 - 00:27:06]
Okay. So, now we will define something interesting that is what we were saying there heuristically. So, what is the interior and the boundary of a set in a metric space. So, given a set --- let me call the set $\Omega$ subset $X$ --- $(X, d)$ is a metric space. I define the interior of $\Omega$. This is called the interior. Maybe I will not put parenthesis. I will say $\operatorname{int} \Omega$ or sometimes you put $\Omega$ with a point like this, or with a point here, depends on the source. The interior --- this is what is called the interior of $\Omega$. This is by definition the largest open set contained in $\Omega$. So, you have your set $\Omega$ and you look at all the open sets that are contained and you get the largest one is the union of all the open sets that contain $\Omega$. So, is the union of the $U$ subset $\Omega$ such that $U$ is open. Right? This is the interior.

Now, I also define the closure that you know from $\mathbb{R}$, right? The closure of some set is the set of points $x$ in $X$ such that there exists some sequence $x_n$ such that the whole sequence... \textit{[Audio cuts off abruptly]}
\end{spoken-clean}

\begin{math-stroke}[Definition: Interior and Closure]
\setcounter{theorem}{45}
\begin{definition}[Interior]\label{def:interior}
The \emph{interior} of a set $\Omega \subseteq X$, denoted $\operatorname{int}(\Omega)$ or $\mathring{\Omega}$, is the largest open set contained within $\Omega$. Formally, it is the union of all open sets contained in $\Omega$:
\[ \operatorname{int}(\Omega) = \bigcup \{ U \subseteq \Omega \mid U \text{ is open} \} \]
\end{definition}

\begin{definition}[Closure]\label{def:closure}
The \emph{closure} of a set $\Omega \subseteq X$, denoted $\overline{\Omega}$, is the set of all points that can be reached as limits of sequences in $\Omega$:
\[ \overline{\Omega} = \{ x \in X \mid \exists (x_n)_{n \ge 0} \subset \Omega \text{ s.t. } x_n \to x \} \]
\end{definition}
\end{math-stroke}

% [SYSTEM] Video complete.
\include{content/thursday-to-merge-part2}

\lecturechapter{Monday}{Feb 23rd}{February 23rd 2026}{Topology of Metric Spaces}

\begin{nice-box}[Context]
This lecture transitions from the basic definitions of metric spaces and sequences to the study of \emph{Topology}. The lecturer, Joaquim Serra, introduces the fundamental concepts of open and closed sets, interior, and closure. These concepts provide the language necessary to discuss continuity and limits in a more abstract and powerful framework, moving beyond the specific Euclidean structure of $\mathbb{R}^n$.
\end{nice-box}

\begin{meta-note}[Initial Board State]
The left chalkboard contains a summary of previous lectures:
\begin{itemize}
    \item $(X, d)$ metric space [$\mathbb{R}^n, d_{\text{Eucl}}$]
    \item $(x_n)_{n \ge 0} \quad x_n \to x$
    \item Cauchy seq. $\mid$ Complete metric space
\end{itemize}
Below this, the plan for today is written:
\begin{itemize}
    \item Today: Open \& closed sets
    \item Continuous functions
\end{itemize}
\end{meta-note}

% ==========================================
% SECTION 1: RECALL AND MOTIVATION
% ==========================================
\section{Recall and Motivation}

\begin{spoken-clean}[00:00:00 - 00:01:50]
We can continue. Uh, good afternoon. We can --- we can continue where we left it last time. So, I remind --- so just a quick summary of what we saw the last lectures. First lecture was mostly introduction, but then we saw metric spaces. And I told you --- because people already asked me, even though I thought the first day --- why we studied metric spaces. So, the model --- whenever it gets too abstract, the model you need to have in mind all the time for this $(X, d)$ metric space is $\mathbb{R}^n$ with the Euclidean distance, okay? So, whenever you have doubts, just plug $\mathbb{R}^n$ where there is an $X$ (a capital $X$) and $d$ is the usual --- the modulus of the difference of the points. 

And, uh, I said that I wanted to do this in this more general framework because later, in the last part of the course, we will use other spaces, okay? And for good reason. It's not --- I mean, we will use to prove some theorems in the course, the space of continuous functions. But, uh, if you want, for the --- this is not something that belongs to mathematicians, right? Let me say very quickly for the physicists: so, in classical mechanics, the points in the space is a vector in $\mathbb{R}^n$, right? But when you will go to quantum mechanics, what is a point? What is the position of a particle in space? Is the wave function. This is a function. So, the functions, the wave functions, are naturally points in your space, right? So, this is what we are saying here. We are trying to say things more general to be able to understand, um, the things that come naturally in the physics, for instance.
\end{spoken-clean}

\begin{spoken-clean}[00:01:50 - 00:02:22]
Okay. So, now, in this framework of metric spaces, we saw what it means a sequence, what it means that a sequence converges, what is a Cauchy sequence, what is a complete metric space, and we proved that $\mathbb{R}^n$ is always a complete metric space using that $\mathbb{R}$ is complete, this last time. And today we will see the definition of open and closed sets in a metric space and the notion of continuous functions in the setup of metric spaces. So, let me start with the definition of open ball.
\end{spoken-clean}

% ==========================================
% SECTION 2: OPEN BALLS AND OPEN SETS
% ==========================================
\section{Open Balls and Open Sets}

\begin{math-stroke}[Definition: Open Ball]
\setcounter{section}{2} \setcounter{subsection}{0} \setcounter{theorem}{37}
\begin{definition}[Open Ball]\label{def:open-ball-v2}
Let $(X, d)$ be a metric space. We define the \emph{open ball} of radius $r > 0$ centered at a point $x \in X$ as the set:
\[ B(x, r) = \{ y \in X \mid d(y, x) < r \} \]
\end{definition}

\begin{explanation-of-steps}
The open ball consists of all points in the space whose distance to the center $x$ is strictly less than the radius $r$. The strict inequality is what makes the ball \qt{open}—it does not include its boundary.
\end{explanation-of-steps}

\begin{ai-note}[Notation]
The lecturer notes that an equivalent notation often found in literature is $B_r(x)$.
\end{ai-note}
\end{math-stroke}

\begin{spoken-clean}[00:03:54 - 00:04:51]
So, if you want a drawing of that, in the Euclidean space --- in a general metric space I cannot draw anything because it's too general, but in the Euclidean space, let's say in $\mathbb{R}^2$, the metric ball would be the ball. So, the --- the open ball would be this, right? \inlinemetanote{Draws a circle on the board} What is inside of the --- of the ball. And it's open because you don't include --- if you want I can --- I can put the boundary of the ball, that we will introduce today also in a general topological space what the boundary means. This surface of the ball is not included, because there is a strict inequality. The points that are at exactly distance $r$, I don't take them. This is the open ball.
\end{spoken-clean}

\begin{math-stroke}[Visualizing the Open Ball in \texorpdfstring{$\mathbb{R}^2$}{R2}]
\begin{center}
\begin{tikzpicture}[scale=1.5]
% \begin{ai-tikz-planner-invisible-content}
% 1. Background: A point x.
% 2. Midground: A dashed circle representing the boundary.
% 3. Foreground: Shaded interior representing the open ball.
% \end{ai-tikz-planner-invisible-content}
    \coordinate (X) at (0,0);
    \draw[dashed, thick, MidnightBlue] (X) circle (1cm);
    \fill[MidnightBlue!10, opacity=0.6] (X) circle (1cm);
    \fill (X) circle (1.5pt) node[below] {$x$};
    \draw[->, thick, gray] (X) -- (0.707, 0.707) node[midway, above left] {$r$};
    \node[MidnightBlue] at (1.2, 0.8) {$B(x, r)$};
\end{tikzpicture}
\end{center}
\end{math-stroke}

\begin{spoken-clean}[00:04:51 - 00:05:52]
Now, next definition is open and closed sets. \inlinemetanote{Writes on the board} So, we say that a set $U$ of my metric space $X$ is open if for every $x$ in $U$, exists some positive radius $r$ (depending on $x$, possibly) such that the ball of radius $r$ centered at $x$ is contained in $U$.
\end{spoken-clean}

\begin{math-stroke}[Definition: Open and Closed Sets]
\setcounter{theorem}{38}
\begin{definition}[Open and Closed Sets]\label{def:open-closed-sets}
Let $(X, d)$ be a metric space.
\begin{itemize}
    \setcounter{enumi}{0} \item A subset $U \subseteq X$ is \emph{open} if every point in $U$ is the center of some open ball entirely contained within $U$:
    \[ \forall x \in U, \exists r > 0 \text{ s.t. } B(x, r) \subseteq U \]
    \setcounter{enumi}{1} \item A subset $A \subseteq X$ is \emph{closed} if its complement is open:
    \[ A \text{ is closed} \iff X \setminus A \text{ is open} \]
\end{itemize}
\end{definition}
\end{math-stroke}

% \begin{ai-global-state-checkpoint-invisible-content}
% timestamp: 00:06:00
% topic: Definitions of open balls, open sets, and closed sets.
% board_state: def:open-ball, def:open-closed-sets
% next_goal: Provide geometric intuition for open sets.
% open_loops: none
% \end{ai-global-state-checkpoint-invisible-content}

\begin{spoken-clean}[00:05:52 - 00:07:10]
So, I will draw some in $\mathbb{R}^2$ some open set. Is essentially any set --- this is a bit heuristically, this is the idea, right? Then we will see concrete open sets. But it's a set that doesn't --- does not include its contour, so its boundary. Again, I didn't define the boundary yet, but now you know --- you have some common sense knowledge of what is a boundary of this set, that is this part --- the skin of the set, right? If the set is the inside without the skin, this is open. Why? Because any point here, you can take a ball, depending on the distance to the boundary, right? Depending on the distance to the skin, you make the ball super small and it will be contained in the set. It's like I said: any point in the interior will be at positive distance from the skin.
\end{spoken-clean}

\begin{math-stroke}[Intuition for Open Sets]
\begin{center}
\begin{tikzpicture}[scale=1.2]
% \begin{ai-tikz-planner-invisible-content}
% 1. Background: A generic blob with a dashed boundary.
% 2. Midground: A point x near the boundary.
% 3. Foreground: A small ball B(x, r) contained within the blob.
% \end{ai-tikz-planner-invisible-content}
    % The Open Set U
    \draw[dashed, thick, fill=ForestGreen!10] 
        (0,0) to[out=90,in=180] (1.5,1.5) 
        to[out=0,in=90] (3,0.5) 
        to[out=-90,in=0] (1.5,-0.5) 
        to[out=180,in=-90] cycle;
    \node at (1.5, 0.5) {\Large $U$};

    % Point x and its ball
    \coordinate (X) at (2.4, 0.3);
    \fill (X) circle (1.5pt) node[below] {$x$};
    \draw[MidnightBlue] (X) circle (0.3cm);
    \node[MidnightBlue, right] at (2.7, 0.3) {\footnotesize $B(x, r) \subset U$};
\end{tikzpicture}
\end{center}
\begin{explanation-of-steps}
The \qt{skin} or boundary of an open set is not part of the set itself. This ensures that no matter how close a point $x$ is to the edge, there is always a tiny bit of \qt{breathing room} around it that remains entirely within $U$.
\end{explanation-of-steps}
\end{math-stroke}

\begin{spoken-clean}[00:07:10 - 00:08:37]
So, if you want a more concrete example, take the set $(0, 1) \times (0, 1)$, the square, open square, subset $\mathbb{R}^2$. This --- prove that it is open. Exercise. So, you take this square... \inlinemetanote{Draws a dashed square on the board} You take a point here. A point in the square will have two coordinates. And now, what is the definition of open? If you take the point, this point is $(x_1, x_2)$, you need to find a radius $r$ positive such that the whole ball centered at this given point belongs to the square. In the case of the square, what is the formula for the radius? Well, you can take $r$ equal, I don't know, something like the minimum between this distance, this distance, this distance, and this distance. So, it will be something like the minimum between $x_1, x_2, 1 - x_1$ and $1 - x_2$. You can check that this radius works. The minimum of these four numbers.
\end{spoken-clean}

\begin{math-stroke}[Example: The Open Unit Square]
Consider the open unit square in $\mathbb{R}^2$:
\[ U = (0, 1) \times (0, 1) = \{ (x_1, x_2) \in \mathbb{R}^2 \mid 0 < x_1 < 1, 0 < x_2 < 1 \} \]
\begin{center}
\begin{tikzpicture}[scale=2]
    % Square
    \draw[dashed, thick, fill=gray!5] (0,0) rectangle (1,1);
    \node[below left] at (0,0) {$0$};
    \node[below right] at (1,0) {$1$};
    \node[above left] at (0,1) {$1$};

    % Point x
    \coordinate (X) at (0.7, 0.4);
    \fill (X) circle (1pt) node[right] {$(x_1, x_2)$};
    
    % Distances to boundaries
    \draw[<->, BrickRed] (0.7, 0) -- (0.7, 0.4) node[midway, right] {\footnotesize $x_2$};
    \draw[<->, BrickRed] (0, 0.4) -- (0.7, 0.4) node[midway, above] {\footnotesize $x_1$};
    \draw[<->, BrickRed] (0.7, 0.4) -- (1, 0.4) node[midway, above] {\footnotesize $1-x_1$};
    \draw[<->, BrickRed] (0.7, 0.4) -- (0.7, 1) node[midway, right] {\footnotesize $1-x_2$};
\end{tikzpicture}
\end{center}
\begin{explanation-of-steps}
To prove $U$ is open, for any $(x_1, x_2) \in U$, we choose the radius:
\[ r = \min(x_1, x_2, 1 - x_1, 1 - x_2) \]
Since $0 < x_1, x_2 < 1$, all four values are strictly positive, so $r > 0$. The ball $B((x_1, x_2), r)$ is then guaranteed to be contained within the square.
\end{explanation-of-steps}
\end{math-stroke}

\begin{spoken-clean}[00:08:37 - 00:10:16]
Okay, second point in the definition is what is a closed set. \inlinemetanote{Writes on the board} So, $A$ subset $X$ is called closed if the complement is open: $X \setminus A$ is open. So, now, this set --- this set is not closed, you can show. Instead, this set... \inlinemetanote{Writes $[0, 1] \times [0, 1]$ on the board} for instance, this set is closed. Now it needs to contain the skin to be closed. Because then, when you pass to the complement, that is by definition open, it need not to contain the boundary. So, the closed sets, we draw them with their skin.
\end{spoken-clean}

\begin{math-stroke}[Example: The Closed Unit Square]
The closed unit square is defined as:
\[ A = [0, 1] \times [0, 1] = \{ (x_1, x_2) \in \mathbb{R}^2 \mid 0 \le x_1 \le 1, 0 \le x_2 \le 1 \} \]
\begin{center}
\begin{tikzpicture}[scale=1.5]
    \draw[thick, fill=MidnightBlue!20] (0,0) rectangle (1,1);
    \node at (0.5, 0.5) {$A$};
    \node[below] at (0.5, -0.2) {Closed Set (includes boundary)};
\end{tikzpicture}
\end{center}
\begin{explanation-of-steps}
The set $A$ is closed because its complement, $\mathbb{R}^2 \setminus A$, is open. If you take any point outside the closed square, there is a minimum distance to the square's boundary, allowing you to fit a small open ball around that point that does not touch $A$.
\end{explanation-of-steps}
\end{math-stroke}

% \begin{ai-global-state-checkpoint-invisible-content}
% timestamp: 00:12:00
% topic: Definitions of open/closed sets and the concept of topology.
% board_state: def:open-closed-sets, ex:open-square, ex:closed-square
% next_goal: Define the topology associated with a metric.
% open_loops: none
% \end{ai-global-state-checkpoint-invisible-content}

\begin{spoken-clean}[00:10:16 - 00:11:35]
Okay, maybe this I can erase. And maybe the last part of this definition, that you will not use --- we will not use actually, but --- but is something that is a common term that you will find in books and it's good that we know, is what is the topology. So, you will see in the notes you will see the word topology many times. You are not yet studying topology, but --- just for information, the word topology, topology associated to the distance $d$ is the --- we can put a $\mathcal{T}$ like this, a sort of calligraphic $\mathcal{T}$, is the collection of all the open sets. This is the topology. So, when we say that something is a topological notion, it has to do with this collection of open sets.
\end{spoken-clean}

\begin{math-stroke}[The Topology of a Metric Space]
\begin{definition}[Topology]\label{def:topology-v2}
The \emph{topology} $\mathcal{T}$ associated with a metric space $(X, d)$ is the collection of all open subsets of $X$:
\[ \mathcal{T} = \{ U \subseteq X \mid U \text{ is open} \} \]
\end{definition}
\end{math-stroke}

% ==========================================
% SECTION 3: PROPERTIES OF OPEN SETS
% ==========================================
\section{Properties of Open Sets}

\begin{spoken-clean}[00:11:35 - 00:13:43]
Now, to see a bit more what are --- because this is still a bit abstract --- to see a bit more practically how we identify open and closed sets, we are going to do it through sequences. Well, let me first do a lemma, maybe, that we will use many times. Lemma: unions and intersections of open sets. This is one of the basic properties. It's very simple, but since we will use many, many times, I prove it now. So, it says that if you have --- I will put it with words and then we will prove with symbols. So, arbitrary unions of open sets are open. This means that if you have some collection $\{U_i\}_{i \in I}$ of $U_i \subseteq X$ open, then the union of this collection (that can be the union of the collection) is also open. Right?
\end{spoken-clean}

\begin{nice-box}[Proposition: Unions and Intersections of Open Sets]
\setcounter{theorem}{42}
\begin{proposition}[Unions and Intersections of Open Sets]
  \label[proposition]{prop:open-sets-properties}
In any metric space $(X, d)$:
\begin{enumerate}
    \setcounter{enumi}{0} \item \emph{Arbitrary Unions:} The union of any collection of open sets is open.
    \label[proposition]{prop:open-sets-properties:1:arbitrary-unions}
    \setcounter{enumi}{1} \item \emph{Finite Intersections:} The intersection of a finite number of open sets is open. 
    \label[proposition]{prop:open-sets-properties:2:finite-intersections}
\end{enumerate}
\end{proposition}
\end{nice-box}

\begin{proof}[Proof of Proposition \ref{prop:open-sets-properties} (Part 1)]
\begin{spoken-clean}[00:13:43 - 00:15:50]
And also finite intersections of open sets are open. Right? So, this means that if $I$ is finite, the intersection --- let's say, let's put it like this: the intersection of $U_i$ for $i$ between $1$ and capital $N$, a bounded number, this is open set. 

The proof of this is just using the definition. So, if I want to check that a union is open, I take a point $x$ in this set that is the union. Let me call this in the union. Right? And then, in particular, $x$ belongs to one of them because it's a union. I call $i_0$ is one of the $i$s. I fix any of them. And then, since this one is open, exists some radius $r$ such that the ball of radius $r$ centered at $x$ is contained in $U_{i_0}$. But this is contained in the union. Which proves that this union is open, right?
\end{spoken-clean}

\begin{math-stroke}[Proof: Arbitrary Unions]
Let $\{U_i\}_{i \in I}$ be a collection of open sets. Let $U = \bigcup_{i \in I} U_i$.
Take any point $x \in U$. By the definition of union:
\[ x \in U \implies \exists i_0 \in I \text{ s.t. } x \in U_{i_0} \]
Since $U_{i_0}$ is open, by definition:
\[ \exists r > 0 \text{ s.t. } B(x, r) \subseteq U_{i_0} \]
Since $U_{i_0} \subseteq \bigcup_{i \in I} U_i = U$, it follows that:
\[ B(x, r) \subseteq U \]
Thus, $U$ is open.
\end{math-stroke}
\end{proof}

% \begin{ai-global-state-checkpoint-invisible-content}
% timestamp: 00:18:00
% topic: Properties of open sets and student interaction on "clopen" sets.
% board_state: prop:open-sets-properties, proof:unions
% next_goal: Address student question about sets that are both open and closed.
% open_loops: none
% \end{ai-global-state-checkpoint-invisible-content}

\begin{student-interaction}[Student Question]
Except for the empty set and the set $X$ itself, can you give us some examples of both closed and open sets?
\end{student-interaction}

\begin{spoken-clean}[00:18:16 - 00:19:15]
We will discuss this when we --- when we discuss connected spaces. Okay? This is a good question. So, the --- the whole space and the empty set are always open and closed, and I will emphasize this when we see connected spaces. So, you need to wait a bit more. I could give an example, but it's not very interesting. It's just a disconnected space like the union, the disjoint union of two intervals or something like that, then each interval is open and closed. But this will not be very interesting. You will see the interest of this question when we see connected spaces. Okay? But it's a fair question.
\end{spoken-clean}

\begin{didactic-insight}[Clopen Sets and Connectivity]
The lecturer hints at a deep topological property: in a \qt{connected} space (like $\mathbb{R}^n$), the only sets that are both open and closed (often called \qt{clopen}) are the empty set and the space itself. The existence of other clopen sets is the defining characteristic of a disconnected space.
\end{didactic-insight}

% ==========================================
% SECTION 4: PROPERTIES OF CLOSED SETS
% ==========================================
\section{Properties of Closed Sets}

\begin{spoken-clean}[00:19:15 - 00:21:15]
So, I am saying that when you take open sets, when you take finitely many intersections of open sets, you get something open. Why? I will say with words and then you try to write down, okay? So, you have finitely many open sets. And now you take a point that is in the intersection of these finitely many. This means that the point is in all of them. And now, for each of the sets, you pick a radius that is different in every time. But there are finitely many radii that you can put your ball in the open set. And if you take the minimum of all these radii, it will be good. With the minimum of the finitely many radii, you will fall inside of the intersection. Okay? You don't need to understand what I said, please try to do the exercise. If there are problems doing this, you have it in the notes. But I think this is a very good opportunity to do as an exercise to see if we understood the definitions. Because it's one line, so it's not --- you don't need to think too much.
\end{spoken-clean}

\begin{math-stroke}[Proof Sketch of \cref{prop:open-sets-properties} \labelcref{prop:open-sets-properties:2:finite-intersections}: Finite Intersections of Open Sets]
\begin{ai-proof-skeleton-invisible-content}
% 1. Take x in the intersection of U_1, ..., U_N.
% 2. For each i, x is in U_i, so there exists r_i > 0 s.t. B(x, r_i) is in U_i.
% 3. Let r = min(r_1, ..., r_N). Since N is finite, r > 0.
% 4. Then B(x, r) is in B(x, r_i) for all i, so B(x, r) is in the intersection.
\end{ai-proof-skeleton-invisible-content}
\begin{short-proof}[Proof Sketch]
Let $U = \bigcap_{i=1}^N U_i$ where each $U_i$ is open. Take $x \in U$.
Then $x \in U_i$ for all $i \in \{1, \dots, N\}$. Since each $U_i$ is open:
\[ \forall i, \exists r_i > 0 \text{ such that } B(x, r_i) \subseteq U_i \]
Let $r = \min(r_1, r_2, \dots, r_N)$. Because we are taking the minimum of a \emph{finite} set of positive numbers, $r > 0$.
Then for all $i$, $B(x, r) \subseteq B(x, r_i) \subseteq U_i$.
Therefore, $B(x, r) \subseteq \bigcap_{i=1}^N U_i = U$, proving that $U$ is open.
\end{short-proof}
\end{math-stroke}

\begin{spoken-clean}[00:21:15 - 00:24:04]
Okay. Now, analogous lemma but for closed sets. \inlinemetanote{Writes on the board} Unions and intersections of closed sets. And it's the same but the opposite. So, it says that finite unions of closed sets are closed and arbitrary intersections of closed sets are closed. And the proof is --- well, you can repeat the same proof or you can reason a bit more --- or a variant of the same proof, or you can reason passing to the complement. So, you can use that being closed is defined as the complement being open. So, if you pass to the complement, the complementation converts closed sets in open and open in closed. So, you can always pass statements about open sets to the complement and you get a statement about closed sets. And this is what we have here. Why? Because $X$ minus a union of sets is the same as the intersection of the complements. Right? So, when you do the complement of a union, you get the intersection of the complements. And the same, the complement of a intersection is equal to the union of the complements. And then you apply the previous result. Because if these sets are closed, these sets will be open, and the finite intersection of open is open.
\end{spoken-clean}

\begin{nice-box}[Lemma: Unions and Intersections of Closed Sets]
\setcounter{theorem}{43}
\begin{proposition}[Unions and Intersections of Closed Sets]\label{prop:closed-sets-properties}
In any metric space $(X, d)$:
\begin{enumerate}
    \setcounter{enumi}{0} \item \emph{Finite Unions:} The union of a finite number of closed sets is closed.
    \setcounter{enumi}{1} \item \emph{Arbitrary Intersections:} The intersection of any collection of closed sets is closed.
\end{enumerate}
\end{proposition}
\end{nice-box}

\begin{math-stroke}[Proof via De Morgan's Laws]
The properties of closed sets follow directly from those of open sets using De Morgan's Laws:
\begin{align*}
X \setminus \left( \bigcup_{i \in I} A_i \right) &= \bigcap_{i \in I} (X \setminus A_i) \\
X \setminus \left( \bigcap_{i \in I} A_i \right) &= \bigcup_{i \in I} (X \setminus A_i)
\end{align*}
\begin{explanation-of-steps}
If each $A_i$ is closed, then each complement $(X \setminus A_i)$ is open. 
\begin{itemize}
    \item For an arbitrary intersection $\bigcap A_i$, its complement is the union $\bigcup (X \setminus A_i)$. Since the arbitrary union of open sets is open, the intersection is closed.
    \item For a finite union $\bigcup_{i=1}^N A_i$, its complement is the finite intersection $\bigcap_{i=1}^N (X \setminus A_i)$. Since the finite intersection of open sets is open, the union is closed.
\end{itemize}
\end{explanation-of-steps}
\end{math-stroke}

% \begin{ai-global-state-checkpoint-invisible-content}
% timestamp: 00:24:00
% topic: Properties of closed sets and the importance of finiteness.
% board_state: prop:closed-sets-properties, eq:de-morgan
% next_goal: Show an example where infinite intersection of open sets is not open.
% open_loops: none
% \end{ai-global-state-checkpoint-invisible-content}

\begin{spoken-clean}[00:24:04 - 00:26:39]
Let me give some example, very easy, that this property when we said finite intersections or finite unions is important that it is finite. That when you get infinite intersections of open sets, you can get a closed set that is not open. So, example: finite is important. So, if you take the simplest --- well, one of the spaces that you are most familiar with that is $\mathbb{R}$ with the Euclidean distance, and you take the sets $U_k = (-1/k, 1/k)$. And you take its intersection, then the intersection for $k=1$ to infinity of $U_k$, this gives you the point $\{0\}$. So, these are open intervals, and it's very easy because the intervals do not contain the endpoints, each of them is open. But when you do their intersection, you get just the point zero, and the point zero is not open because you cannot put any ball that contains anything. The point zero is closed but is not open. So, these are open and this is not open! That's why it's important that the intersection is finitely many. And passing to the complement you get a counterexample for the infinite unions of closed sets.
\end{spoken-clean}

\begin{math-stroke}[Example: Infinite Intersection of Open Sets]
\setcounter{theorem}{44}
\begin{example}[Infinite Intersection of Open Sets]\label{ex:infinite-intersection}
Consider the metric space $(\mathbb{R}, d_{\text{Eucl}})$. Define a sequence of open intervals:
\[ U_k = \left( -\frac{1}{k}, \frac{1}{k} \right) \quad \text{for } k \in \{1, 2, 3, \dots\} \]
Each $U_k$ is an open set. However, their infinite intersection is:
\[ \bigcap_{k=1}^{\infty} U_k = \{0\} \]
\begin{center}
\begin{tikzpicture}[scale=2]
    \draw[->] (-1.5,0) -- (1.5,0) node[right] {$\mathbb{R}$};
    \draw[ultra thick, MidnightBlue] (-1,0) -- (1,0);
    \node[below, MidnightBlue] at (-1,0) {$-1$};
    \node[below, MidnightBlue] at (1,0) {$1$};
    \node[above, MidnightBlue] at (0, 0.1) {$U_1$};

    \draw[ultra thick, BrickRed] (-0.5,0) -- (0.5,0);
    \node[below, BrickRed] at (-0.5,0) {\tiny $-1/2$};
    \node[below, BrickRed] at (0.5,0) {\tiny $1/2$};
    \node[above, BrickRed] at (0, 0.2) {$U_2$};

    \draw[ultra thick, ForestGreen] (-0.25,0) -- (0.25,0);
    \node[above, ForestGreen] at (0, 0.3) {$U_4$};
    
    \fill[black] (0,0) circle (1.5pt) node[below] {$0$};
\end{tikzpicture}
\end{center}
\begin{explanation-of-steps}
As $k \to \infty$, the intervals shrink around the origin. The only point contained in \emph{every} interval is $0$. The set $\{0\}$ is a singleton set, which is closed in $\mathbb{R}$ but not open, as any open ball $B(0, r) = (-r, r)$ contains points other than $0$.
\end{explanation-of-steps}
\end{example}
\end{math-stroke}

% ==========================================
% SECTION 5: INTERIOR AND CLOSURE
% ==========================================
\section{Interior and Closure}

\begin{spoken-clean}[00:26:39 - 00:28:30]
Okay, so now we will define something interesting that is what we were saying there heuristically: so what is the interior and the boundary of a set in a metric space. So, given a set $\Omega$ in my metric space $X$, I define the interior of $\Omega$. This is called the interior of $\Omega$. This is by definition the largest open set contained in $\Omega$. So, you have your set $\Omega$ and you look at all the open sets that are contained and you get the largest one is the union of all the open sets that contain $\Omega$. This is the interior.
\end{spoken-clean}

\begin{math-stroke}[Definition: Interior and Closure]
\setcounter{theorem}{45}
\begin{definition}[Interior and Closure]\label{def:interior-closure}
Let $(X, d)$ be a metric space and $\Omega \subseteq X$ be any subset.
\begin{itemize}
    \setcounter{enumi}{0} \item The \emph{interior} of $\Omega$, denoted $\operatorname{int}(\Omega)$ or $\Omega^\circ$, is the largest open set contained in $\Omega$. Formally, it is the union of all open subsets of $\Omega$:
    \[ \operatorname{int}(\Omega) = \bigcup \{ U \subseteq \Omega \mid U \text{ is open} \} \]
    \item The \emph{closure} of $\Omega$, denoted $\overline{\Omega}$, is the smallest closed set containing $\Omega$.
\end{itemize}
\end{definition}
\end{math-stroke}

\begin{spoken-clean}[00:28:30 - 00:29:59]
Now, I also define the closure, that you know from $\mathbb{R}$, the closure of some set is the set of points $x$ in $X$ such that there exists some sequence $x_n$ such that the whole sequence... \textit{[Audio cuts off abruptly]}
\end{spoken-clean}

\begin{math-stroke}[Sequential Characterization of Closure]
A point $x \in X$ belongs to the closure $\overline{\Omega}$ if and only if it can be reached as the limit of a sequence of points from $\Omega$:
\[ \overline{\Omega} = \{ x \in X \mid \exists (x_n)_{n \ge 0} \subseteq \Omega \text{ s.t. } x_n \to x \} \]
\textit{[Audio cuts off abruptly]}
\end{math-stroke}

% [SYSTEM] Video complete.
% ==========================================
% SECTION 5: INTERIOR AND CLOSURE (CONTINUED)
% ==========================================

\begin{spoken-clean}[00:00:00 - 00:01:32]
\inlinemetanote{The lecturer is writing the definitions of interior, closure, and boundary on the right chalkboard}
\end{spoken-clean}

\begin{math-stroke}[Definition: Interior, Closure, and Boundary]
\setcounter{theorem}{44}
\begin{definition}[Interior, Closure, and Boundary]\label{def:interior-closure-boundary}
Given a set $\Omega \subseteq X$ in a metric space $(X, d)$:
\begin{itemize}
    \setcounter{enumi}{0} \item The \emph{interior} of $\Omega$, denoted $\operatorname{int}(\Omega)$ or $\Omega^\circ$, is the largest open set contained in $\Omega$. Formally, it is the union of all open subsets of $\Omega$:
    \[ \operatorname{int}(\Omega) = \Omega^\circ = \bigcup \{ U \subseteq \Omega \mid U \text{ is open} \} \]
    \item The \emph{closure} of $\Omega$, denoted $\overline{\Omega}$, is the smallest closed set containing $\Omega$. Formally, it is the intersection of all closed sets containing $\Omega$:
    \[ \overline{\Omega} = \bigcap \{ A \supseteq \Omega \mid A \text{ is closed} \} \]
    Equivalently, it is the set of all points that can be reached as limits of sequences in $\Omega$:
    \[ \overline{\Omega} = \{ x \in X \mid \exists (x_n)_{n \ge 0} \subseteq \Omega \text{ s.t. } x_n \to x \} \]
    \item The \emph{boundary} of $\Omega$, denoted $\partial \Omega$, is the set difference between the closure and the interior:
    \[ \partial \Omega = \overline{\Omega} \setminus \Omega^\circ \]
\end{itemize}
\end{definition}
\end{math-stroke}

\begin{spoken-clean}[00:01:32 - 00:02:29]
Given a set $\Omega \subset X$, $(X, d)$ is a metric space. I define the interior of $\Omega$. This is called the interior of $\Omega$. This is by definition the largest open set contained in $\Omega$. So, you have your set $\Omega$ and you look at all the open sets that are contained and you get the largest one is the union of all the open sets that contain $\Omega$ (i.e., actually that are contained in $\Omega$). This is the interior.
\end{spoken-clean}

\begin{spoken-clean}[00:02:29 - 00:03:43]
Now, I also define the closure, that you know from $\mathbb{R}$, the closure of some set is the set of points $x$ in $X$ such that there exists some sequence $x_n$ such that the whole sequence is contained in $\Omega$ and then the sequence converges to $x$. You can also define this as the intersection of all the closed sets that are bigger than $\Omega$ (i.e., that contain $\Omega$). So, the closure is the smallest closed set that contains $\Omega$.
\end{spoken-clean}

\begin{spoken-clean}[00:03:43 - 00:05:33]
And then their difference, this minus this, this is called the boundary and is denoted with this symbol, the del symbol. This is called the boundary of $\Omega$. And this boundary of $\Omega$ is the \qt{skin}. Right? And if you do any example in the plane, you will see what this is doing. So, the closure is taking all this set with the skin, including the inside. And then you are removing the inside. So, when you draw the --- for this set over here... \inlinemetanote{Points at the square example he is about to draw}
\end{spoken-clean}

\begin{math-stroke}[Example: The Square \texorpdfstring{$[0, 1) \times [0, 1)$}{[0, 1) x [0, 1)}]
Consider the set $\Omega = [0, 1) \times [0, 1) \subset \mathbb{R}^2$.
\begin{center}
\begin{tikzpicture}[scale=1.5]
% \begin{ai-tikz-planner-invisible-content}
% 1. Background: Axes.
% 2. Midground: The square with mixed boundaries.
% 3. Foreground: Labels.
% \end{ai-tikz-planner-invisible-content}
    \draw[->] (-0.5,0) -- (2,0);
    \draw[->] (0,-0.5) -- (0,2);
    
    % The Set Omega
    \draw[thick, MidnightBlue] (0,1) -- (0,0) -- (1,0); % Solid boundaries
    \draw[dashed, thick, MidnightBlue] (1,0) -- (1,1) -- (0,1); % Dashed boundaries
    \fill[MidnightBlue!10, opacity=0.6] (0,0) rectangle (1,1);
    
    \node[MidnightBlue] at (0.5, 0.5) {$\Omega$};
    \node[below] at (0,0) {$0$};
    \node[below] at (1,0) {$1$};
    \node[left] at (0,1) {$1$};
\end{tikzpicture}
\end{center}

\begin{explanation-of-steps}
The set $\Omega$ includes its bottom and left edges (solid lines) but excludes its top and right edges (dashed lines).
\end{explanation-of-steps}
\end{math-stroke}

\begin{spoken-clean}[00:05:33 - 00:07:55]
So, how do you draw this example? It would be something like this, right? \inlinemetanote{Draws the square on the board} Is the inside of this square, but you are missing this part of the boundary and this part of the boundary. Now, can you see that these points satisfy this definition here? If I take a point here $x$ here, is on the boundary --- so is on the closure, sorry, is on the closure. Because I can take a sequence of points from the inside, from the red part (i.e., the shaded interior), that approaches $x$, right? But these points that are not in my set because I am excluding them, they are also on the closure. This point here $x$ over here is also on the closure. This belongs to the closure because I can take a sequence from here that approaches my point. Right?
\end{spoken-clean}

\begin{math-stroke}[Visualizing Interior, Closure, and Boundary]
For the set $\Omega = [0, 1) \times [0, 1)$:
\begin{center}
\begin{tikzpicture}[scale=1.2]
    % Interior
    \begin{scope}[shift={(0,0)}]
        \draw[dashed, thick, MidnightBlue] (0,0) rectangle (1,1);
        \fill[MidnightBlue!10] (0,0) rectangle (1,1);
        \node at (0.5, -0.3) {$\Omega^\circ = (0, 1)^2$};
    \end{scope}
    
    % Closure
    \begin{scope}[shift={(2.5,0)}]
        \draw[thick, MidnightBlue] (0,0) rectangle (1,1);
        \fill[MidnightBlue!10] (0,0) rectangle (1,1);
        \node at (0.5, -0.3) {$\overline{\Omega} = [0, 1]^2$};
    \end{scope}
    
    % Boundary
    \begin{scope}[shift={(5,0)}]
        \draw[thick, BrickRed] (0,0) rectangle (1,1);
        \node at (0.5, -0.3) {$\partial \Omega$};
    \end{scope}
\end{tikzpicture}
\end{center}
\begin{explanation-of-steps}
The interior $\Omega^\circ$ is the open square $(0, 1) \times (0, 1)$, which excludes all edges. The closure $\overline{\Omega}$ is the closed square $[0, 1] \times [0, 1]$, which includes all edges. The boundary $\partial \Omega$ is the set of all four edges of the square.
\end{explanation-of-steps}
\end{math-stroke}

\begin{spoken-clean}[00:07:55 - 00:08:52]
So, again all of these are perfectly common sense notions but that we can formalize and make perfectly rigorous just using these few axioms. And then, yes, I will prove maybe I will prove a lemma. \inlinemetanote{The lecturer erases the left chalkboard and prepares to write a new lemma}
\end{spoken-clean}

% \begin{ai-global-state-checkpoint-invisible-content}
% timestamp: 00:06:00
% topic: Definitions of Interior, Closure, and Boundary with the square example.
% board_state: def:interior-closure-boundary, ex:square-topology
% next_goal: State and prove the lemma characterizing open/closed sets via sequences.
% open_loops: none
% \end{ai-global-state-checkpoint-invisible-content}

% ==========================================
% SECTION 6: CHARACTERIZING OPEN AND CLOSED SETS VIA SEQUENCES
% ==========================================
\section{Characterizing Open and Closed Sets via Sequences}

\begin{spoken-clean}[00:08:52 - 00:10:12]
This lemma is called open and closed through sequences. How we determine if a set is open and closed using convergent sequences. So... \inlinemetanote{Starts writing the lemma on the board}
\end{spoken-clean}

\begin{math-stroke}[Lemma: Open and Closed Sets through Sequences]
\setcounter{theorem}{47}
\begin{lemma}[Open and Closed Sets through Sequences]\label{lem:open-closed-sequences}
Let $(X, d)$ be a metric space.
\begin{enumerate}
    \setcounter{enumi}{0} \item A subset $U \subseteq X$ is \emph{open} if and only if for every sequence $(x_n)_{n \ge 0}$ in $X$ that converges to a point $x \in U$, the sequence elements $x_n$ belong to $U$ \emph{eventually}.
    \[ U \text{ is open} \iff \left( \forall (x_n) \to x \in U \implies x_n \in U \text{ eventually} \right) \]
    \item A subset $A \subseteq X$ is \emph{closed} if and only if for every sequence $(x_n)_{n \ge 0}$ contained in $A$ that converges to a point $x \in X$, the limit $x$ also belongs to $A$.
    \[ A \text{ is closed} \iff \left( \forall (x_n)_{n \ge 0} \subseteq A \text{ s.t. } x_n \to x \implies x \in A \right) \]
\end{enumerate}
\end{lemma}
\end{math-stroke}

\begin{spoken-clean}[00:10:12 - 00:12:10]
So, this concept of \qt{up to finitely many exceptions} is very useful and is very, uh, nice to --- to say the thing shortly. And we give a name to this. And you have a more formal definition in the --- in the notes. So, when something happens --- when you have a sequence and some property, like for instance being in $U$, happens, is true up to finitely many exceptions, we say that this property happens \qt{eventually}.
\end{spoken-clean}

\begin{math-stroke}[Definition: Eventually]
\setcounter{theorem}{45}
\begin{definition}[Eventually]\label{def:eventually}
A property $P(x_n)$ is said to hold \emph{eventually} for a sequence $(x_n)_{n \ge 0}$ if it holds for all but finitely many terms. Formally:
\[ \exists N \in \mathbb{N} \text{ s.t. } \forall n \ge N, P(x_n) \text{ is true} \]
\end{definition}
\end{math-stroke}

\begin{proof}[Proof of Lemma \ref{lem:open-closed-sequences} (Part 1)]
\begin{spoken-clean}[00:12:10 - 00:15:50]
Proof of (1). So, I need to do two implications. First I will do this side ($\implies$) that is the one I draw here \inlinemetanote{points at the drawing of the open set $U$}. Assume that $U$ is open. And now take a point --- I'm just going to put in symbols the drawing that I did here. Take $x$ in $U$, any point. Then there --- by definition of open set, since $U$ is open, exists some $r$ positive such that the ball of radius $r$ centered at $x$ is contained in $U$. 

And now, since by assumption $x_n$ is converging to $x$, by definition of convergence, what was the definition of convergence? Remember: given any $\epsilon$ --- I called it in the definition --- this was the $\epsilon$ (i.e., the radius $r$), right? Given any $\epsilon$, now is $r$ but is the same. For $n$ large enough, bigger than a certain capital $N$, my sequence will belong to the ball of radius $\epsilon$ (i.e., $r$).
\end{spoken-clean}

\begin{math-stroke}[Forward Implication \texorpdfstring{$\implies$}{=>}]
Assume $U$ is open and $x_n \to x \in U$.
Since $U$ is open:
\[ \exists r > 0 \text{ s.t. } B(x, r) \subseteq U \]
Since $x_n \to x$, by the definition of convergence (using $\epsilon = r$):
\[ \exists N \in \mathbb{N} \text{ s.t. } \forall m \ge N, d(x_m, x) < r \]
This is equivalent to saying:
\[ \forall m \ge N, x_m \in B(x, r) \]
Since $B(x, r) \subseteq U$, it follows that:
\[ \forall m \ge N, x_m \in U \]
Thus, $x_n \in U$ eventually.
\end{math-stroke}

\begin{spoken-clean}[00:15:50 - 00:18:16]
Now I will prove the second implication ($\impliedby$). And I will prove this by contraposition. $A \implies B$ is equivalent to $\neg B \implies \neg A$, right? We know that. So, I will prove that \qt{not left} (i.e., $U$ is not open) implies \qt{not right} (i.e., there exists a sequence $x_n \to x \in U$ such that $x_n \notin U$ infinitely often).

So, now, let's do a bit of logic. What is the negation of $U$ open? Remember now, every time: what is the definition of $U$ open? $U$ open is: for every point in my set $U$, exists some radius such that the ball is contained in $U$. Now we do the negation. The negation of \qt{for every point} is \qt{exists a point}. So, $U$ is not open means exists $x$ in $U$ such that for every radius $r$ positive, the ball centered at $x$ with radius $r$ is not contained in $U$.
\end{spoken-clean}

\begin{math-stroke}[Backward Implication \texorpdfstring{$\impliedby$}{<=} (via Contraposition)]
We prove that if $U$ is not open, then there exists a sequence $x_n \to x \in U$ such that $x_n \notin U$ for all $n$.
        
\emph{Negation of Openness:}
\[ U \text{ is not open} \iff \exists x \in U \text{ s.t. } \forall r > 0, B(x, r) \not\subseteq U \]
The condition $B(x, r) \not\subseteq U$ means that the ball must contain at least one point from the complement $X \setminus U$:
\[ \forall r > 0, B(x, r) \cap (X \setminus U) \neq \emptyset \]
\end{math-stroke}

\begin{spoken-clean}[00:18:16 - 00:19:15]
\inlinemetanote{A student asks a question about notation} Ah, you mean this? \inlinemetanote{points at the $\subseteq$ symbol on the board} Is the same. You mean this? \inlinemetanote{points at the $\subset$ symbol} Is the same meaning. So, some people use the equality --- since it means the same, I save one line, right? But many --- this is very standard, many people just use this. So, I will use always this \qt{contained} without the equality sign.
\end{spoken-clean}

\begin{student-interaction}[Student Question]
You write the subset symbol without the extra line at the bottom. Does it have a meaning, like that it cannot be equal, or is that just notation?
\end{student-interaction}

\begin{spoken-clean}[continued]
It's just notation. In this course, $\subset$ and $\subseteq$ are used interchangeably to mean subset.
\end{spoken-clean}

\begin{spoken-clean}[00:19:15 - 00:21:15]
Okay, so back to the proof. Taking $r = 1/n$ for $n \ge 1$, I produced a sequence of points --- so these balls are smaller and smaller --- and these points $x_n$ (i.e., the $x_r$ for $r = 1/n$) are in balls closer and closer to $x$. So, they are converging to the point $x$. Right? So, I found a sequence that is converging to the point $x$ and that comes from the outside. But then, this is false (i.e., it contradicts the \qt{right} side of our implication). Because I'm saying that up to finitely many exceptions, the sequence should be contained in $U$. And all of the elements of this sequence are outside of $U$. So, this is false. This contradicts that $x_n$ belongs to $U$ eventually. And it ends the proof of this implication.
\end{spoken-clean}

\begin{math-stroke}[Constructing the Counter-Sequence]
For each $n \in \{1, 2, 3, \dots\}$, choose a radius $r_n = \frac{1}{n}$. By our assumption, there exists a point:
\[ x_n \in B(x, 1/n) \cap (X \setminus U) \]
This construction yields a sequence $(x_n)_{n \ge 1}$ such that:
\begin{enumerate}
    \setcounter{enumi}{0} \item $x_n \notin U$ for all $n$.
    \item $d(x_n, x) < \frac{1}{n} \to 0$, so $x_n \to x$.
\end{enumerate}
Since $x \in U$, the \qt{right} side of our lemma would require $x_n \in U$ eventually. However, our sequence is entirely outside $U$, which is a contradiction. Thus, $U$ must be open.
\end{math-stroke}
\end{proof}

\begin{spoken-clean}[00:21:15 - 00:24:04]
So, this proof by contraposition is very useful. Now, part (2) of the lemma is very similar. I leave it as an exercise. Try to do it, again. If you understood well the definitions and you can draw a picture of what's going on, this shouldn't be too difficult. But it can be a bit difficult at the beginning because there are a lot of new definitions. The best is to practice with these exercises. 

Let me put another exercise before continuity. Prove that if $(X, d)$ is a complete metric space and $A$ is a closed subset, then $A$ with the same distance restricted is complete. Example: we know that $\mathbb{R}^n$ is a complete metric space. Then any closed subset of $\mathbb{R}^n$ will be a complete metric space when you restrict the distance. And now you have all the ingredients to prove that. It's not too difficult. It's a good exercise.
\end{spoken-clean}

\begin{math-stroke}[Exercise: Completeness of Closed Subsets]
\begin{nice-box}[Exercise]
Prove the following proposition:
\begin{proposition}\label{prop:closed-subset-complete}
Let $(X, d)$ be a complete metric space. If $A \subseteq X$ is a closed subset, then the metric subspace $(A, d|_A)$ is also a complete metric space.
\end{proposition}
\end{nice-box}
\end{math-stroke}

\begin{spoken-clean}[00:24:04 - 00:25:30]
Okay, so I see the usual problem: when I see people that are very serious, I don't know if they are completely lost, if they find it super boring, or what is the problem. So, if someone has questions and wants to participate, we can slow down and ask questions. Otherwise, we continue with other stuff. We still need to do continuity, right? Well, let's make the pause. \inlinemetanote{The lecturer prepares for the break}
\end{spoken-clean}

\begin{lecture-break}[15-Minute Break]
The lecturer announces a break. The video resumes with the lecturer addressing the class after the break.
\end{lecture-break}

\begin{spoken-clean}[00:25:30 - 00:26:39]
Okay, I will start if you don't listen... well. So, may you forgive me for the stereotype, but when I studied mathematics in Barcelona, which is a Mediterranean country, people started on time the lecture. So, they were silent. We are in Switzerland, you should be like the champions of starting on time. And I'm seeing the contrary. I'm very surprised. Okay. And you know this is recorded, so this will be the proof of the misbehavior will be there for the history.
\end{spoken-clean}

% [SYSTEM] Video complete.
\begin{spoken-clean}[00:00:00 - 00:00:08]
...standard, many people just use this \inlinemetanote{points at the $\subseteq$ symbol on the board}. So, I will use always this \qt{contained} without the equality sign. \inlinemetanote{The lecturer begins erasing the chalkboards to prepare for the next section}
\end{spoken-clean}

\begin{meta-note}[Topic Transition]
The lecturer erases the center and right chalkboards. He briefly writes down the exercise mentioned at the end of the previous segment before transitioning to the main topic of the second half of the lecture: Continuity.
\end{meta-note}

\begin{math-stroke}[Exercise: Completeness of Closed Subsets]
\begin{nice-box}[Exercise]
\setcounter{proposition}{49}
\begin{proposition}\label{prop:closed-subset-complete-v2}
Let $(X, d)$ be a complete metric space. If $A \subseteq X$ is a closed subset, then the metric subspace $(A, d|_A)$ is also a complete metric space.
\end{proposition}
\end{nice-box}
\end{math-stroke}

\begin{spoken-clean}[00:00:24 - 00:01:25]
Okay, let's do this because, um, maybe you will find more exciting continuous functions. Since this proof over here \inlinemetanote{points at the board} is super similar to the previous one, and is again --- you see that these proofs are one-line proofs. It's just using the definitions and manipulating around. So, I give this proof --- I left this proof as an exercise. Try to do it, again. If you understood well the definitions and you can draw a picture of what's going on, this shouldn't be too difficult. 

But it can be a bit difficult at the beginning because there are a lot of new definitions. You need to learn them. But the best is to practice with these exercises. To try to do them. Because if I do them, um, it looks like a bunch of symbols that all work sometimes, but --- but you don't interiorize the definitions. The best --- the best way is to try by yourself.
\end{spoken-clean}

\begin{spoken-clean}[00:01:25 - 00:03:41]
Okay, so I will --- I will do something more exciting now, that is continuity. Because you already know... \inlinemetanote{The lecturer erases the board and writes the heading \qt{Continuity}} Okay. Now we start to have the problems with the shadows. Maybe... so, someone, I don't remember who, complained about shadows. Now I can write here, right? Is just when the other blackboard is annoying. Okay. So, now we will see continuity.
\end{spoken-clean}

% ==========================================
% SECTION 8: CONTINUITY IN METRIC SPACES
% ==========================================
\section{Continuity in Metric Spaces}

\begin{spoken-clean}[00:03:41 - 00:05:47]
For continuity, we will define continuous maps. So, we will be in this setup. $f$ will be a map between two sets $X$ and $Y$. And they are both metric spaces. So, we will have some $X$ with its distance $d_X$, and $Y$ with its distance $d_Y$. So, many times for us, these sets will be the Euclidean space of some dimension and the Euclidean space with some other dimension, right? But it can be whatever. 

So, this is the setup. Both --- both are different metric spaces. Well, different or the same; they could be the same. Now, definition. We say that a map like this is $f$ is continuous if one of the following three equivalent properties hold true.
\end{spoken-clean}

\begin{math-stroke}[Definition: Continuity]
\setcounter{section}{2} \setcounter{subsection}{1} \setcounter{definition}{52}
\begin{definition}[Continuity]\label{def:continuity}
Let $(X, d_X)$ and $(Y, d_Y)$ be metric spaces. A mapping $f: X \to Y$ is said to be \emph{continuous} if it satisfies any of the following three equivalent conditions:
\end{definition}
\end{math-stroke}

\begin{spoken-clean}[00:05:47 - 00:07:10]
So, now I will state the three definitions of continuity, because at least with two of them you are very familiar from --- from Analysis 1, right? But I will state now postulating that they are equivalent, and later we will prove immediately that they are indeed equivalent. Right? So, the definition number one of continuity is the $\epsilon$-$\delta$ continuity that you are familiar. 

So, you know this very well for functions of one real variable. Now we need to abstract the same notion in metric spaces. So, how was the definition of --- of $\epsilon$-$\delta$? So, can someone remind me in the real line what is the definition of continuity with $\epsilon$-$\delta$? Does someone remember? \inlinemetanote{The lecturer retrieves the catch-box microphone} Yeah, good. \inlinemetanote{Tosses microphone to student}
\end{spoken-clean}

\begin{student-interaction}[Student Answer]
I think it was: a function is continuous at $x_0$ if for every $\epsilon > 0$, there exists a $\delta > 0$ such that for all $x$, if $|x - x_0| < \delta$, then $|f(x) - f(x_0)| < \epsilon$.
\end{student-interaction}

\begin{spoken-clean}[continued]
So, up to --- is very good. You said what is continuous at a point. For a continuous function is continuous at all the points, you were going to say that. But I will use directly the continuous at --- at every point. So, I will start exactly as you said. He said: well, first I will add --- since I want to be continuous at every point, I will add \qt{for every point}. For all $x$ in $X$. Okay. Next is what you said. For every $x$, given --- for all $\epsilon$ positive, exists $\delta$ positive. And now the rest of your definition used intervals and stuff that are not available in metric spaces. So, how we translate?
\end{spoken-clean}

\begin{math-stroke}[The Three Faces of Continuity]
\setcounter{proposition}{53}
\begin{proposition}[The Three Faces of Continuity]\label{prop:three-faces-continuity}
Let $f: (X, d_X) \to (Y, d_Y)$ be a map between metric spaces. The following are equivalent:
\begin{enumerate}
    \setcounter{enumi}{0} \item \emph{\texorpdfstring{$\epsilon$}{epsilon}-\texorpdfstring{$\delta$}{delta} Continuity:}
    \[ \forall x \in X, \forall \epsilon > 0, \exists \delta > 0 \text{ s.t. } \forall x' \in X, d_X(x, x') < \delta \implies d_Y(f(x), f(x')) < \epsilon \]
\end{enumerate}
\end{math-stroke}

\begin{spoken-clean}[00:08:23 - 00:10:23]
So, the notion will be that if the distance between $x$ and $x'$ is less than $\delta$, then the distance between $f(x)$ and $f(x')$ is less than $\epsilon$. And what is this distance? This distance is the distance in the space here, $X$, and this distance --- these guys belong to $Y$, so this is the distance in $Y$. Right? What else could be? Cannot be anything else. So, if you think about this, is the only way to generalize. 

Another way to say this, because this is a bit annoying to write, is a bit long, is to say that the ball --- let me put it like this: this thing, the bracket red, is the same, exactly the same as saying that the $f$ --- when I compute $f$ the image of the ball of radius $\delta$ centered at $x$ will be contained in the ball of radius $\epsilon$ centered at $f(x)$.
\end{spoken-clean}

\begin{math-stroke}[Geometric Interpretation of \texorpdfstring{$\epsilon$}{epsilon}-\texorpdfstring{$\delta$}{delta}]
The $\epsilon$-$\delta$ condition can be expressed concisely using open balls:
\begin{equation} \label{eq:continuity_balls}
f(B_\delta(x)) \subseteq B_\epsilon(f(x))
\end{equation}
\begin{explanation-of-steps}
This states that the image of a sufficiently small neighborhood around $x$ (the $\delta$-ball) must be entirely contained within a given neighborhood around the image $f(x)$ (the $\epsilon$-ball).
\end{explanation-of-steps}
\end{math-stroke}

\begin{spoken-clean}[00:10:23 - 00:12:18]
So, this is the $\epsilon$-$\delta$ definition. Second definition is with sequences. Right? Continuity with sequences. And --- and this is the same as in $\mathbb{R}$ essentially. So, if you have $x_n$ is a sequence and $x_n$ is converging to a point $x$, then the sequence of the images, that is $f(x_n)$, is convergent in the space $Y$ with limit $f(x)$. So, and a compact way to write this is that $f(x_n)$ this sequence in $n$ converges when $n$ goes to infinity to $f(x)$. And for this essentially you don't even need anything new. So, is really the same definition of $\mathbb{R}$, you plug it here and it makes sense.
\end{spoken-clean}

\begin{math-stroke}[Sequential Continuity]
\begin{enumerate}
    \setcounter{enumi}{1} \item \emph{Sequential Continuity:}
    \[ \forall (x_n)_{n \ge 0} \subset X, \quad x_n \to x \implies f(x_n) \to f(x) \]
\end{enumerate}
\end{math-stroke}

\begin{spoken-clean}[00:12:18 - 00:14:15]
And then there is the new notion of continuity, that is the topological notion of continuity, the number three, that this one is new but will be very easy to prove that is the same. And we will use this all the time, that's why I'm giving. That says that for every $V$ in $Y$ open, the pre-image $f^{-1}(V)$ that is a subset of $X$ is open. 

A map, a function is continuous is whenever you take an open set in the target and you see it back --- so you go $f^{-1}$ --- you get an open set here in the --- in the domain. Now, maybe is not obvious that the three are equivalent, but we will show that.
\end{spoken-clean}

\begin{math-stroke}[Topological Continuity]
\begin{enumerate}
    \setcounter{enumi}{2} \item \emph{Topological Continuity:}
    \[ \forall V \subseteq Y \text{ open} \implies f^{-1}(V) \subseteq X \text{ is open} \]
\end{enumerate}
\end{math-stroke}

% ==========================================
% SECTION 9: PROVING THE EQUIVALENCES
% ==========================================
\section{Proving the Equivalences}

\begin{spoken-clean}[00:14:15 - 00:16:35]
Okay, let's give the proof of this proposition. So, first I need to start to show --- so I will start with the $\epsilon$-$\delta$ continuity, right? And I will show that (1) implies (2). So, first step I will show that (1) implies (2). By assumption, I am in (1). So, my function $f$ assume that $f$ from $X$ to $Y$ is $\epsilon$-$\delta$ continuous. And let $x_n$ be a sequence converging to a point $x$. And we want to show that $f(x_n)$ converges to $f(x)$. Right? 

So, I need to use that is $\epsilon$-$\delta$ continuous. So, the first thing is to find the $\epsilon$. So, you give me the $\epsilon$ positive. What I want to prove --- let me put it like this: I want to prove that the distance between $f(x_n)$ and $f(x)$ is less than $\epsilon$ eventually for every $\epsilon$ positive. You see, I have rephrased here I have rephrased the notion of convergence. So, what I want to prove really is that $f(x_n)$ converges to $f(x)$. This is my goal.
\end{spoken-clean}

\begin{proof}[Proof of Proposition \ref{prop:three-faces-continuity}]
\begin{math-stroke}[Forward Implication \texorpdfstring{$(1) \implies (2)$}{(1) => (2)}]
\begin{ai-proof-skeleton-invisible-content}
% 1. Assume f is epsilon-delta continuous.
% 2. Let x_n -> x.
% 3. Goal: Show f(x_n) -> f(x).
% 4. For any epsilon > 0, find delta > 0 from (1).
% 5. Since x_n -> x, x_n is in B_delta(x) eventually.
% 6. By (1), f(x_n) is in B_epsilon(f(x)) eventually.
\end{ai-proof-skeleton-invisible-content}
Assume $f$ is $\epsilon$-$\delta$ continuous. Let $(x_n)_{n \ge 0}$ be a sequence in $X$ such that $x_n \to x$.
We want to show that $f(x_n) \to f(x)$, which means:
\[ \forall \epsilon > 0, \quad d_Y(f(x_n), f(x)) < \epsilon \text{ eventually} \]
Given $\epsilon > 0$, by the $\epsilon$-$\delta$ continuity of $f$:
\[ \exists \delta > 0 \text{ s.t. } f(B_\delta(x)) \subseteq B_\epsilon(f(x)) \]
Since $x_n \to x$, by the definition of convergence:
\[ x_n \in B_\delta(x) \text{ eventually} \]
Applying the function $f$ to these terms:
\[ f(x_n) \in f(B_\delta(x)) \subseteq B_\epsilon(f(x)) \text{ eventually} \]
This implies $d_Y(f(x_n), f(x)) < \epsilon$ eventually, which proves that $f(x_n) \to f(x)$.
\end{math-stroke}

\begin{spoken-clean}[00:16:35 - 00:18:12]
But I can rephrase this goal like this: I want to prove that the distance between --- sorry, last ten minutes start to... so this is with $f$. Now is correct. So, what I want to prove is that $f(x_n)$ converges to $f(x)$. This is my goal. But I am putting it in this way because here at least I see the $\epsilon$. So, given the $\epsilon$ I can find the $\delta$. So, now, given $\epsilon$ positive, by $\epsilon$-$\delta$ continuity exists some $\delta$ positive such that the ball of radius $\delta$ centered at $x$ when I apply $f$ to this ball will be contained in the ball of radius $\epsilon$ centered at $f(x)$.
\end{spoken-clean}

\begin{spoken-clean}[00:18:12 - 00:19:50]
But remember my assumption: my assumption is that $x_n$ was converging to $x$. So, since $\delta$ is positive, $x_n$ belongs to this set (i.e., $B_\delta(x)$) eventually. When I say eventually is like for $n$ very large, $x_n$ will belong to this set. But then, when I read here, the image then this implies that $f(x_n)$ belongs to the ball of radius $\epsilon$ centered at $f(x)$ eventually. So, for $n$ large. And this is what I wanted to show. Right? So, this proof shows that (1) implies (2).
\end{spoken-clean}

\begin{spoken-clean}[00:19:50 - 00:21:15]
So, since this is maybe very symbolic, let me do a picture of what is going on. So, the pictures are never rigorous but they help a lot to --- to do some intuition. And even if I don't finish the proof today of all the equivalences, is good to go step by step. So, some people work very good with these proofs that are symbols. For some people like myself, for me are very difficult to think directly with these symbols. So, when I write this kind of proof before writing to the notes or something, what I do is I do the drawing. First I do the drawing and then I translate the drawing in the symbols. Directly the symbols for some people is easy, for some people is more difficult. So, it depends on which way you reason.
\end{spoken-clean}

\begin{math-stroke}[Visualizing Continuity]
\begin{center}
\begin{tikzpicture}[scale=1.2]
% \begin{ai-tikz-planner-invisible-content}
% 1. Background: Two metric spaces X and Y.
% 2. Midground: Point x in X and f(x) in Y.
% 3. Foreground: Sequence x_n approaching x, and f(x_n) approaching f(x).
% 4. Annotations: delta-ball in X and epsilon-ball in Y.
% \end{ai-tikz-planner-invisible-content}
    % Space X
    \draw[thick] (0,0) ellipse (1.5cm and 1cm);
    \node at (-1, 0.7) {$X$};
    \coordinate (X) at (0.5, -0.2);
    \fill (X) circle (1.5pt) node[below] {$x$};
    
    % delta-ball
    \draw[dashed, MidnightBlue] (X) circle (0.4cm);
    \node[MidnightBlue, below] at (0.5, -0.6) {\footnotesize $B_\delta(x)$};

    % Space Y
    \begin{scope}[shift={(5,0)}]
        \draw[thick] (0,0) ellipse (1.5cm and 1cm);
        \node at (1, 0.7) {$Y$};
        \coordinate (FX) at (-0.5, -0.2);
        \fill (FX) circle (1.5pt) node[below] {$f(x)$};
        
        % epsilon-ball
        \draw[dashed, BrickRed] (FX) circle (0.6cm);
        \node[BrickRed, below] at (-0.5, -0.8) {\footnotesize $B_\epsilon(f(x))$};
    \end{scope}

    % Mapping
    \draw[->, thick] (1.8, 0.2) to[bend left=20] node[midway, above] {$f$} (3.2, 0.2);

    % Sequence
    \foreach \i in {1, 2, 3, 4}
        \fill (0.5 - 0.8/\i, -0.2 + 0.5/\i) circle (1pt);
    \node at (-0.3, 0.4) {\footnotesize $x_n$};
    
    \foreach \i in {1, 2, 3, 4}
        \fill (4.5 - 0.4/\i, -0.2 + 0.3/\i) circle (1pt);
    \node at (4.1, 0.2) {\footnotesize $f(x_n)$};
\end{tikzpicture}
\end{center}
\end{math-stroke}

\begin{spoken-clean}[00:22:32 - 00:24:26]
So, what is going on in this proof? We have our first metric space $X$ and our second metric space $Y$. And we have a continuous map from $X$ to $Y$. And then what we wanted to show is that whenever we have a sequence here approaching $x$, the corresponding sequence of images will approach the $f(x)$. And we had to reason with the $\epsilon$-$\delta$. 

So, what this that this sequence converges here means literally that whenever I pick any ball --- can be very, very tiny this ball, very small --- then if I wait enough, all the elements of the sequence will fall inside of my ball. Right? If the ball is smaller, I need to wait more and more. But it always falls. And now, this is what I want to prove. So, to find this ball and to find that if I wait long enough I will be here, I need to use that the $x_n$ is convergent to $x$. So, I need to go back with this ball and see what happens with this ball, right? So, what happens is that this ball here comes from some set here around $x$ that may not be a ball. But what I know is that here there is a ball of radius $\delta$. This is the $\epsilon$-$\delta$. The blue ball. That when I compute the image of this ball, maybe it becomes a \qt{potato} or something, but is inside of my original ball. Right? And since this is a ball, the yellow points eventually will be in this ball, and then when I take the image, the points will be in my ball because the blue is inside of the bigger ball. So, with drawings is a bit clearer for some people perhaps.
\end{spoken-clean}

\begin{spoken-clean}[00:24:26 - 00:24:53]
\inlinemetanote{A student asks a question} Ah, there is some... \inlinemetanote{The lecturer retrieves the catch-box microphone}
\end{spoken-clean}

\begin{student-interaction}[Student Question]
Can I give back the cube eventually?
\end{student-interaction}

\begin{spoken-clean}[continued]
\inlinemetanote{The lecturer laughs} Ah, you need to get --- this was the question, to get back the cube? No. So, or someone else has a question. Who has a question? No, I had the cube the whole time and I was waiting for you to... ah, okay, okay.
\end{spoken-clean}

\begin{spoken-clean}[00:24:53 - 00:26:13]
Okay, now let's start to prove that (2) implies (3). To prove that (2) implies (3), we will actually prove that \qt{not 3} implies \qt{not 2}. So, again contraposition. Remember: what was 3? 3 was saying that for every $V$ subset $Y$ open, $f^{-1}(V)$ subset $X$ is open. Now, what is \qt{not 3}? The negation of this. The negation of \qt{for all} is \qt{exists}. So, \qt{not 3} means exists some $V$ subset $Y$ open such that $f^{-1}(V)$ is not open. Right?
\end{spoken-clean}

\begin{math-stroke}[Forward Implication \texorpdfstring{$(2) \implies (3)$}{(2) => (3)} (via Contraposition)]
\begin{ai-proof-skeleton-invisible-content}
% 1. Assume (3) is false.
% 2. Goal: Show (2) is false.
% 3. If (3) is false, there exists an open set V in Y such that f^{-1}(V) is not open.
% 4. Since f^{-1}(V) is not open, there exists x in f^{-1}(V) and a sequence x_n -> x such that x_n is not in f^{-1}(V) for all n.
% 5. Then f(x_n) is not in V for all n.
% 6. But f(x) is in V. Since V is open, f(x_n) cannot converge to f(x).
% 7. This contradicts (2).
\end{ai-proof-skeleton-invisible-content}
We prove $\neg (3) \implies \neg (2)$.
Assume (3) is false. Then there exists an open set $V \subseteq Y$ such that its pre-image $U = f^{-1}(V)$ is \emph{not open} in $X$.
\end{math-stroke}

\begin{spoken-clean}[00:26:13 - 00:28:12]
So, what does it mean not open? Not open, we said before. Not open, it means that exists some point $x$ in $f^{-1}(V)$ such that and there is a sequence $x_n$ convergent to $x$ --- this is what I did before, remember when I said a set is not open, there are a sequence of balls that have some point in the complement, right? So, there is a sequence $x_n$ convergent to $x$ and all the points $x_n$ do not belong or belong to the complement. Right? So, this gives me a sequence. And now, I will try to use this sequence to contradict (2).
\end{spoken-clean}

\begin{math-stroke}[Constructing the Counter-Sequence]
Since $U = f^{-1}(V)$ is not open, by Lemma \ref{lem:open-closed-sequences}:
\[ \exists x \in U \text{ and a sequence } (x_n)_{n \ge 0} \subset X \setminus U \text{ s.t. } x_n \to x \]
By the definition of the pre-image:
\begin{itemize}
    \item $x \in f^{-1}(V) \implies f(x) \in V$.
    \item $x_n \notin f^{-1}(V) \implies f(x_n) \notin V$ for all $n$.
\end{itemize}
\end{math-stroke}

\begin{spoken-clean}[00:28:12 - 00:29:34]
So, this sequence is converging to $x$, so I need to try to show that the images do not converge to $f(x)$. But then $f(x_n)$ belongs all belong to the complement of the image (i.e., actually the complement of $V$). And you notice that $V$ because $x$ belongs to $f^{-1}(V)$, $f(x)$ belongs to $V$. $f(x)$ belongs to an open set. And what I said before about open sets? That a set is open if and only if whenever you have a sequence eventually belongs to your set. But here all the sequence is in the complement of the set. So, $f(x_n)$ cannot converge to $f(x)$. Right? I use the characterization of open set by sequences. This point is inside my open set and the points are outside. So, these points can never converge to this one. Right? So, this gives that the negation of (2).
\end{spoken-clean}

\begin{math-stroke}[Contradicting Sequential Convergence]
We have $f(x) \in V$ and $f(x_n) \in Y \setminus V$ for all $n$.
Since $V$ is open, by Lemma \ref{lem:open-closed-sequences}, any sequence converging to $f(x)$ must eventually be contained in $V$.
However, our sequence $(f(x_n))$ is entirely contained in the complement $Y \setminus V$.
Therefore, $f(x_n) \not\to f(x)$.
This contradicts condition (2), completing the proof that $(2) \implies (3)$.
\end{math-stroke}
\end{proof}

\begin{spoken-clean}[00:29:34 - 00:29:59]
And since it is almost time, actually this is the easiest of the three, but the next day we will show that (3) implies (1) and we will close the --- we will close the circle. Okay? And since it is maybe the easiest, you can try to do as an exercise. If you succeed, less work for me. \inlinemetanote{The students applaud as the lecturer packs up his equipment}
\end{spoken-clean}

% [SYSTEM] Video complete.

%% ==========================================
% Session Info:
% System Prompt loaded: True
% History loaded: True
% 
% AI Studio User Prompt:
% [Dateien] 
% ==========================================

```latex
\section{Course Introduction and Logistics}

\begin{spoken-clean}[00:00:00 - 00:01:07]
Good morning. Good morning, thank you. Please, please, please. Uh, good morning and welcome to Analysis 2. Uh, I have here a projection with my name, Joaquim Serra, the name of the coordinator, Enric Florit-Simon. Now, this is my introduction. Um, and then you know the times of the lectures obviously because you are here. Um, and I wanted to share before starting with the mathematical content some practical informations that are very important because we are a large group and the communication and everything needs to be done in a more or less structured way, otherwise --- otherwise we collapse.

So, as I guess you did so as a general rule, uh, everything that applied to Analysis 1 applies to Analysis 2, unless otherwise stated. We will try to follow the same rules and then you know what to expect. In particular you will have this, uh, meta-for page \textit{[points at the projector screen]} that I can see that from from up there is a bit difficult to see, right? But I will try to put it bigger here.
\end{spoken-clean}

\begin{meta-note}[Projected Content: Course Website]
The lecturer displays the course website for \qt{Analysis II: mehrere Variablen Spring 2026}. It lists the lecturer (Joaquim Serra), coordinator (Enric Florit-Simon), lecture times (Mo 08:15-10:00 in ETA F 5, We 08:15-10:00 in HG F 7, Th 16:15-18:00 in ETA F 5), and information regarding exercise classes, communication via the Course Forum, recordings, and the final exam in August.
\end{meta-note}

\begin{spoken-clean}[00:01:07 - 00:02:55]
Um, so you know that you need to register for the exercise classes. Uh, for the communication, please try to use whenever possible the forum. We will look at the forum and and reply, and also your your peers can also reply to questions because many --- so then we can use the collective knowledge, right? Um, there will be recordings of the lecture, but not until the later on, okay? Because uh is important for the lecture and for you, but mainly --- I mean for you is your interest uh if you want to come or not, but for the lecture is very good that you come to the lecture, you ask questions, you stop me when I'm going too fast, these kind of things. You know, it's very important. So it's a --- and also is where you meet with your peers and you can communicate, ask questions to your peers. It's it's it's better for the learning.

Uh, we will tell you more about the exam later on, right? Today I will tell you a bit more what is important to be successful in the course, what to expect about this course, and uh the details and the rules of the exam will come later on. You don't need to know now. The bonus, there is a zero --- uh twenty-five bonus will be in questions that you will answer in some exercise classes. We will tell you in advance which ones, and uh essentially you need to be in the in the exercise class and answer it. And do not worry because uh the rule will be thought, in any case it's just $0.25$ points, uh but the rules will be thought so that if you miss one day or casuistics like this, it's not that you get zero points, right? Okay.
\end{spoken-clean}

\begin{spoken-clean}[00:02:55 - 00:04:15]
Uh, lecture notes. There will be lecture notes similar to those in Analysis 1, and we will put them in this page, hopefully by the end of each week, because I I update them after the lecture, mainly depending on your questions, depending on your feedback. And that's it. Uh, I think that the rules are mainly those and the organization. Is there some question uh about organization? So essentially as I said, um if you have doubts, uh is like in Analysis 1, unless I say something that is surprising to you as in Analysis 1. Is it okay, this? Yeah. Okay now, I will switch --- I have some other things I want to project later, but now because we cannot record at the same time blackboard and projection, I will switch one moment to the blackboard, so I can mute I think, yeah. Okay. Now we need the light. \textit{[The lecturer turns off the projector and prepares the blackboard]} Okay, I think that there is already on. Okay.
\end{spoken-clean}

\section{Analysis I vs. Analysis II: A Conceptual Overview}

\begin{spoken-clean}[00:04:15 - 00:05:30]
Okay now, just to start, um we will give a very quick --- so quick that it's almost a cartoon --- but a very quick overview of what is Analysis 1 and Analysis 2, right? It's very very quick. Analysis 1... is... sorry but for me it's very difficult to write big and then it's so --- so if you don't see something, I will try to say it out loud everything I write, but if if something --- I mean Analysis is a... so uh Analysis 1, calculus of one variable, right? 

And then if I if if we got to choose a representative theorem of Analysis 1, that a bit represents the course, uh what would you choose? So can you tell me, I know it's the first day, but let's start a bit of communication. Can you tell me at least two answers --- so two people can tell me what is the representative --- one representative theorem of Analysis 1? Yeah? \textit{[Student answers]} Yeah, I agree. Some other representative theorem --- he said the Fundamental Theorem of Calculus. Makes sense, right? Okay after --- is there some other guess? Is --- yeah, there? \textit{[Student answers]} The Intermediate Value Theorem. Okay.
\end{spoken-clean}

\begin{math-stroke}[The Fundamental Theorem of Calculus in 1D]
In Analysis I, the central object of study is the calculus of a single real variable.
\[
\text{Analysis 1: Calculus of One Variable}
\]
The most representative result is the \emph{Fundamental Theorem of Calculus} (FTC):
\begin{equation} \label{eq:ftc_1d}
\int_a^b f'(x) \, dx = f(b) - f(a)
\end{equation}

\begin{explanation-of-steps}
The lecturer initially writes $f(a) - f(b)$ on the board, but a student quickly spots the error. He corrects it to the standard $f(b) - f(a)$. He emphasizes that while the formula looks compact, understanding it required defining the derivative, the Riemann integral, and the structure of real numbers $\mathbb{R}$.
\end{explanation-of-steps}
\end{math-stroke}

\section{The Fundamental Theorem of Calculus in Higher Dimensions}

\begin{spoken-clean}[00:05:30 - 00:06:23]
I think we need to use this for the recording otherwise I will repeat. Later on when I get better with the throwing the \textit{[referring to the catch-box microphone]} we will use that. So, okay, so fundamental --- so we we said two theorems that are um that are indeed important. Perhaps the most representative I would agree and and I saw many people saying yes, so it's the Fundamental Theorem of Calculus, right? And it's this \textit{[points at Equation \ref{eq:ftc_1d}]}. Is it correct? Yeah, looks like, no?

Now, to to unfold and to write and to understand every detail about this equation over here, you needed to introduce a lot of things. You needed to define what is that (i.e., the derivative $f'(x)$), right? What is this (i.e., the integral $\int_a^b \dots dx$)? And even --- and you spent a lot of time, right? What is that (i.e., the real number $\mathbb{R}$)? It was not --- it was not easy. Yeah, there is some mistake? \textit{[Corrects the board]} Good, well spotted. $f(b) - f(a)$. Now? 

Okay so this is a fairly compact theorem, yet the ingredients you had to learn to unfold all of this and to understand the meaning of every symbol here were a bit of work. Now, how will Analysis 2 look like? So this will be calculus, still calculus, in uh more than one variable. Well, since you already know one, let's say at least two variables, because otherwise it's already covered by Analysis 1.
\end{spoken-clean}

\begin{math-stroke}[Preview: The Fundamental Theorem in Higher Dimensions]
Analysis II extends these concepts to multiple variables, typically in $\mathbb{R}^n$.
\[
\text{Analysis 2: Calculus in } \ge 2 \text{ Variables}
\]
The higher-dimensional counterparts to the FTC are known by various names:
\begin{itemize}
    \item Gauss's Theorem
    \item Stokes's Theorem
    \item Divergence Theorem
    \item Green's Theorem
\end{itemize}

\begin{center}
\begin{tikzpicture}[scale=1.2]
    % This TikZ diagram is extracted from the blackboard at 00:09:30
    % 3D Axes
    \draw[->, thick, gray!80] (0,0,0) -- (3,0,0) node[right] {$y$};
    \draw[->, thick, gray!80] (0,0,0) -- (0,3,0) node[above] {$z$ (or $x_3$)};
    \draw[->, thick, gray!80] (0,0,0) -- (0,0,4) node[below left] {$x$ (or $x_1$)};

    % The "Potato" Domain Omega
    \draw[thick, MidnightBlue, fill=MidnightBlue!10, opacity=0.6] 
        (1,1,1) to[out=20,in=160] (3,1.5,1) 
        to[out=-20,in=70] (3.5,0.5,3) 
        to[out=160,in=-20] (1.5,0,3) 
        to[out=70,in=-20] cycle;
    \node[MidnightBlue] at (2,0.8,2) {\Large $\Omega$};
    
    % The Surface S (The "Skin")
    \node[BurntOrange] at (3.8,1.2,2) {Surface $S = \partial \Omega$};
    
    % Vector Field F (The "Cereal Field")
    \foreach \x in {0.5, 1.5, 2.5}
        \foreach \y in {0.5, 1.5, 2.5}
            \draw[->, BrickRed, opacity=0.4] (\x,\y,2) -- (\x+0.3, \y+0.4, 2.2);
    \node[BrickRed] at (0.5, 2.8, 2) {Vector Field $F = (f_1, f_2, f_3)$};

    % Normal Vector nu
    \draw[->, very thick, RoyalPurple] (3.2, 1.1, 1.8) -- (3.8, 1.4, 1.8) node[right] {$\nu$};
\end{tikzpicture}
\end{center}

The generalized theorem relates the integral of a derivative over a volume to the flux across its boundary:
\begin{equation} \label{eq:ftc_nd}
\underbrace{\int_{\Omega} \sum_{i=1}^n \frac{\partial f_i}{\partial x_i} \, dx}_{\text{Integral of "Derivative" over Volume}} = \underbrace{\int_{S} F \cdot \nu \, dS}_{\text{Flux across Boundary Surface}}
\end{equation}
\end{math-stroke}

\begin{spoken-clean}[00:06:23 - 00:07:45]
And some representative theorem that would be the equivalent, so the counterpart of this one, is the Gauss --- it has many names and it's always the same theorem, so but --- so it's the Gauss, Stokes, uh Divergence, sometimes it's called the Divergence Theorem, uh and that's it, I think Green, sometimes Green Theorem. 

And now I will write the theorem and you will not understand exactly what each thing means, I will tell you intuitively, and then the goal of the lecture will be at the end of the lecture you will be able to understand in full precision everything that I wrote here. Okay? So the theorem goes as follows. So we have --- we have the space, the space, let's say I will put it in the space $\mathbb{R}^3$. Three coordinates of space, right? This thing: $x, y, z$ or $x_1, x_2, x_3$. Sometimes $x_1, x_2, x_3$. Right? 

And then in the --- in this space we are given a domain, we are given some or some surface that encloses some interior, so this --- this is the interior of the surface and the exterior. And this is a surface. Right? The surface I call it $S$. What is a surface? It could be a sphere, it could be a sort of the \qt{skin of a potato}. It's a surface. Two-dimensional object. It has and --- and what I'm --- this kind of surface will have a interior and exterior, and the interior of the surface, so the --- if the surface is the skin of the potato, the --- the filling of the potato, the white part, right? is $\Omega$, I call it $\Omega$.
\end{spoken-clean}

\begin{didactic-insight}[The Potato Analogy]
The lecturer uses the \qt{potato} analogy to ground abstract topological concepts. The \qt{skin} represents the 2D boundary surface $S$, while the \qt{filling} represents the 3D interior volume $\Omega$. This physical metaphor helps students visualize the relationship between a set and its boundary in higher dimensions.
\end{didactic-insight}

\begin{spoken-clean}[00:07:45 - 00:09:01]
And then the equivalent of this theorem in this setup is this: Given a vector field --- now you don't know what it is a vector field, let me tell you, $F = (f_1, f_2, f_3)$. These are three --- three functions of space, right? At each point of space you have three numbers. Here you have three numbers, here you have three numbers, yeah. This is a vector field $F$. And how we do we draw vector fields? Uh, we draw it like this, right? It's like a --- a field of some cereal, right? but in every point you have an arrow instead of a --- a plant. 

Okay, so this is the vector field $F$. And the theorem says that the integral in $\Omega$ of the derivative of the field, the $i$-component with respect to the $x_i$ variable, and you sum over $i$, this is equal to the flux of the field across the boundary, across the $S$. And it's called like this: $\int_S F \cdot \nu \, dS$. This is the same theorem you have got here in one dimension, in several variables.
\end{spoken-clean}

\begin{spoken-clean}[00:09:01 - 00:10:45]
Now, the small problem is that maybe some of you in physics or --- or in previous courses have seen some of this notation, probably not, I don't expect you to --- we --- and we will need to introduce each of these notations. These partial derivatives, what is a domain, what is a surface? What is this surface mathematically? How do you model a surface? How do you compute a integral over a surface? You know how to compute a integral over the real line, right? over a segment, but over a surface? This will need some definition and some work. And these objects that are the partial derivatives. 

So here in --- in one equation you have a bit the goal and the syllabus of the course. How will --- how will we proceed? So first we will start looking at the space. In Analysis 1 you did everything in $\mathbb{R}$. Now we will need to introduce $\mathbb{R}^3$, but we are mathematicians so we will --- well, some of us. Uh, some of you are physicist or are proto-physicist, would like to become physicist. So we will use this because we don't want to spend time proving the same theorem in --- in the plane $\mathbb{R}^2$, then in $\mathbb{R}^3$ in the space, and then when you go to do general relativity you need $\mathbb{R}^4$, right? We do everything in $\mathbb{R}^n$. $n$ is the dimension.
\end{spoken-clean}

\begin{spoken-clean}[00:10:45 - 00:11:00]
So what is $\mathbb{R}^n$? Is the set of $n$-tuples of real numbers. This is the definition, okay? And what is this? This is the same as the set of points we will denote a point $x$, but now $x$ will carry $n$ coordinates: $x_1, x_2, \dots, x_n$. Now this is the point. And these $x_i$ are real numbers. This you know probably because you are doing well or you have done linear algebra, right? So this should be quite familiar. But still it's good that we recall. Okay.
\end{spoken-clean}

[SYSTEM] Segment complete. Please prompt "Continue" for the remainder of the segment.```
%% ==========================================
% Session Info:
% System Prompt loaded: True
% History loaded: True
% 
% AI Studio User Prompt:
% continue
% ==========================================

```latex
\begin{spoken-clean}[00:11:00 - 00:12:30]
Now, the first part of the course we will be studying $\mathbb{R}^n$ and the properties of $\mathbb{R}^n$. The same a bit as in $\mathbb{R}$ you studied the --- the axiom of the supremum, right? Or the --- the fact that you can take uh limits of Cauchy sequences, these kind of facts, right? The same we will need to study in $\mathbb{R}^n$. Then there will be some new things or --- or things that you could also study in $\mathbb{R}$ but become very trivial and now will be much more interesting, is that the topology of $\mathbb{R}^n$. 

So first, the $\mathbb{R}^n$ as a space, as Euclidean space. So this is a bit the syllabus. Two, the topology of $\mathbb{R}^n$. You will see what does this mean in detail, the topology. Then three, this object over here \textit{[points at the blackboard]}: partial derivatives. So when you have --- when you have several variables, you can differentiate a quantity with respect to each of them, right? And differential and so --- the differentiation of functions of several variables. Then integration. Integration in $\mathbb{R}^n$, here, in a domain of $\mathbb{R}^n$, and on surfaces. 

But because we want to learn also how to integrate on surfaces, we will do a bit of introduction to differential geometry. Introduction to surfaces and \qt{submanifolds}. Now, I know it's a lot of words. But we will introduce one by one. So surface is a usual two-dimensional object, right? Typically smooth. When it is a three-dimensional object in $\mathbb{R}^4$, for instance, we call this a submanifold. You can also call it surface but --- but it's more precise the word submanifold. Or --- or or some two-dimensional object in a five-dimensional space. Any combination of dimensions. And then when you know all of this, you will need to know how to integrate. Integration over $\mathbb{R}^n$ and surfaces.
\end{spoken-clean}

\begin{math-stroke}[The Analysis II Syllabus]
The course is structured into five primary logical blocks:
\begin{enumerate}
    \item $\mathbb{R}^n$ as Euclidean Space
    \item Topology of $\mathbb{R}^n$
    \item Partial Derivatives and Differentiation
    \item Introduction to Surfaces and \qt{Submanifolds}
    \item Integration over $\mathbb{R}^n$ and Surfaces
\end{enumerate}

\begin{explanation-of-steps}
The ultimate goal of the course is to prove the generalized Fundamental Theorem of Calculus (Equation \ref{eq:ftc_nd}). To reach this, we must first understand the geometry of the space ($\mathbb{R}^n$), the behavior of functions within it (Topology and Differentiation), and finally how to measure volumes and \qt{skins} (Integration).
\end{explanation-of-steps}
\end{math-stroke}

\section{The Role of Proofs in Analysis}

\begin{spoken-clean}[00:12:30 - 00:13:50]
And hopefully then you will know all of these elements. This element --- this is the normal vector, I can tell you now what is this. This is just the --- the vector, the unit vector, the vector of length one that is normal to the surface, right? If you have a surface, this $\nu$ is just the normal to the surface. So this is easy. Um, but still we will --- we will define this in --- in detail. And after we know all of this, we will be able to prove that equation over here and work with it. And see why it is useful, right? For many things. Okay. 

So this is a rough plan of the course. Now I want to --- I wanted to do a very small survey about your attitude towards proofs, okay? This is important to mention. So, I will put two columns. Column left is: I find proofs --- so proofs, I mean rigorous, long, very careful mathematical proofs with all details. These I mean proofs. So, here I will say: I find proof um essential, illuminating. You see with the --- with this sound I don't need to run the surface \textit{[referring to the chalk sound]}, the survey. Um, essential, illuminating, um useful to clarify, um the best part of any course. And on this column over here I will put: I find the proof unnecessary, confusing. Any other idea? Um, I don't know when I have finished the proof or not, right? Because it looks like it sounds convincing but --- mysterious, no clear rules of the game.
\end{spoken-clean}

\begin{math-stroke}[Survey: Attitude Towards Proofs]
The lecturer contrasts two common perspectives on mathematical proofs:
\begin{center}
\begin{tabular}{l | l}
\textbf{Positive Perspective} & \textbf{Negative Perspective} \\ \hline
Essential & Unnecessary \\
Illuminating & Confusing \\
Useful to clarify & Mysterious \\
The best part of the course & No clear rules of the game
\end{tabular}
\end{center}
\end{math-stroke}

\begin{spoken-clean}[00:13:50 - 00:15:00]
And now, okay, this is extreme, I'm making this on purpose. This is extreme, right? This is extreme. I know that any --- most of you would be maybe in the middle. But I would like to know which of you lean more towards this opinion or more towards this opinion, okay? So, let's start with this one. So who is here? Who loves proofs and... \textit{[Students cheer and raise hands]} Okay. And here? \textit{[Students cheer and raise hands]} Okay, it's interesting. So it's roughly --- there are a bit more people that love proofs, but --- but okay. So there is no correct answer, right? So I would be --- well, depending on where --- in between. 

So why am I asking this? So does it coincide? Is it --- does it have some correlation with physicist and mathematician or not? Sometimes? What do you think? It could make sense, not always. There are physicist that... but it could make sense. So I think that um about this discussion of proofs, I wanted to take ten minutes today to tell you two things. First, why I believe --- many people can have many answers --- but why I believe that the proofs are important and what will we test in this course, right? Because it's --- it's a bit correlated of course. 

The first thing is that in --- in algebra many times you have a theorem that is very beautiful, very clean, and you can apply in many situations without modifying it, right? In analysis, we will see that, sometimes the theorem that you learned doesn't exactly apply to the situation you --- you need. Then you need to know how to modify a little bit the theorem. And the only way to do that is to know the proof and then modify a bit some step here and some step there.
\end{spoken-clean}

\begin{didactic-insight}[Computable Truth and the \qt{Lean} Prover]
The lecturer concludes by noting that in the age of AI, proofs are increasingly viewed as \emph{computable}. He references \qt{Lean}, a functional programming language and proof assistant. This shifts the definition of a proof from a \qt{convincing argument} or an \qt{opinion} to a piece of code that either compiles (is correct) or does not. This modern perspective reinforces the necessity of rigorous detail.
\end{didactic-insight}

\begin{meta-note}[End of Introduction]
The lecturer concludes the introductory remarks and prepares to begin the first formal topic of the syllabus: the structure of $\mathbb{R}^n$.
\end{meta-note}

[SYSTEM] Video complete.
```

%% \documentclass[11pt]{article}

% % ==========================================
% % PREAMBLE & PACKAGES
% % ==========================================
% \usepackage[utf8]{inputenc}
% \usepackage[T1]{fontenc}
% \usepackage{amsmath, amssymb, amsthm, mathtools}
% \usepackage{xcolor}
% \usepackage[margin=1in]{geometry}
% \usepackage{hyperref}
% \usepackage{tcolorbox}
% \tcbuselibrary{skins, breakable}

% % Fail-Safe TikZ Configuration
% \usepackage{tikz}
% \usetikzlibrary{calc, arrows.meta, positioning, shapes.geometric, decorations.pathreplacing}
% \usepackage{pgfplots}
% \pgfplotsset{compat=1.18}

% % Custom Blackboard Colors
% \definecolor{profwhite}{RGB}{255, 255, 255}
% \definecolor{proforange}{RGB}{255, 165, 0}
% \definecolor{profyellow}{RGB}{255, 255, 0}
% \definecolor{profgreen}{RGB}{0, 255, 0}
% \definecolor{profred}{RGB}{255, 0, 0}
% \definecolor{profblue}{RGB}{0, 150, 255}

% % Math operators
% \DeclareMathOperator{\id}{id}
% \DeclareMathOperator{\Vol}{Vol}

% % Standard Math Environments
% \newtheorem{theorem}{Theorem}
% \newtheorem{lemma}[theorem]{Lemma}
% \theoremstyle{definition}
% \newtheorem{definition}{Definition}

% % ==========================================
% % CUSTOM META-ENVIRONMENTS FOR AI PARSING
% % ==========================================

% % 1. Clean Spoken Protocol (Grammar-corrected transcript)
% \newtcolorbox{spoken_clean}[1][]{
%     colback=blue!5!white,
%     colframe=blue!75!black,
%     title=\textbf{Cleaned Spoken Protocol [Time: #1]},
%     breakable,
%     fonttitle=\bfseries,
%     fontupper=\sffamily
% }

% % 2. AI Stroke-by-Stroke Note (Chronological board tracking)
% \newtcolorbox{ainote}[1][]{
%     colback=gray!10!white,
%     colframe=gray!50!black,
%     title=\textbf{Board State \& Physical Actions: #1},
%     breakable,
%     fonttitle=\bfseries
% }

% % 3. Best Teacher / Didactic Insight (Meta-commentary on the "Why")
% \newtcolorbox{didactic_insight}[1][Pedagogical Concept]{
%     colback=profgreen!5!white,
%     colframe=profgreen!60!black,
%     title=\textbf{Best Teacher Insight: #1},
%     breakable,
%     fonttitle=\bfseries
% }

% % 4. Semantic Formula Box (For highlighted chalk formulas)
% \newtcolorbox{orangeformula}[1][Highlighted Board Formula]{
%     colback=profwhite,
%     colframe=proforange,
%     fonttitle=\bfseries\sffamily,
%     coltitle=proforange,
%     coltext=proforange,
%     title=#1,
%     boxrule=1.5pt,
%     arc=4pt,
%     center title
% }

% \begin{document}

\section{Hyper-Detailed Pedagogical Blueprint \\ Real Analysis: Change of Variables \& Rigorous Foundations}

\emph{Director's Cut \& AI Extraction Protocol}

\textit{Lecture Segment: 00:00 -- 15:00}

\begin{didactic_insight}[The "Wall" of Multivariable Calculus]
This segment serves as the bridge between computational calculus and rigorous Real Analysis. The professor is actively shifting the students' mindset away from "memorizing integration formulas" and towards "understanding the geometric distortions of space via linear algebra (Jacobians)." Notice how carefully he handles domains and inverse mappings—this rigor is the core of the discipline.
\end{didactic_insight}

% ==========================================
% SECTION 1: THE "FINAL BOSS" SETUP
% ==========================================
\section{Reminder: The Setup for the Change of Variables}

\begin{spoken_clean}[00:00:11 - 00:01:29]
I hope you all had a great holiday and have recovered from the first half of the semester so we can continue with our work! Let me start by reminding you of what we did last time. In our previous lecture, we stated and began to prove what we playfully call the "final boss of integration": the Change of Variables formula. Let us review the setup. 

As always, we begin with our domain $U \subset \mathbb{R}^n$. We apply a diffeomorphism—which we call $\Phi$—that maps $U$ to another domain $V$. 
\end{spoken_clean}

\begin{ainote}[Board State 1: The Diffeomorphism]
The professor writes the title "Change of vars." and begins defining the mathematical environment on the far left board. He writes the mapping $U \to V$ and explicitly writes "diffeo" in orange chalk beneath the arrow.
\end{ainote}

\textbf{Mathematical Setup (Stroke-by-Stroke):}
Let $U, V \subset \mathbb{R}^n$. We define the mapping between these spaces:
\[
U \xrightarrow[\text{\textcolor{proforange}{diffeo}}]{\Phi} V
\]

\textbf{Redundant Explanation of Step:} 
We establish a transformation $\Phi$ mapping a domain $U$ to a codomain $V$. The orange annotation "diffeo" indicates this mapping is a \textit{diffeomorphism}. A diffeomorphism is a heavily constrained, "well-behaved" mapping: it must be a bijection (one-to-one and onto), it must be differentiable everywhere, and its inverse $\Phi^{-1}$ must also be differentiable everywhere. This ensures the space is smoothly warped without tearing or folding.

\begin{center}
\begin{tikzpicture}[scale=1.2]
    % Domain U
    \draw[thick, fill=blue!10] (0,0) ellipse (1.5cm and 1cm);
    \node at (0, 0.5) {\Large $U$};
    \node[below] at (0, -1.2) {Original Space $\mathbb{R}^n$};

    % Domain V
    \draw[thick, fill=red!10, shift={(5,0)}] (0,0) to[out=45,in=135] (1.5,0.5) to[out=-45,in=45] (1,-1) to[out=-135,in=-45] (-1,-0.5) to[out=135,in=-135] cycle;
    \node at (5, 0.5) {\Large $V$};
    \node[below] at (5, -1.2) {Distorted Space $\mathbb{R}^n$};

    % Mapping Arrow
    \draw[->, very thick, profblue] (1.8, 0) to[bend left=20] node[midway, above] {\textbf{Diffeomorphism $\Phi$}} (3.2, 0);
\end{tikzpicture}
\end{center}

\begin{spoken_clean}[00:01:29 - 00:02:45]
We then consider a Jordan measurable set $A \subset \mathbb{R}^n$, with the strict condition that its closure, $\overline{A}$, must be entirely contained within the domain $U$ of our diffeomorphism. Next, we need a continuous function to integrate. We define a function $f$ mapping from the closure $\overline{A}$ to the real numbers $\mathbb{R}$.
\end{spoken_clean}

\begin{ainote}[Board State 2: Domain Correction]
The professor writes $A \subset \mathbb{R}^n$ J. m. st. $\overline{A} \subset U$. When defining the function $f$, he initially says "from $U$ or from $A$ to $\mathbb{R}$", but immediately catches his imprecision and writes $\overline{A} \to \mathbb{R}$ continuous.
\end{ainote}

\textbf{Mathematical Setup (Continued):}
We restrict our attention to a specific subset of $U$:
\[
A \subset \mathbb{R}^n \text{ is Jordan measurable, such that } \overline{A} \subset U
\]
We define our integrand function:
\[
f: \overline{A} \to \mathbb{R} \quad (\text{continuous})
\]

\textbf{Redundant Explanation of Step:}
Why the closure $\overline{A}$? By requiring the \textit{closure} of $A$ to be strictly contained within the open set $U$, we ensure that the boundary of $A$ behaves nicely under the mapping $\Phi$, and we avoid pathological singularities that could exist on the absolute edge of the domain $U$. We then ensure $f$ is continuous over this closed, bounded region, guaranteeing Riemann integrability.

% ==========================================
% SECTION 2: THE CLUNKY THEOREM
% ==========================================
\section{The Clunky Formulation of the Theorem}

\begin{spoken_clean}[00:02:45 - 00:03:50]
With this setup, the theorem states that the integral of $f(x)$ over the region $A$ is equivalent to the integral evaluated over the mapped image, $\Phi(A)$. However, to calculate this, we must evaluate our function $f$ at the pre-image, $\Phi^{-1}(y)$. Furthermore, the differential element $dy$ must be corrected by the distortion of the volume, which is exactly the absolute value of the determinant of the Jacobian of $\Phi$, evaluated at $\Phi^{-1}(y)$.
\end{spoken_clean}

\begin{ainote}[Board State 3: The Clunky Formula]
The professor writes out the full Change of Variables formula exactly as it was derived in the previous lecture. It is notationally heavy due to the inverse mappings inside the arguments.
\end{ainote}

\textbf{The Initial Theorem (Inline Annotated):}
\begin{equation} \label{eq:cov_clunky}
\underbrace{\int_A f(x) \, dx}_{\text{Integral in Original Space}} = \int_{\Phi(A)} f(\underbrace{\Phi^{-1}(y)}_{\text{Pre-image under }\Phi}) \, \underbrace{\frac{1}{|\det J\Phi(\Phi^{-1}(y))|}}_{\text{Correction for Volume Distortion}} \, dy
\end{equation}

\textbf{Redundant Explanation of Step:}
To evaluate the integral in the distorted space $V$ (where the variable is $y$), every $x$ in the original function must be replaced with the mapping that takes $y$ back to $x$, which is $\Phi^{-1}(y)$. Because the transformation $\Phi$ stretches and squishes space, the resulting "volume pixels" ($dy$) are distorted. To preserve the total "mass" of the integral, we must divide by the scaling factor of that distortion. The scaling factor of a linear transformation is the determinant of its matrix—hence, we divide by the determinant of the Jacobian matrix of $\Phi$.

\begin{spoken_clean}[00:03:50 - 00:06:50]
I previously explained the geometric intuition behind why this theorem is true, but we didn't have time to complete the formal proof. I propose that we leave the proof for now. Instead, we will look at some concrete examples to see how we apply this theorem in actual calculations. Once you see it in action, you will understand why it is worth the effort to prove it. We will revisit the formal proof in two weeks during our revision buffer week. 

*(The professor then spends a few minutes discussing an upcoming interactive teaching session on Thursday with guest coaches from ETH. He asks students to attend in person for the interactive exercises.)*
\end{spoken_clean}

% ==========================================
% SECTION 3: REFINING THE NOTATION
% ==========================================
\section{Rewriting the Theorem via the Inverse Map}

\begin{spoken_clean}[00:06:50 - 00:08:20]
Let us rewrite this theorem into a format that is much easier to use in practice. Instead of constantly writing $\Phi^{-1}$ and dragging it through our equations, let us define a new mapping $\Psi$ to be exactly equal to the inverse, $\Phi^{-1}$. 

If we substitute this into our theorem, the integration region $\Phi(A)$ remains, but the function evaluation $f(\Phi^{-1}(y))$ simply becomes $f(\Psi(y))$. But what happens to that clunky determinant in the denominator? 
\end{spoken_clean}

\begin{ainote}[Board State 4: Defining Psi]
The professor moves to the center board to clean up the notation. He defines $\Psi$ and prepares to algebraically manipulate the Jacobian fraction.
\end{ainote}

\textbf{Simplifying the Notation:}
Let us define the inverse map explicitly:
\[
\Psi = \Phi^{-1}
\]
Substituting this directly into Equation \ref{eq:cov_clunky}, we get an intermediate step:
\[
\int_A f(x) \, dx = \int_{\Phi(A)} f(\Psi(y)) \frac{1}{|\det J\Phi(\Psi(y))|} \, dy
\]

\begin{spoken_clean}[00:08:20 - 00:10:45]
Remember the lemma we proved regarding the differential of an inverse function. If we compose $\Phi$ with its inverse $\Psi$, we get the identity map. 

By applying the chain rule, the differential of this composition evaluated at a point $y$ is the differential of $\Phi$ multiplied by the differential of $\Psi$. Because the differential of the identity map is just the identity matrix, this means that the Jacobian matrix of $\Phi$ multiplied by the Jacobian matrix of $\Psi$ equals the unit (identity) matrix.

Because the determinant is a multiplicative property for matrices, the determinant of one matrix must be the exact inverse (reciprocal) of the other. Therefore, instead of dividing by the determinant of $J\Phi$, we can simply multiply by the determinant of $J\Psi$!
\end{spoken_clean}

\begin{ainote}[Board State 5: Chain Rule Derivation]
The professor formally derives the relationship between the determinants on the board, line by line, mapping the abstract function composition to concrete matrix algebra.
\end{ainote}

\textbf{Derivation of the Jacobian Reciprocal Relation:}
\begin{align*}
\underbrace{\Phi \circ \Psi = \id}_{\text{Function Composition}} &\implies \underbrace{D\Phi_{\Psi(y)} \circ D\Psi_y = \id}_{\text{Taking the Differential (Chain Rule)}} \\
&\implies \underbrace{J\Phi(\Psi(y))}_{\text{Matrix 1}} \cdot \underbrace{J\Psi(y)}_{\text{Matrix 2}} = \underbrace{I_n}_{\text{Identity Matrix}}
\end{align*}

\textbf{Redundant Explanation of Step:}
The chain rule in multivariable calculus dictates that the derivative of a composition of functions is the matrix product of their Jacobian matrices. Since the composition yields the identity function ($y \mapsto y$), the product of their Jacobians must yield the identity matrix $I_n$. 

Taking the determinant of both sides, and using the property that $\det(AB) = \det(A)\det(B)$ and $\det(I_n) = 1$:
\begin{align*}
\det\Big( J\Phi(\Psi(y)) \cdot J\Psi(y) \Big) &= \det(I_n) \\
\det(J\Phi(\Psi(y))) \cdot \det(J\Psi(y)) &= 1 \\
\implies \underbrace{\frac{1}{\det J\Phi(\Psi(y))}}_{\text{Clunky Denominator}} &= \underbrace{\det J\Psi(y)}_{\text{Clean Numerator}}
\end{align*}

\begin{didactic_insight}[The Power of Linear Algebra]
The professor takes a moment here to show how beautifully linear algebra serves calculus. The geometric reality that "stretching space forward by factor $k$, then pulling it backward, returns it to normal" is flawlessly encoded in the algebraic fact that the determinants of inverse matrices multiply to 1.
\end{didactic_insight}

\begin{spoken_clean}[00:10:45 - 00:12:15]
This gives us a highly practical, informal way to remember the substitution rule. Just like in Analysis 1, when you substitute $x = \Psi(y)$, you must replace the differential. But instead of replacing $dx$ with the 1D derivative $\Psi'(y) \, dy$, you replace it with the absolute determinant of the Jacobian matrix, $|\det J\Psi(y)| \, dy$.
\end{spoken_clean}

\begin{ainote}[Board State 6: The Orange Rule]
To emphasize its importance for practical computation, the professor writes this finalized substitution rule on the board in bright orange chalk.
\end{ainote}

\begin{orangeformula}[The Practical Substitution Rule]
When substituting $x = \Psi(y)$, the differential transforms as:
\[
\underbrace{dx}_{\text{Original Volume}} = \underbrace{|\det J\Psi(y)|}_{\text{Scaling Factor}} \, \underbrace{dy}_{\text{New Volume}}
\]
Yielding the clean theorem:
\[
\int_A f(x) \, dx = \int_{\Phi(A)} f(\Psi(y)) |\det J\Psi(y)| \, dy
\]
\end{orangeformula}

% ==========================================
% SECTION 4: EXAMPLE - AREA OF A UNIT DISK
% ==========================================
\section{Example: Area of a Unit Disk \& The Definition of $\pi$}

\begin{spoken_clean}[00:12:15 - 00:15:00]
Let us compute a concrete example. But before we do, I must ask you: how was the constant $\pi$ officially defined in your Analysis 1 course? 

*(A student answers correctly: "It is defined as the unique number in the open interval $(0, 4)$ such that the sine of $\pi$ is zero.")*

Exactly! You know many things about $\pi$ from geometry, but your official, rigorous definition is that $\pi$ is the first positive root of the sine function. 

The sine and cosine functions themselves are defined purely analytically (usually via power series), and their defining properties are that the derivative of sine is cosine, the derivative of cosine is negative sine, $\sin(0) = 0$, and $\cos(0) = 1$. Those relations are definitive.
\end{spoken_clean}

\begin{ainote}[Board State 7: Defining Pi rigorously]
The professor steps to the far right board to lay down the absolute axioms of Analysis 1 before computing the area of a circle. He wants to prove the area is $\pi$ *without* using circular geometric logic.
\end{ainote}

\textbf{Rigorous Axioms from Analysis 1:}
The trigonometric functions are defined purely by their analytical properties (not triangles):
\begin{align}
\sin'(x) &= \cos(x) \label{ax:sin_deriv} \\
\cos'(x) &= -\sin(x) \label{ax:cos_deriv} \\
\sin^2(x) + \cos^2(x) &= 1 \label{ax:pythagoras}
\end{align}
The constant $\pi$ is defined analytically as the unique root:
\[
\sin(\pi) = 0 \quad (\text{for the first positive root})
\]

\textbf{Redundant Explanation of Step:}
By establishing these purely algebraic definitions, the professor is ensuring that when we compute the integral of a disk and get the symbol $\pi$, we are arriving at it through rigorous calculus and limits, not by relying on a high-school geometric formula ($A = \pi r^2$) which assumes the conclusion we are trying to prove.

\begin{didactic_insight}[Avoiding Circular Logic]
This is a brilliant teaching moment. It is very easy for students to take geometry for granted. By forcing them to recall that $\pi$ is just "the root of a specific power series", the professor elevates the upcoming Change of Variables computation from a "boring high-school geometry problem" into a "profound demonstration of analytical power." 
\end{didactic_insight}

\textit{Note: The video segment ends right as the professor finishes writing these trigonometric axioms on the board. In the next segment, he will define the polar coordinate mapping $\Psi(y_1, y_2)$ and use the Orange Substitution Rule to calculate the Jacobian determinant and prove the area of the unit disk is exactly $\pi$.}


%\include{content/4-15-monday-first_half-part2}  
%\include{content/4-15-monday-first_half-part3}

%\section{Monday, April 15th: Second Half of the Lecture \texorpdfstring{\\}{\&}
    Cavalieri's Principle: Statement and Proof via Dyadic Cubes}

\emph{AI Extraction Protocol}

\textit{Segment: 00:00 -- 15:15}


% \begin{meta-note}[Previous Board State]
% As the video begins, there is a pre-existing formula on the left blackboard from the previous session:
% $\mu_3(B^3_1) = \mu_3(\Psi(D)) = \int_{(0,1) \times (0,2\pi) \times (\frac{\pi}{2}, \pi)} y_1^2 \cos y_3 \, dy = \dots$ with the words ``Slicing idea'' written above it. On the right board, the heading ``Cavalieri principle'' is written alongside a 3D sketch of a body $A \subset \mathbb{R}^3$ being sliced into parallel disks, with the formula: $\text{volume}(A) = \sum_{\text{all slices}} \dots$ 
% \end{meta-note}

% ==========================================
% SECTION 1: GEOMETRIC SETUP OF CAVALIERI'S PRINCIPLE
% ==========================================
\subsection{Geometric Setup of Cavalieri's Principle}

\begin{spoken-clean}[00:00:08 - 00:02:20]
Let me continue with the statement. We are trying to formalize this theorem that I informally described as ``cutting the potato.'' But now we are trying to be more serious, and I will introduce the proper notations. 

This theorem, you will see, is a bit the opposite of the Change of Variables. It is a very simple statement, but it has a very hard proof. The setup involves very complicated statements, but the proof itself is actually very easy, right? 

Let me do a drawing of the measure space. We have our set $A$, right, that is contained in a box. Here we will put the X-axis and the Y-axis. The X-axis actually represents $\mathbb{R}^k$, and the Y-axis represents $\mathbb{R}^{n-k}$. The whole space is $\mathbb{R}^n$, which is the product of the two. 

Okay, now we are saying that our set is contained in some bounding box. It is a very large box, going from $-c$ to $c$. Now, let's look inside this box. If we pick some point $x$ here on the X-axis, I will consider the slice. The relevant piece of the set is this vertical slice over here, right? This is what I am calling $A_x$. The set $A_x$ will be the projection of this slice onto the vertical axis.
\end{spoken-clean}

\begin{nice-box}[Geometric Visualization Setup]
The professor draws the bounding box and the sets $A$ and $A_x$ to geometrically define the slicing operation.
\end{nice-box}

\begin{math-stroke}[Geometric Visualization Setup]
\begin{center}
\begin{tikzpicture}[scale=1.5]
  % Axes with defined dimensions
  \draw[->, thick, profgreen] (0,0) -- (4,0) node[right] {$x \in \mathbb{R}^k$};
  \draw[->, thick, profgreen] (0,0) -- (0,3) node[above] {$y \in \mathbb{R}^{n-k}$};
  \node[above right, profgreen] at (3.5, 2.8) {$\mathbb{R}^n$};
  
  % Bounding Box [-c, c]^n
  \draw[dashed, gray] (0.5, 0.5) rectangle (3.5, 2.5);
  \node[below, gray] at (0.5, 0) {$-c$};
  \node[below, gray] at (3.5, 0) {$c$};

  % The Set A
  \draw[thick, profblue, fill=profblue!5] (1.5,1.5) to[out=45,in=135] (2.5,2.0) to[out=-45,in=45] (3.0,1.2) to[out=-135,in=-45] (2.0,0.8) to[out=135,in=-135] cycle;
  \node[profblue] at (2.4, 1.5) {$A$};

  % The Slice
  \draw[thick, profred] (1.8, 0.95) -- (1.8, 2.05);
  \node[below, profgreen] at (1.8, 0) {$x$};
  \draw[dashed, profgreen] (1.8, 0) -- (1.8, 0.95);
  
  % Projection Ax
  \draw[thick, proforange] (0, 0.95) -- (0, 2.05);
  \node[left, proforange] at (0, 1.5) {$A_x$};
  \draw[dashed, proforange, ->] (1.8, 2.05) -- (0, 2.05);
  \draw[dashed, proforange, ->] (1.8, 0.95) -- (0, 0.95);
\end{tikzpicture}
\end{center}
\end{math-stroke}

\begin{spoken-clean}[00:02:20 - 00:04:50]
So, you take this set, and you project what? You take the $y$-coordinates. Given the $x$, you take the $y$-coordinates of the pairs $(x,y)$ that belong to $A$. Since $x$ is fixed, what you are looking at is: you go to $x$, you look at which $y$-coordinates you have inside the set, and this is the set $A_x$. 

Now, this set here is very nice because in the drawing it is just a 1D interval, right? But we might be in very high dimensions. So $A_x$ is just a subset of $\mathbb{R}^{n-k}$ a priori, and it is a bounded set. We know that it is strictly contained within $(-c, c)^{n-k}$. 

Because we know the entire set $A$ is Jordan measurable, we might hope the slice is too. But maybe for some pathological slices, the set is not measurable because the outer and inner measures do not coincide. However, we can surely consider the marginals! We can always measure the outer measure $\mu_{n-k}^{\text{out}}$ and the inner measure $\mu_{n-k}^{\text{in}}$. 

The measure from the inside and the outside of the set is always well-defined, and this defines the functions $\overline{\varphi}(x)$ and $\underline{\varphi}(x)$, right? So, this is the setup.
\end{spoken-clean}

\begin{math-stroke}[Formal Definition of the Slice and Marginals]
Let $\mathbb{R}^n = \mathbb{R}^k \times \mathbb{R}^{n-k}$, with coordinates $(x,y) \in \mathbb{R}^n$. \\
Let $A \subset \mathbb{R}^n$ be Jordan measurable, with $A \subset (-c, c)^n$. \\
Define the slices: for $x \in (-c, c)^k$,
\[
A_x = \{ y \in \mathbb{R}^{n-k} \mid (x,y) \in A \} \subset (-c, c)^{n-k}
\]
We define the upper and lower marginal densities $\overline{\varphi}, \underline{\varphi} : [-c, c]^k \to \mathbb{R}$ as:
\begin{align*}
\overline{\varphi}(x) &= \mu_{n-k}^{\text{out}}(A_x) \\
\underline{\varphi}(x) &= \mu_{n-k}^{\text{in}}(A_x)
\end{align*}

\begin{explanation-of-steps}
Even if a solid body $A$ is perfectly measurable, a specific 1D slice could theoretically intersect a fractal boundary or a cloud of points, making that exact slice non-measurable. To bypass this entirely, the marginals are rigorously defined using the universally safe inner and outer measures.
\end{explanation-of-steps}
\end{math-stroke}

% ==========================================
% SECTION 2: THE THEOREM STATEMENT
% ==========================================
\section{The Theorem Statement}

\begin{spoken-clean}[00:04:50 - 00:07:30]
And the theorem says that under these assumptions, both functions $\overline{\varphi}$ and $\underline{\varphi}$ mapping from $[-c,c]^k \to \mathbb{R}$ are Riemann integrable. 

And the integral of them over the set $[-c, c]^k$—the integral of $\underline{\varphi}(x)\,dx$ is equal to the integral of $\overline{\varphi}(x)\,dx$, and both are equal to the measure of $A$, $\mu_n(A)$.

So, this integral is just you summing up the $\overline{\varphi}$. The $\overline{\varphi}$ is the measure of each slice (i.e., $\mu_{n-k}^{\text{out}}(A_x)$)! 
\end{spoken-clean}

\begin{student-question}
Maybe for some slices this measure is not well-defined because the outer and the inner measure do not coincide!
\end{student-question}

\begin{spoken-clean}[00:06:15 - 00:07:30]
Do not worry! Just put either the outer or the inner. I don't care, because when you integrate the outer, or the inner, or anything in between, you will always recover the exact measure of your set $A$. 

Okay, let's prove that. To do the proof, I will do another drawing.
\end{spoken-clean}

\begin{nice-box}[Highlighted Theorem]
\begin{theorem}[Cavalieri's Principle]
If $A \subset \mathbb{R}^n$ is Jordan measurable, then the marginal functions $\overline{\varphi}, \underline{\varphi} : [-c, c]^k \to \mathbb{R}$ are Riemann integrable, and:
\[
\int_{[-c,c]^k} \underline{\varphi}(x) \, dx = \int_{[-c,c]^k} \overline{\varphi}(x) \, dx = \mu_n(A)
\]
\end{theorem}
\end{nice-box}

% ==========================================
% SECTION 3: THE PROOF SETUP VIA DYADIC CUBES
% ==========================================
\section{The Proof Setup: Dyadic Cubes}

\begin{spoken-clean}[00:07:30 - 00:11:00]
Always when we want to prove something about Jordan measures, remember that we are approximating our set from the inside and the outside. We will take an approximation from the inside, right? I'm approximating from the inside, and then there will be the approximation from the outside. 

Let me describe a bit with a picture of what we are going to do in the proof, and later we will do the symbols. Since our set is Jordan measurable, we can take a small pixel size of $2^{-p}$. We can take an inner and an outer approximation over our set by dyadic cubes. So, by pixels! This is the inner, this is the outer. 

The difference between these two is less than epsilon. We can make the volume of the red outer approximation minus the volume of the yellow inner approximation less than epsilon for whatever precision we want. 

Now we will see what happens when we take a slice of this pixelated set. Let's say we take $x$ and we slice, right? What we see is that the slice of $A$ contains the slice of the yellow, which is contained in the slice of the red. And when you start to move $x$, whenever you move $x$ less than the size of the pixel, you are not jumping to a new cube!
\end{spoken-clean}

\begin{math-stroke}[Geometric Proof Setup: Inner and Outer Approximations]
\begin{center}
\begin{tikzpicture}[scale=1.5]
    % Axes
    \draw[->, thick, profgreen] (0,0) -- (4,0) node[right] {$x$};
    \draw[->, thick, profgreen] (0,0) -- (0,3) node[above] {$y$};
    
    % The Smooth Set A (White Outline)
    \draw[thick, gray, dashed] (1.5,1.5) to[out=45,in=135] (2.5,2.0) to[out=-45,in=45] (3.0,1.2) to[out=-135,in=-45] (2.0,0.8) to[out=135,in=-135] cycle;
    \node[gray] at (2.8, 1.5) {$A$};

    % Inner Approximation F (Yellow / Jagged)
    \draw[thick, profyellow] (1.7, 1.3) -- (1.7, 1.8) -- (2.3, 1.8) -- (2.3, 1.6) -- (2.6, 1.6) -- (2.6, 1.2) -- (2.0, 1.2) -- (2.0, 1.0) -- (1.8, 1.0) -- (1.8, 1.3) -- cycle;
    \node[profyellow, right] at (2.6, 1.4) {$F$ (Inner)};

    % Outer Approximation G (Red / Jagged)
    \draw[thick, profred] (1.2, 1.0) -- (1.2, 2.2) -- (1.6, 2.2) -- (1.6, 2.4) -- (2.8, 2.4) -- (2.8, 2.0) -- (3.3, 2.0) -- (3.3, 1.0) -- (2.3, 1.0) -- (2.3, 0.6) -- (1.5, 0.6) -- (1.5, 1.0) -- cycle;
    \node[profred, right] at (3.3, 2.1) {$G$ (Outer)};
    
    % The Vertical Slice showing step-nature
    \draw[thick, cyan] (2.1, 0) -- (2.1, 2.6);
    \node[below, cyan] at (2.1, 0) {$x$};
\end{tikzpicture}
\end{center}
\end{math-stroke}

\begin{spoken-clean}[00:11:00 - 00:15:15]
Because you are not jumping to a new cube, the marginal densities computed for the inner and outer approximations are dyadic step functions! 

This is the key observation. They are dyadic step functions in $x$, and they will approximate the Riemann integral. The dyadic step functions are what we use to approximate the integral. 

So, this is how we will try to proceed. Let's write this formally. The idea is that given a dyadic cube—I will take one of size $2^{-p}$ in $\mathbb{R}^n$. Let me call this $Q_p^n(z)$. So $Q_p^n(z) = z + 2^{-p}[0,1)^n$, right? Since there will be cubes of several dimensions, let me put the dimension here as an index $n$. 

Now, if $z$ is of the form $(a,b)$ where $a \in \mathbb{R}^k$ and $b \in \mathbb{R}^{n-k}$—meaning in this splitting I take a $z$ that has coordinates $a$ and $b$—then this $n$-dimensional cube $Q_p^n(z)$ is equal to $Q_p^k(a) \times Q_p^{n-k}(b)$, right? 

So, I am saying something very complicated to write, but very trivial geometrically. On the picture, I am saying that the $n$-dimensional cube over here is always the product of a base cube over here times the vertical cube over there!
\end{spoken-clean}

\begin{math-stroke}[Dyadic Cube Definitions]
A dyadic cube of side length $2^{-p}$ in $\mathbb{R}^n$ anchored at $z$ is defined as:
\[
Q_p^n(z) = z + 2^{-p}[0, 1)^n
\]
If the anchor point splits across our axes such that $z = (a,b)$ where $a \in \mathbb{R}^k$ and $b \in \mathbb{R}^{n-k}$, then the $n$-dimensional cube factors across the dimensions:
\[
Q_p^n(z) = Q_p^k(a) \times Q_p^{n-k}(b)
\]
Where, generally for any dimension $m$:
\[
Q_p^m(i) = i + 2^{-p}[0, 1)^m \subset \mathbb{R}^m
\]

\begin{explanation-of-steps}
The geometric idea here is simple despite the rigorous notation. Imagine building a replica of the shape entirely out of LEGO bricks (``dyadic cubes''). If you slice this LEGO shape vertically at position $x$, the cross-section does not change at all as long as your slice stays within the width of that specific column of bricks. It only jumps to a new value when $x$ crosses into the next column. This is why the marginal densities become \textit{step functions}. 

Mathematically, the formula $Q_p^n(z) = Q_p^k(a) \times Q_p^{n-k}(b)$ proves that an $n$-dimensional brick is constructed simply by taking a $k$-dimensional base square on the floor, and extruding it upwards along an $(n-k)$-dimensional vertical line. This factorization is what rigorously justifies isolating the variables to integrate them one dimension at a time.
\end{explanation-of-steps}
\end{math-stroke}

\begin{didactic-insight}[Calculus via Combinatorics]
By replacing the continuous, smooth curve with a grid of literal squares (dyadic cubes), the nasty continuous calculus problem drops away. As long as your $x$-coordinate stays within the width of a single vertical column of squares, the vertical slice does not change its length. It is piecewise constant. A terrifying integration problem becomes the simple integration of a step function!
\end{didactic-insight}
 
%\section{Cavalieri's Principle: Integration of Step Functions and The Squeeze}

\emph{AI Extraction Protocol}

\textit{Segment: 15:15 -- 30:42}


\begin{meta-note}[Segment Continuation]
Continuing directly from the previous segment, the professor has just established that the $n$-dimensional set $A$ can be approximated from the inside and the outside using a discrete grid of dyadic cubes ("pixels"). He now transitions to the algebraic formalization of these geometric bounds.
\end{meta-note}

% ==========================================
% SECTION 5: FORMALIZING THE APPROXIMATIONS
% ==========================================
\section{Formalizing the Approximations}

\begin{spoken-clean}[15:15 - 18:22]
Let us formalize these geometric bounds. Given any arbitrary precision $\varepsilon > 0$, we can choose a sufficiently fine pixel size—let us say, governed by some integer $p \in \mathbb{N}$ determining the side length $2^{-p}$. By the very definition of Jordan measurability, there must exist an inner dyadic set $F$ and an outer dyadic set $G$ such that our set $A$ is completely sandwiched between them.

Furthermore, we are guaranteed that the total volume of the large outer approximation minus the total volume of the small inner approximation is strictly less than $\varepsilon$ (i.e., $\mu_n(G) - \mu_n(F) < \varepsilon$). 

Now, let us examine the vertical cross-sections of these pixelated approximations. For any base coordinate $x$ residing within our bounding box $[-c, c]^k$, we define the slice of the outer approximation, $G_x$, as the set of all vertical $y$-coordinates such that the composite point $(x,y)$ belongs to $G$. We define the slice of the inner approximation, $F_x$, in the exact same manner. Recall that both $F$ and $G$ are simply finite unions of dyadic cubes of pixel size $2^{-p}$.
\end{spoken-clean}

\begin{ai-note}[Establishing Jordan Bounds]
The professor moves to a clean section of the center board to begin the formal algebraic proof. He writes down the standard definition of a Jordan measurable set using the dyadic approximations $F$ and $G$, followed by the set-builder notation for their respective vertical slices.
\end{ai-note}

\begin{math-stroke}[Inner and Outer Dyadic Sets and Slices]
Given $\varepsilon > 0$, $\exists p \in \mathbb{N}$ and dyadic sets $F, G$ (composed of cubes of side length $2^{-p}$) such that:
\[
F \subset A \subset G \quad \text{and} \quad \mu_n(G) - \mu_n(F) < \varepsilon
\]
For all $x \in [-c, c]^k$, we define the vertical slices:
\begin{align*}
G_x &= \{ y \in \mathbb{R}^{n-k} \mid (x,y) \in G \} \\
F_x &= \{ y \in \mathbb{R}^{n-k} \mid (x,y) \in F \}
\end{align*}
\end{math-stroke}

% ==========================================
% SECTION 6: PIXELS TO STEP FUNCTIONS
% ==========================================
\section{Pixels to Step Functions}

\begin{spoken-clean}[18:22 - 21:50]
Let us define a new function, $f(x)$, which represents the $(n-k)$-dimensional volume of the inner slice $F_x$. Similarly, let $g(x)$ represent the volume of the outer slice $G_x$. 

Because these sets are constructed entirely from a uniform grid of dyadic cubes, the geometry of the vertical slice does not change continuously as you move $x$. As long as your coordinate $x$ remains within the boundaries of a single base pixel $Q_p^k(a)$, the volume of the vertical slice intersecting it remains perfectly constant. 

What does it mean for a function to be piecewise constant across a grid of cubes? It means that $f(x)$ and $g(x)$ are, by definition, dyadic step functions. Additionally, since our entire set is contained within the bounding box, these functions trivially evaluate to zero outside of $[-c, c]^k$. Consequently, we have successfully bounded our continuous problem with adequate step functions, which are exactly the tools we need to approximate the Riemann integral.
\end{spoken-clean}

\begin{didactic-insight}[The Geometric Leap]
This explicitly bridges the physical drawing of "pixels" to the algebraic definition of a step function. Because $F$ is constructed purely of vertical columns of uniform width $2^{-p}$, moving $x$ horizontally within the boundary of a single column does not change the vertical cross-section $F_x$. Thus, the volume function $f(x)$ is perfectly flat (constant) across that entire base pixel.
\end{didactic-insight}

\begin{math-stroke}[Step Function Property]
Define the slice volume functions $f, g: [-c, c]^k \to \mathbb{R}$ as:
\begin{align*}
f(x) &:= \mu_{n-k}(F_x) \\
g(x) &:= \mu_{n-k}(G_x)
\end{align*}
Because $F$ and $G$ are finite unions of dyadic cubes, $f(x)$ and $g(x)$ are constant on every base cube $Q_p^k(a)$. Furthermore, $f(x) = g(x) = 0$ outside of $[-c, c]^k$. \\
$\implies f$ and $g$ are \textbf{dyadic step functions}.
\end{math-stroke}

% ==========================================
% SECTION 7: THE COMBINATORIAL COUNTING ARGUMENT
% ==========================================
\section{The Combinatorial Counting Argument}

\begin{spoken-clean}[21:50 - 25:30]
Now, let us integrate these step functions. What is the integral of $f(x) \, dx$ over our base space? 

Because we are operating in the pixelated world of dyadic cubes, integrating the volume of the slices over the base domain degenerates into a simple counting argument! You are merely looking at each column and asking, ``How many pixels are stacked here?'' When you integrate, you are simply summing those discrete columns together to recover the total number of pixels in the entire shape. 

There are no limits here, and no continuous approximation errors. It is pure combinatorics. Therefore, the integral of our inner step function $f(x)$ evaluates exactly to the total $n$-dimensional volume of the inner approximation, $\mu_n(F)$. By identical reasoning, the integral of $g(x)$ evaluates exactly to $\mu_n(G)$.
\end{spoken-clean}

\begin{ai-note}[Integrating the Step Functions]
He writes the integral of $f(x)$ and $g(x)$ across the base space. He explicitly equates the integral over the base space to the total volumetric measure of the multi-dimensional approximation sets.
\end{ai-note}

\begin{math-stroke}[Combinatorial Integration]
Because $f$ and $g$ are dyadic step functions, their Riemann integrals are precisely the sum of the volumes of their constituent $n$-dimensional cubes:
\begin{align*}
\int_{[-c,c]^k} f(x) \, dx &= \mu_n(F) \\
\int_{[-c,c]^k} g(x) \, dx &= \mu_n(G)
\end{align*}
\end{math-stroke}

\begin{didactic-insight}[Calculus Becomes Counting]
This is a beautiful pedagogical moment. Multi-dimensional integration is demystified by showing that when dealing with step functions constructed from a uniform grid, the Riemann integral becomes simple addition. Taking the volume of a 2D slice and multiplying it by the tiny 1D width of the base pixel ($dx$) yields the 3D volume of that column. Summing them all up gives the total volume.
\end{didactic-insight}

% ==========================================
% SECTION 8: THE CHAIN OF INEQUALITIES
% ==========================================
\section{The Chain of Inequalities}

\begin{spoken-clean}[25:30 - 27:30]
We now rely on our fundamental geometric inclusion. Because the inner approximation $F$ is contained within $A$, which is in turn contained within $G$, this strict nesting applies to every single vertical slice. Thus, $F_x \subset A_x \subset G_x$. 

This geometric nesting forces a strict ordering on the measures of these slices. The measure of the inner slice, $f(x)$, must be less than or equal to the inner measure of $A_x$, which is our lower marginal $\underline{\varphi}(x)$ (i.e., $f(x) \le \underline{\varphi}(x)$). In turn, the outer measure of $A_x$, which is our upper marginal $\overline{\varphi}(x)$, must be less than or equal to the measure of the outer slice, $g(x)$ (i.e., $\overline{\varphi}(x) \le g(x)$). 

We have successfully constructed a chain of inequalities for our functions. A foundational property of Riemann integration is that if functions are strictly ordered, their integrals maintain that exact order. We will apply the lower Riemann integral to the lower marginal, and the upper Riemann integral to the upper marginal, preserving the entire inequality chain from end to end.
\end{spoken-clean}

\begin{ai-note}[Establishing the Squeeze]
The professor moves back to the center board. He uses the geometric inclusion to establish an ordered chain of marginal densities. He then applies the lower and upper Riemann integrals to this chain, drawing a long series of inequalities.
\end{ai-note}

\begin{math-stroke}[Squeezing the Marginals]
From the inclusion $F \subset A \subset G$, we have for all slices $x$:
\[
F_x \subset A_x \subset G_x
\]
Taking the respective measures preserves the ordering:
\[
\underbrace{\mu_{n-k}(F_x)}_{f(x)} \le \underbrace{\mu_{n-k,\text{in}}(A_x)}_{\underline{\varphi}(x)} \le \underbrace{\mu_{n-k,\text{out}}(A_x)}_{\overline{\varphi}(x)} \le \underbrace{\mu_{n-k}(G_x)}_{g(x)}
\]
Integrating across $[-c, c]^k$ preserves the inequality chain:
\[
\int f(x) \, dx \le \underline{\int} \underline{\varphi}(x) \, dx \le \overline{\int} \overline{\varphi}(x) \, dx \le \int g(x) \, dx
\]
\end{math-stroke}

% ==========================================
% SECTION 9: CONCLUSION OF CAVALIERI'S PRINCIPLE
% ==========================================
\section{Conclusion of the Proof}

\begin{spoken-clean}[27:30 - 30:42]
Let us substitute the combinatorial volumes we found earlier into the bounding ends of this inequality. We know that the integral of $f(x)$ is exactly $\mu_n(F)$, and the integral of $g(x)$ is exactly $\mu_n(G)$. 

Therefore, our unknown marginal integrals are permanently trapped between the volume of the inner approximation and the volume of the outer approximation. But remember our initial premise! We constructed our pixels specifically such that the difference between the large outer volume and the small inner volume is strictly less than $\varepsilon$. 

What happens as we send $\varepsilon$ to zero? The outer and inner bounds squeeze tightly together. Because the marginal integrals are trapped in the middle, and because they are entirely independent of our arbitrary choice of $\varepsilon$, they are forced to equal the exact same value. 

This simultaneously proves two things: first, because the lower and upper integrals coincide, the marginal functions are definitively Riemann integrable. Second, because they are squeezed between bounds that converge to the true volume of the shape, their integral evaluates exactly to $\mu_n(A)$. And that completes the proof of Cavalieri's Principle.
\end{spoken-clean}

\begin{nice-box}[Substituting and Concluding]
He substitutes $\mu_n(F)$ and $\mu_n(G)$ into the bounding ends of the integral inequality chain. He circles the $\varepsilon$ condition, shows the bounds collapsing, and writes the final conclusion, boxing the result.
\end{nice-box}

\begin{math-stroke}[The Final Inequality Chain]
Substituting the known combinatorial integrals:
\[
\mu_n(F) \le \underline{\int} \underline{\varphi}(x) \, dx \le \overline{\int} \overline{\varphi}(x) \, dx \le \mu_n(G)
\]
We know by our initial choice of dyadic approximations that:
\[
\mu_n(G) - \mu_n(F) < \varepsilon
\]
\end{math-stroke}

\begin{orange-formula}
\begin{theorem}[Conclusion of Cavalieri's Principle]
Since the upper and lower integrals of the marginal functions are squeezed between $\mu_n(F)$ and $\mu_n(G)$ for any $\varepsilon > 0$, and the length of that interval vanishes as $\varepsilon \to 0$, we conclude:
\[
\int_{[-c,c]^k} \underline{\varphi}(x) \, dx = \int_{[-c,c]^k} \overline{\varphi}(x) \, dx = \mu_n(A)
\]
Therefore, $\underline{\varphi}$ and $\overline{\varphi}$ are Riemann integrable, completing the proof. $\blacksquare$
\end{theorem}
\end{orange-formula}
 
%
% ==========================================
% LECTURE METADATA
% ==========================================
\section{From Cavalieri to Fubini: Integrating Functions via Hypographs}

\emph{AI Extraction Protocol}

\textit{Segment: 30:42 -- 45:00 (End of Lecture)}

\begin{meta-note}[Segment Transition]
The professor has just finished the rigorous dyadic-cube proof of Cavalieri's Principle for $n$-dimensional sets. He erases the center and right chalkboards to transition to the ultimate goal of the lecture: applying this geometric slicing principle to the calculus of multi-variable functions. 
\end{meta-note}

% ==========================================
% SECTION 10: THE TRANSITION TO FUNCTIONS
% ==========================================
\section{From Sets to Functions: The Hypograph}

\begin{spoken-clean}[30:42 - 34:15]
We have successfully and rigorously established Cavalieri's Principle for the volumes of multi-dimensional sets. But we are taking a class on Analysis, not just geometry, so we must now ask: how do we extend this logic to the integration of functions? 

You will recall that much earlier in the semester, when we first defined the Riemann integral of a positive function $f$ over a domain $A$, we defined it strictly in terms of the measure of its hypograph. 

What is the hypograph geometrically? It is the $(n+1)$-dimensional solid region situated directly beneath the graph of the function and directly above the base domain $A$. If we wish to integrate a function $f(x,y)$ over a product space $\mathbb{R}^k \times \mathbb{R}^{n-k}$, we are effectively seeking the $(n+1)$-dimensional volume of this overarching solid block. 

Therefore, integrating a function is fundamentally identical to measuring a physical shape. We can simply apply Cavalieri's Principle directly to this $(n+1)$-dimensional hypograph. By taking vertical slices of this solid, we reduce a massive multi-dimensional volume calculation into an iterated sequence of much simpler lower-dimensional integrals. This elegant corollary is what we call Fubini's Theorem.
\end{spoken-clean}

\begin{ai-note}[Recalling the Hypograph Connection]
The professor moves to the freshly cleaned center board. He writes down the formal definition of an integral via the hypograph, establishing the crucial bridge between the pure geometry of Cavalieri and the operational calculus of Fubini.
\end{ai-note}

\begin{math-stroke}[Integral as a Measure of the Hypograph]
Let $f: A \to \mathbb{R}_{\ge 0}$ be a Riemann integrable function over a Jordan measurable set $A \subset \mathbb{R}^n$.
The hypograph of $f$ is defined as the $(n+1)$-dimensional set:
\[
\text{Hypo}(f) = \{ (x, z) \in \mathbb{R}^n \times \mathbb{R} \mid x \in A, \ 0 \le z \le f(x) \}
\]
By our foundational definition, the integral of the function is identically the measure of this set:
\[
\int_A f(x) \, dx = \mu_{n+1}(\text{Hypo}(f))
\]
\end{math-stroke}

\begin{didactic-insight}[The Geometric Unification]
This is the grand unification of the lecture. By framing the integral not as an abstract analytic operation, but as a literal physical volume (the hypograph), the "potato slicing" intuition mastered with Cavalieri's Principle can be seamlessly applied directly to abstract function integration.
\end{didactic-insight}

% ==========================================
% SECTION 11: STATEMENT OF FUBINI'S THEOREM
% ==========================================
\section{Statement of Fubini's Theorem}

\begin{spoken-clean}[34:15 - 38:00]
Let us formalize this into a proper theorem. Suppose we partition our base domain space $\mathbb{R}^n$ into a Cartesian product $\mathbb{R}^k \times \mathbb{R}^{n-k}$, utilizing coordinate vectors $x$ and $y$ respectively. Let $A$ be our domain of integration, and $f(x,y)$ be our continuous, integrable function hovering above it.

If we freeze a base coordinate $x \in \mathbb{R}^k$, treating it strictly as a constant, we can restrict our function to the corresponding vertical slice of the domain, $A_x$. We shall denote this restricted, lower-dimensional function as $f_x(y)$. 

Under Fubini's Theorem, assuming the requisite integrability conditions hold over the individual slices, the total $n$-dimensional integral of $f$ over the entire domain $A$ is exactly equivalent to the iterated integral! 

Operationally, what does this mean? You first integrate the restricted function $f_x(y)$ with respect to $y$ over the slice $A_x$. Because $x$ was held constant, the result of this inner integral is a new marginal function purely in terms of $x$. Subsequently, you integrate this resulting marginal function with respect to $x$ over the base projection of the domain.
\end{spoken-clean}

\begin{ai-note}[Formalizing the Sliced Function]
The professor meticulously writes out the notation for the restricted function $f_x(y)$. He taps the board with his chalk under the $x$ variable to emphasize to the students that $x$ must be treated purely as a constant scalar during the inner integration step. He then writes the full statement of Fubini's Theorem in bright orange chalk.
\end{ai-note}

\begin{math-stroke}[Function Restriction and Iteration]
Given a function $f: A \to \mathbb{R}$ where the domain $A \subset \mathbb{R}^k \times \mathbb{R}^{n-k}$. \\
For a fixed, constant $x \in \mathbb{R}^k$, define the slice domain $A_x = \{y \mid (x,y) \in A\}$. \\
We define the restricted slice function $f_x: A_x \to \mathbb{R}$ as:
\[
f_x(y) = f(x,y) \quad \text{(treating } x \text{ strictly as a constant)}
\]
\end{math-stroke}

\begin{orange-formula}
\begin{theorem}[Fubini's Theorem]
Let $A = X \times Y$ be a compact, Jordan measurable domain, and $f: A \to \mathbb{R}$ be integrable. The multiple integral can be evaluated sequentially as an iterated integral:
\[
\int_A f(x,y) \, d(x,y) = \int_X \left( \int_{A_x} f(x,y) \, dy \right) dx = \int_Y \left( \int_{A_y} f(x,y) \, dx \right) dy
\]
\end{theorem}
\end{orange-formula}

% ==========================================
% SECTION 12: VISUALIZING THE ITERATED INTEGRAL
% ==========================================
\section{Visualizing the Iterated Integral}

\begin{spoken-clean}[38:00 - 42:30]
To truly grasp the mechanics of Fubini's Theorem and prevent algebraic mistakes, one must visualize the sequential nature of this volume accumulation. 

When we evaluate the inner integral — integrating $f(x,y)$ strictly with respect to $y$ while holding $x$ frozen at some value $x_0$ — what are we calculating geometrically? We are calculating the exact 2D cross-sectional area of the solid hypograph at that specific coordinate $x_0$. 

This inner integral effectively collapses the massive $(n+1)$-dimensional solid into a one-dimensional density function residing on the base axis. The outer integral then sweeps this cross-sectional area across the entire length of the domain $X$, accumulating those differential slices to arrive at the total total measure of the solid. 

The symmetry of the theorem ensures that the order of accumulation—whether sweeping left-to-right along the $x$-axis or front-to-back along the $y$-axis—is irrelevant to the final numerical volume, provided your integration boundaries $A_x$ and $A_y$ are properly accounted for.
\end{spoken-clean}

\begin{ai-note}[The Hypograph Cross-Section]
The professor moves to the right board to draw a 3D perspective of a curved surface floating above a rectangular domain in the $xy$-plane. He uses red chalk to shade a specific 2D planar cross-section beneath the surface, labeling its area as the exact result of the inner integral, $\int f(x_0, y) dy$.
\end{ai-note}

\begin{math-stroke}[Geometric Visualization of Fubini's Theorem]
\begin{center}
\begin{tikzpicture}[scale=1.5]
    % Axes (Muted to push them to the background, keeping focus on the surfaces)
    \draw[->, thick, gray!80!black] (0,0,0) -- (4,0,0) node[right] {$y$ axis};
    \draw[->, thick, gray!80!black] (0,0,0) -- (0,3,0) node[above] {$z = f(x,y)$};
    \draw[->, thick, gray!80!black] (0,0,0) -- (0,0,4) node[below left] {$x$ axis};

    % Domain in xy-plane
    \draw[dashed, profblue, thick] (1,0,1) -- (3,0,1) -- (3,0,3) -- (1,0,3) -- cycle;
    \node[profblue] at (2,-0.3,3.5) {Base Domain $A = X \times Y$};

    % Surface
    \draw[thick, fill=proforange!20, opacity=0.8] (1,2,1) to[out=20,in=160] (3,2.5,1) to[out=-20,in=70] (3,1.5,3) to[out=160,in=-20] (1,1,3) to[out=70,in=-20] cycle;
    \node[proforange!80!black] at (2,2.8,1) {Surface $z = f(x,y)$};

    % The Slice at constant x_0
    \draw[thick, profred, fill=profred!20, opacity=0.7] (1,0,2) -- (3,0,2) -- (3,2.0,2) to[out=150,in=-10] (1,1.5,2) -- cycle;
    
    % Slice Annotations
    \draw[dashed, profred, thick] (0,0,2) -- (1,0,2);
    \node[left, profred] at (0,0,2) {$x_0$};
    \node[profred, fill=white, inner sep=1pt] at (2, 0.8, 2) {Area $= \int f(x_0, y) \, dy$};

\end{tikzpicture}
\end{center}

\begin{explanation-of-steps}
The visual clarifies a crucial geometric concept: why the bounds of integration change when swapping $dx$ and $dy$. The red slice represents the inner integral: evaluating the area of the 2D plane at a constant $x_0$. As $x_0$ moves along the $x$-axis, the boundaries of that slice might grow or shrink if the base domain $A$ is not a perfect rectangle (e.g., a triangle or a circle). The outer integral then integrates this continuously changing 1D area function along the $x$-axis to get the final 3D volume.
\end{explanation-of-steps}
\end{math-stroke}

\begin{spoken-clean}[42:30 - 45:00]
So, Fubini's Theorem is not magic; it is simply Cavalieri's Principle applied to the volume under a graph. By treating one variable as a constant, you reduce a multivariable problem into a sequence of single-variable calculus problems that you already know how to solve from Analysis I. 

We will spend the next lecture doing heavy computational examples of this, swapping the bounds of integration, and looking at how this interacts with our Change of Variables theorem from last week. Please review your single-variable integration techniques tonight, because you will be doing them twice as much now. Have a good afternoon, everyone.
\end{spoken-clean}


%\section{Integration, Volumes, and Fubini's Theorem}

% Dont change this Gemini: 
\begin{nice-box}[Board State \& Physical Actions: Polar Coordinates Substitution in $\mathbb{R}^3$]
The mapping $\Psi: D \to \mathbb{R}^3$ for spherical (polar) coordinates is defined as:
\[
\Psi(y_1, y_2, y_3) = 
\begin{pmatrix} 
y_1 \cos y_2 \cos y_3 \\ 
y_1 \sin y_2 \cos y_3 \\ 
y_1 \sin y_3 
\end{pmatrix}
\]
Where the domain $D$ is the open box:
\[
D = (0, 1) \times (0, 2\pi) \times \left(-\frac{\pi}{2}, \frac{\pi}{2}\right)
\]
The volume transformation leverages the Jacobian determinant:
\[
\underbrace{dx}_{\text{Original Volume Element}} = \underbrace{|\det J\Psi(y)| \, dy}_{\text{Scaled Volume Element}} = y_1^2 \cos y_3 \, dy
\]

\begin{orangeformula}
\begin{theorem}[Fubini's Theorem Formulation]
By decomposing the space as $\mathbb{R}^n \equiv \mathbb{R}^k \times \mathbb{R}^{n-k}$, we can express the integral of $f$ over a bounded box $[-C,C]^n$ via its marginals:
\[
\int_{[-C,C]^n} f(x,y) \, d(x,y) = \int_{[-C,C]^k} \overline{\varphi}(x) \, dx = \int_{[-C,C]^k} \underline{\varphi}(x) \, dx
\]
where the upper and lower marginals are:
\[
\underline{\varphi}(x) = \underline{\int}_{[-C,C]^{n-k}} f(x,y) \, dy \quad \text{and} \quad \overline{\varphi}(x) = \overline{\int}_{[-C,C]^{n-k}} f(x,y) \, dy
\]
\end{theorem}
\end{orangeformula}
\end{nice-box}

\begin{spoken_clean}[00:00:03 - 00:02:08]
The plan for today is to finish our discussion of integration and provide more details on the proof of the Change of Variables formula. We will return to the formal proof later, but for now, I want to finish these parts, especially the examples, so that you can begin computing them. We will see many examples: some with me, some in the problem sessions, and some by yourselves. Once we do this, all this abstract theory will land on solid ground, and you will be able to understand it much better.

As motivation from Monday's lecture, we were trying to compute the volume of a ball in $\mathbb{R}^3$. We computed the volume of the disk in $\mathbb{R}^2$ with a trick, but when we try to do the same in $\mathbb{R}^3$, we must apply the Change of Variables formula using polar coordinates. Look at the board: these are the polar coordinates. We established that in polar coordinates, the volume element $dx$ is replaced by the new volume element $dy$, scaled by the determinant of the Jacobian. If you compute the Jacobian matrix and then its determinant, you will see exactly this value.
\end{spoken_clean}

\begin{spoken_clean}[00:02:08 - 00:04:18]
So, we need to compute this integral. I mentioned last time that the image of our polar coordinate domain $D$ is not the completely full ball, but it equals the full ball up to a measure-zero set (i.e., the boundary planes that are excluded from the open domain $D$ have no volume). Therefore, the volume of this domain equals the volume of the ball itself. We can compute this volume using the Change of Variables formula as an integral. But here we got stuck, because we did not know how to evaluate this multivariable integral.

Then we introduced the Cavalieri principle. This is a classic example of ``divide and conquer'' in mathematics. If you do not know how to compute an integral over a full volume, you compute it over a slice, thereby reducing the problem by one dimension to yield an easier integral. We stated the theorem for volumes of bodies, and then we stated the same theorem for integrals of functions. That was a bit complicated, but now, we will establish a corollary that is much easier to apply in practice.
\end{spoken_clean}

\begin{didactic_insight}[Divide and Conquer Algorithm]
The Cavalieri principle is beautifully framed not just as a geometric trick, but as an algorithmic strategy. A seemingly impossible 3D volume integral is broken down (``divide'') into a continuous sequence of 2D slices (``conquer''), systematically reducing the dimensional complexity step-by-step.
\end{didactic_insight}

\begin{spoken_clean}[00:04:18 - 00:06:07]
Let me remind you of our setting. We are splitting the space $\mathbb{R}^n$ as a Cartesian product of $\mathbb{R}^k$ and $\mathbb{R}^{n-k}$. The $x$ variable lives in $\mathbb{R}^k$, and the $y$ variable lives in $\mathbb{R}^{n-k}$. Now, suppose we have a function that is Riemann integrable. We can assume that every Riemann integrable function is zero outside of a bounded set (since we are integrating over bounded domains), right? Let's say that the large box where the function is non-zero is $[-C, C]^n$. 

You can define these marginals—the lower and upper marginals. This means you fix one $x$, and then you look at the vertical slice on top of that $x$. You now have a function of $y$. Because the $x$ is fixed, you can integrate this new function with respect to $y$. Sometimes, when you fix a specific $x$, you might be unlucky (e.g., the slice might intersect a pathological boundary), and the resulting function of $y$ is no longer strictly Riemann integrable. However, what you can always do is define the lower and upper integrals. These are universally well-defined. Then, you plug them back in: you integrate this lower or upper marginal with respect to $x$. You do not need to care which one you choose, as both will yield the exact integral of $f$.
\end{spoken_clean}

\begin{math_stroke}[Defining the Marginals]
    
\begin{equation}
    \begin{aligned}
        &\int_{[-C,C]^k} \underbrace{\underline{\varphi}(x)}_{\text{``Define the lower integral... always well-defined''}} \, dx \\
        &= \int_{[-C,C]^k} \underbrace{\overline{\varphi}(x)}_{\text{``or the upper marginal... both yield the exact integral''}} \, dx \\
        &= \int f    
    \end{aligned}
\end{equation}
\end{math_stroke}

\begin{spoken_clean}[00:06:07 - 00:08:02]
This is the general theorem. But now, we will look at a practical corollary so that we do not have to deal with so many bars on the top and bottom of our integral signs. We will consider a very simple situation. Let $K$ be a closed box, which we can express as a product of intervals forming a subset of $\mathbb{R}^n$. Furthermore, we are given a continuous function $f$ mapping from $K$ to $\mathbb{R}$. 

We know, because we proved earlier that uniformly continuous functions over a Jordan measurable set are Riemann integrable, that this will work out perfectly (i.e., we do not need the upper and lower integrals because continuity guarantees integrability). This box is Jordan measurable, and the function is continuous over a compact set. Therefore, it is uniformly continuous, and thus, it is automatically Riemann integrable. I no longer need to be concerned about verifying this assumption. I simply want to state the following: the integral over the box of $f(x) \, dx$ can be computed iteratively. You first integrate with respect to $x_n$.
\end{spoken_clean}

\begin{orangeformula}
\begin{corollary}[The Practical Fubini Iteration]
Let $K = [a_1, b_1] \times [a_2, b_2] \times \dots \times [a_n, b_n]$ be a closed box in $\mathbb{R}^n$. For any continuous function $f: K \to \mathbb{R}$, the integral evaluates via nested iteration:
\[
\int_K f(x) \, dx = \int_{a_1}^{b_1} \dots \underbrace{ \left( \int_{a_n}^{b_n} f(x_1, \dots, x_n) \, dx_n \right) }_{\text{``First, I integrate with respect to the } n\text{-th variable''}} \dots \, dx_1
\]
\end{corollary}
\end{orangeformula}

\begin{spoken_clean}[00:08:02 - 00:10:22]
For now, we will put parentheses around the integrals, but later on, I will erase them to save time. The iteration continues with the integral over the second variable, all the way up to the last variable, $x_n$. What I mean is this: I take my function, which depends on $n$ variables, and I first integrate it with respect to the $n$-th variable. The other variables are held constant; they are essentially frozen. 

Once I integrate with respect to this variable (i.e., after evaluating the definite integral from $a_n$ to $b_n$), the result only depends on the remaining variables, from $x_1$ up to $x_{n-1}$. Then, I integrate again, this time with respect to the next variable, and I continue this process until I get a final number. Let me apply it to a concrete example, so that you can see what this implies in reality. You first evaluate the interval from $a_1$ to $b_1$, then $a_2$ to $b_2$, and finally $a_3$ to $b_3$. The last integral you write down is actually the first one you compute. It acts like a parenthesis. An integral is naturally closed by its corresponding differential (e.g., $dx_3$). That is exactly why you do not need the parentheses; the $dx_3$ implicitly closes the innermost integral block, $dx_2$ closes the next one, and so on. 
\end{spoken_clean}

\begin{spoken_clean}[00:10:22 - 00:12:46]
Let's look at an example. What is the integral of $y_1^2 \cos(y_3)$? Since the boundary has measure zero, the boundary of the set forms a box, and I can apply our theorem directly. Whether we integrate over the open set or the closed set is entirely the same (since the boundary is a set of measure zero and does not contribute to the integral's value), so please do not complain that our domain here is open! Notice I am also just changing the variable names. Instead of calling them $x$, I will call them $y$, but the math is exactly the same. The intervals are $y_1$ in $[0, 1]$, $y_2$ in $[0, 2\pi]$, and $y_3$ in $[-\pi/2, \pi/2]$. Now, we are ready to start computing.
\end{spoken_clean}

\begin{math_stroke}[Setting up the Iterated Integral]
Mapping the polar coordinate domain $D$ into the iterated framework:
\[ 
\text{Volume}(B_3) = \underbrace{\int_{0}^{1}}_{\text{``} y_1 \text{ in } [0, 1] \text{''}} \underbrace{\int_{0}^{2\pi}}_{\text{``} y_2 \text{ in } [0, 2\pi] \text{''}} \underbrace{\int_{-\pi/2}^{\pi/2} y_1^2 \cos(y_3) \, \underbrace{dy_3}_{\text{``implicitly closes the inner block''}}}_{\text{``The last integral you write down is the first one you compute''}} \, dy_2 \, dy_1 
\]
\end{math_stroke}

\begin{spoken_clean}[00:12:46 - 00:14:05]
First, I need to compute the inner integral, from $-\pi/2$ to $\pi/2$, of $y_1^2 \cos(y_3) \, dy_3$. Just as we do with partial derivatives, when you integrate with respect to $y_3$, you must remember that $y_1$ is fixed. So, you can think of it as a constant for now, and pull it completely out of the integral. This leaves us with $y_1^2$ multiplied by the integral of the cosine function. The primitive of cosine is sine, and we evaluate it between $-\pi/2$ and $\pi/2$ (i.e., $\sin(\pi/2) - \sin(-\pi/2) = 1 - (-1) = 2$).
\end{spoken_clean}

\begin{math_stroke}[Step 1: The Inner Integral]
\[
\int_{-\pi/2}^{\pi/2} y_1^2 \cos(y_3) \, dy_3 = \underbrace{y_1^2}_{\text{``think of it as a constant... pull it out''}} \cdot \underbrace{\Big[ \sin(y_3) \Big]_{-\pi/2}^{\pi/2}}_{\text{``i.e., } \sin(\pi/2) - \sin(-\pi/2) = 2\text{''}} = 2y_1^2 
\]
\end{math_stroke}

\begin{spoken_clean}[00:14:05 - 00:15:00]
Once we have the result of this first integral, we insert it right back into the next layer. So, now, we need to compute the integral between $0$ and $2\pi$ of $2y_1^2 \, dy_2$. This is trivial because our integrand is now a constant with respect to $y_2$. It factors out of the integral, and you simply multiply it by the length of the interval, $2\pi$, giving you $4\pi y_1^2$. 

Finally, we reach the third step. We need to compute the integral from $0$ to $1$ of this remaining expression. The primitive of $y_1^2$ is $y_1^3 / 3$. Evaluating this between $y_1 = 0$ and $y_1 = 1$, I get exactly $\frac{4}{3}\pi$. That is the volume of the unit ball! Yes, we already knew this, and we remember being told this, right? But the earth is not flat, and some people still doubt it. Therefore, we must prove the things we were told in high school for ourselves to ensure the math is sound. And as you can see, we have successfully proven it.
\end{spoken_clean}

\begin{math_stroke}[Step 2 and 3: Resolving the Final Volume]
\textbf{Step 2 (Middle Integral):}
\[ 
\int_{0}^{2\pi} \underbrace{2y_1^2}_{\text{``trivial, because our integrand is now a constant''}} \, dy_2 = 2y_1^2 \Big[ y_2 \Big]_{0}^{2\pi} = 4\pi y_1^2 
\]

\textbf{Step 3 (Outer Integral):}
\[ 
\int_{0}^{1} 4\pi y_1^2 \, dy_1 = 4\pi \underbrace{\left[ \frac{y_1^3}{3} \right]_0^1}_{\text{``The primitive of } y_1^2 \text{ is } y_1^3 / 3 \text{''}} = 4\pi \left( \frac{1}{3} - 0 \right) = \frac{4\pi}{3} 
\]
\end{math_stroke}
%\section{Alternative Methods and the Proof of Fubini in 3D}

\begin{nice-box}[Cartesian Slicing vs. 3D Box Splitting]
The left board contains a diagram of a 2D cross-section of a sphere sliced at height $x_3$. The right board sets up the formal 3D box definitions for the proof of Fubini's theorem.
\end{nice-box}

\begin{math_stroke}[Cartesian Slicing and Box Definition]
\textbf{Left Board: Cartesian Slicing of the Sphere} \\
Using the Pythagorean theorem, the radius of the slice is written as:
\[
r^2 + x_3^2 = 1 \implies r = \sqrt{1 - x_3^2}
\]
The Cavalieri volume integral is established:
\[
V = \int_{-1}^{1} \pi (1 - x_3^2) \, dx_3
\]

\textbf{Right Board: 3D Box Definition for Fubini's Proof}
\[
K = [a_1, b_1] \times [a_2, b_2] \times [a_3, b_3] \subset \mathbb{R}^3
\]
The splitting strategy ($k=1$):
\[
\mathbb{R}^3 \equiv \underbrace{\mathbb{R}^1}_{x_1} \times \underbrace{\mathbb{R}^2}_{(x_2, x_3)}
\]
\end{math_stroke}
\begin{spoken_clean}[15:00 - 18:38]
So, is it clear how to apply the theorem in practice? Is it useful? Yes, absolutely. Now, we will move on to prove the corollary we just used. We understand what this corollary says, and the proof itself is not particularly difficult, although the formal statement can be a bit intimidating to read at first. Because this is the kind of exercise that everyone needs to know how to do, we must go through it. Are there any questions before we begin the proof?
\end{spoken_clean}

\begin{student_question}
A student raises their hand and asks: "Would it not be easier to calculate the volume of the unit sphere just using the slicing lemma directly, without doing a change of variables? Just slicing the sphere into circles?"
\end{student_question}

\begin{spoken_clean}[18:38 - 21:05]
You are asking if we could calculate the volume of the unit sphere without doing a change of variables, simply by slicing it into circles. It is true, you can certainly slice the sphere like that. You can do it, but it is not necessarily easier! Let's look at what that entails. You can slice it horizontally, along the $x_3$-axis, which we can call the height. 

The Cavalieri principle says that you can integrate from $-1$ to $1$, and in each slice, you compute the area of the resulting disk. By using Pythagoras, if the hypotenuse is $1$, and the vertical distance is $x_3$, the radius squared of this disk is exactly $1 - x_3^2$. So, you can also evaluate this integral to find the volume. However, notice that you must remember how to compute the primitive of this specific polynomial function. You always need to do some extra work here. In contrast, with the polar coordinates we just used, the primitives were completely trivial.
\end{spoken_clean}

\begin{math_stroke}[Direct Cartesian Slicing]
Setting up the integral based on the student's suggestion:
\[
\text{Volume} = \underbrace{\int_{-1}^{1}}_{\text{``integrate from } -1 \text{ to } 1 \text{''}} \underbrace{\pi \big(1 - x_3^2\big)}_{\text{``the radius squared of this disk is exactly } 1 - x_3^2 \text{''}} \, \underbrace{dx_3}_{\text{``along the } x_3 \text{-axis''}}
\]
\end{math_stroke}

\begin{didactic_insight}[The Hidden Cost of Cartesian Integration]
Acknowledging geometric intuition highlights a vital pragmatic truth: Cartesian slicing often leads to integrating roots and complex polynomials. Polar coordinates, while requiring the initial setup of a Jacobian determinant, radically simplify the actual calculus, reducing the integrands to basic trigonometric functions and constants.
\end{didactic_insight}

\begin{spoken_clean}[21:05 - 24:15]
Okay, let us move on to the proof of the corollary. Because this is Analysis II, there is often a big leap from Analysis I—where you can visualize everything very well—to abstract $n$-dimensional spaces, where some people struggle. That is completely normal! 

For concreteness, I will do this proof in three dimensions. The proof in your notes is in $n$ dimensions, and it is exactly the same, but with a general $n$. This is an exercise I am telling you to do all the time: whenever a proof is too abstract in $n$ dimensions, just write it down in 2 or 3 dimensions, and you will see exactly what is going on. Simplify. How do we proceed in 3D? We will use Fubini's theorem with $k = 1$. We have our box, $K$, and we split the space such that the first variable is in one dimension, and the other two are in a two-dimensional box.
\end{spoken_clean}

\begin{math_stroke}[First Dimensional Split ($k=1$)]
Applying the general Fubini theorem to the 3D box $K$:
\[
\int_K f = \underbrace{\int_{a_1}^{b_1}}_{\text{``first variable is in one dimension''}} \underbrace{\Phi_1(x_1)}_{\text{``marginal function''}} \, dx_1
\]
\end{math_stroke}

\begin{spoken_clean}[24:15 - 26:10]
The integral of our function over the box is equal to the integral over $[a_1, b_1]$ of the marginal function, which we will call $\Phi_1(x_1)$. Now, what exactly is this marginal? It is the integral over the remaining two-dimensional box, $[a_2, b_2] \times [a_3, b_3]$, of our function with respect to $x_2$ and $x_3$. 

Because $f$ is continuous over the entire box, if we fix $x_1$, the restriction of $f$ to this 2D plane is also a continuous function. Since it is a continuous function over a closed, two-dimensional box, it is automatically integrable. We do not need to worry about upper and lower integrals here; the integral is perfectly well-defined.
\end{spoken_clean}

\begin{math_stroke}[Expanding the Marginal $\Phi_1$]
\[
\Phi_1(x_1) = \underbrace{\int_{[a_2, b_2] \times [a_3, b_3]}}_{\text{``remaining two-dimensional box''}} \underbrace{f(x_1, x_2, x_3)}_{\text{``if we fix } x_1 \text{, the restriction is continuous''}} \, \underbrace{d(x_2, x_3)}_{\text{``with respect to } x_2 \text{ and } x_3 \text{''}}
\]
\end{math_stroke}

\begin{spoken_clean}[26:10 - 28:00]
Now, we need to compute this inner, two-dimensional integral. But we can simply apply the exact same idea again! To compute this integral, we use Fubini's theorem once more with $k = 1$, but this time, in dimension 2. We reduce the problem by one more dimension. 

We integrate over $[a_2, b_2]$ with respect to $x_2$, and the new marginal inside is the integral over $[a_3, b_3]$ of $f$ with respect to $x_3$. Here, both $x_1$ and $x_2$ are fixed. We integrate with respect to $x_3$, we get a number, and then we move $x_2$. 
\end{spoken_clean}

\begin{math_stroke}[Second Dimensional Split (Applying Fubini to the 2D slice)]
\[
\int_{[a_2, b_2] \times [a_3, b_3]} f \, d(x_2, x_3) = \underbrace{\int_{a_2}^{b_2}}_{\text{``integrate over } [a_2, b_2] \text{''}} \left( \underbrace{\int_{a_3}^{b_3} f(x_1, x_2, x_3) \, dx_3}_{\text{``both } x_1 \text{ and } x_2 \text{ are fixed''}} \right) \underbrace{dx_2}_{\text{``then we move } x_2 \text{''}}
\]
\end{math_stroke}

\begin{spoken_clean}[28:00 - 30:00]
If you substitute this all back into our first equation, you see that it gives you exactly the triple nested formula I wrote before: an integral from $a_1$ to $b_1$, $a_2$ to $b_2$, and $a_3$ to $b_3$. End of the proof! 

If you look at the notes, it is the exact same proof in $n$ dimensions. It simply introduces notation for these successive domains that become smaller and smaller as we integrate out variables one by one. There are two things in every theorem: the theoretical proof, and the practical application. To pass the exam, you need to at least be able to apply the theorem without hesitation. If you have some doubts about the abstract proof, that understanding will come with time, but understanding how to apply it is absolutely crucial.
\end{spoken_clean}

\begin{nice-box}
\begin{theorem}[Concluded 3D Fubini Iteration]
Substituting the nested marginals back into the original split yields the final iterated form:
\[
\int_K f = \underbrace{\int_{a_1}^{b_1} \int_{a_2}^{b_2} \int_{a_3}^{b_3} f(x_1, x_2, x_3) \, dx_3 \, dx_2 \, dx_1}_{\text{``exactly the triple nested formula I wrote before''}}
\]
This recursive logic extends to $n$ dimensions.
\end{theorem}
\end{nice-box}
%% ==========================================
% SECTION 5: CONCLUDING THE PROOF & FINAL REMARKS
% ==========================================

\begin{spoken_clean}[30:32 - 32:44]
Now you substitute this back, and this yields exactly the formula I wrote before: the triple integral over the bounds $a_1$ to $b_1$, $a_2$ to $b_2$, and $a_3$ to $b_3$ of $f(x_1, x_2, x_3) \, dx_3 \, dx_2 \, dx_1$. And that concludes the proof! 

With this in mind, if you consult the lecture notes, you will see it is the exact same proof in $n$ dimensions. It simply introduces notation for these successive domains that become smaller and smaller as we integrate out variables one by one, and you get the exact same result.

As with every theorem, there are two aspects: the theoretical proof and the practical application. To pass the exam, you need to at least be able to apply the theorem without hesitation. If you still have some doubts about the abstract proof, that understanding will come with time, but understanding how to apply it is absolutely crucial. Are there any questions about this? If it is okay, we can move on.
\end{spoken_clean}

\begin{didactic_insight}[Bridging Abstraction and Calculation]
The professor explicitly acknowledges the cognitive load of $n$-dimensional abstraction. By explicitly giving students "permission" to focus on the mechanical application of Fubini's theorem first, he lowers their anxiety, ensuring they don't get stuck on the rigorous $n$-dimensional notation while they are still building their intuition in 3D space.
\end{didactic_insight}


% ==========================================
% SECTION 6: DIFFERENTIATION UNDER THE INTEGRAL SIGN
% ==========================================

\section{Differentiation Under the Integral Sign}

\begin{nice-box}[Board State: Transitioning to New Topic (32:56 - 34:07)]
The professor erases the Fubini proof and prepares a fresh board. He writes the new major section header: \textbf{Theorem: Differentiation under the integral sign}.
\end{nice-box}

\begin{spoken_clean}[32:56 - 37:37]
To wrap up our unit on integration, I will introduce a final, crucial tool: differentiation under the integral sign. This is something you will absolutely need, not just in this class, but extensively in physics and partial differential equations. 

Consider a scenario in physics where you are integrating a density of charge, say $q(x,t)$, which depends on both space $x$ and time $t$. The charge density varies over time, but you integrate it over a fixed spatial volume to find the total mass or total charge. If you then want to compute the rate of change of that total mass with respect to time, you must take the derivative of the integral. Because the spatial volume is fixed, you can pass this time derivative directly inside the integral sign, applying it as a partial derivative to the density function itself. You will use this constantly in physics, often without even saying why it is true. Today, we will establish exactly why.
\end{spoken_clean}

\begin{math_stroke}[Formal Statement of the Theorem (37:40 - 41:03)]
The professor formalizes the physical analogy into rigorous mathematics. 
He defines a compact spatial set $K$ and a time interval $[a,b]$.

\textbf{Setup:}
Let $K \subset \mathbb{R}^n$ be a compact set. Let $f: K \times [a,b] \to \mathbb{R}$ defined by $(x,t) \mapsto f(x,t)$.

\textbf{Assumptions:}
Assume $f$ is continuous, and the partial derivative with respect to the last variable, $\frac{\partial f}{\partial t}$, is also continuous on the closed set $K \times [a,b]$.

\textbf{Conclusion:}
For every $t_0 \in (a,b)$:
\[
\frac{d}{dt} \bigg|_{t_0} \int_K f(x,t) \, dx = \int_K \frac{\partial f}{\partial t}(x,t_0) \, dx
\]
\end{math_stroke}

\begin{spoken_clean}[37:40 - 41:03]
Let us formalize this. Let $K$ be a compact subset in $\mathbb{R}^n$, representing our space. We define a real-valued function $f$ mapping from $K \times [a,b]$—where the interval represents our time domain—to $\mathbb{R}$. We assume not only that $f$ is continuous, but also that its partial derivative with respect to the time variable, $\frac{\partial f}{\partial t}$, is continuous on this closed product space. 

Under these conditions, for any time $t_0$ in the open interval, the total derivative with respect to $t$ of the integral of $f(x,t)$ over $K$ is exactly equal to the integral over $K$ of the partial derivative of $f$ evaluated at $t_0$. When you take a function that depends on space and time, and you integrate with respect to space, the only free variable remaining is time. So, taking the derivative on the outside is a standard Analysis 1 derivative of a single variable, but pushing it inside turns it into a partial derivative.
\end{spoken_clean}

\begin{redundant_explanation}[Swapping Operations]
The core of this theorem is the commutative property of limit operations under specific continuity conditions. The integral is a limit of Riemann sums, and the derivative is a limit of a difference quotient. Swapping "$\frac{d}{dt}$" and "$\int$" is only strictly valid because $\frac{\partial f}{\partial t}$ is uniformly continuous on the compact domain $K \times [a,b]$.
\end{redundant_explanation}

% ==========================================
% SECTION 7: PROOF OF THEOREM (PART 1)
% ==========================================

\section{Proof: Differentiation Under the Integral Sign}

\begin{nice-box}[Board State: The Difference Quotient (41:03 - 44:55)]
The professor begins the proof by writing out the formal definition of the derivative using the limit definition, letting $h \to 0$. He defines an auxiliary function $F(t)$ to simplify the notation on the board.
\end{nice-box}

\begin{math_stroke}[Setting up the Limit Definition]
Let us define a 1-dimensional function for the spatial integral:
\[
F(t) = \int_K f(x,t) \, dx
\]
We want to compute $F'(t_0)$. By the definition of the derivative:
\[
F'(t_0) = \lim_{h \to 0} \frac{F(t_0 + h) - F(t_0)}{h}
\]
Substituting the integrals back into the limit:
\[
F'(t_0) = \lim_{h \to 0} \frac{\int_K f(x, t_0 + h) \, dx - \int_K f(x, t_0) \, dx}{h}
\]
By the linearity of the integral, we combine the terms:
\[
F'(t_0) = \lim_{h \to 0} \int_K \frac{f(x, t_0 + h) - f(x, t_0)}{h} \, dx
\]
\end{math_stroke}

\begin{spoken_clean}[41:03 - 44:55]
Let us begin the proof. We want to compute the derivative with respect to $t$ at $t_0$. By definition, this is the limit as $h$ approaches zero of the difference quotient. To make this cleaner, let's introduce a new function $F(t)$ representing the spatial integral of $f(x,t)$. 

We are simply computing $F'(t_0)$, which is the limit of $\frac{F(t_0 + h) - F(t_0)}{h}$. Because the integral is a linear operator, we can group these two integrals into one. We bring the difference quotient inside, dividing the difference $f(x, t_0 + h) - f(x, t_0)$ by $h$. 

This is starting to look much better—it strongly resembles the definition of the partial derivative! All that remains is to formally justify that we can pass the limit as $h$ goes to zero inside the integral. We will finish that rigorous justification after the break, showing that this converges exactly to the integral of the partial derivative.
\end{spoken_clean}

\begin{meta_note}[Lecture Break]
At 44:55, the professor pauses the mathematical derivation, having successfully motivated the difference quotient. The rigorous bounding of the limit via the Mean Value Theorem and Uniform Continuity will occur after the break.
\end{meta_note}

%\include{content/4-16-wednesday-second-half-part-1}
%\section{Example: The Gaussian Integral}

\begin{nice-box}[Problem Setup: Evaluating the Gaussian Integral]
The professor moves to a fresh blackboard to provide a capstone example. This problem combines the new theory of improper integrals, Fubini's theorem, and the Change of Variables formula.
\end{nice-box}

\begin{exercise}[The Gaussian Integral]
Compute the integral over the entirety of $\mathbb{R}^2$ of the 2D Gaussian function:
\[
\int_{\mathbb{R}^2} e^{-x^2-y^2} \, dx \, dy
\]
\end{exercise}

\begin{proof}[Step-by-Step Evaluation of the Gaussian Integral]

\textbf{Step 1: Exhausting $\mathbb{R}^2$ with Expanding Balls}
\begin{spoken_clean}[16:58 - 19:22]
Now, we will do a practical example so that you see exactly what I mean by exhausting an open set. I will do an example to practice both improper integrals and the change of variables formula. Let us compute the integral over $\mathbb{R}^2$ of the function $e^{-x^2 - y^2}$. 

This is an improper integral because $\mathbb{R}^2$ is an open set; it is not a compact set. However, the function $e^{-x^2 - y^2}$ is strictly positive everywhere. Therefore, we can safely apply our new definition for non-negative functions. How will we exhaust $\mathbb{R}^2$? I will exhaust it by taking a sequence of expanding balls. We will take the limit as $R$ goes to infinity of the integral over the ball of radius $R$, centered at zero, of this function. This sequence of nested balls clearly exhausts the entire space.
\end{spoken_clean}

\begin{math_stroke}[Exhaustion by Expanding Balls]
Because the integrand is non-negative, we apply the sequential evaluation property, exhausting $\mathbb{R}^2$ using closed balls $B_R(0)$:
\[
\int_{\mathbb{R}^2} e^{-x^2-y^2} \, dx \, dy = \underbrace{\lim_{R \to \infty} \int_{B_R(0)} e^{-x^2-y^2} \, dx \, dy}_{\text{``limit as } R \text{ goes to infinity over the ball of radius } R \text{''}}
\]
\end{math_stroke}

\textbf{Step 2: Transforming to Polar Coordinates}
\begin{spoken_clean}[19:22 - 24:00]
Now, how do I actually compute this integral over the ball? I will use a change of variables, specifically polar coordinates. As we saw before, $x^2 + y^2$ becomes $r^2$. The Jacobian determinant of this transformation gives me an extra factor of $r$ (i.e., $|\det J| = r$). The limits of integration for the angle $\phi$ are $0$ to $2\pi$, and for the radius $r$, from $0$ up to $R$. 

We can then apply Fubini's corollary to evaluate it iteratively. The integral with respect to $\phi$ is trivial and just gives $2\pi$. What is the primitive of $r e^{-r^2}$? It is simply $-\frac{1}{2} e^{-r^2}$. Evaluating this primitive from $0$ to $R$, we plug in the boundaries and obtain exactly $\pi(1 - e^{-R^2})$ (since evaluating at $R$ and $0$ gives $-\frac{1}{2}e^{-R^2} - (-\frac{1}{2}) = \frac{1}{2}(1 - e^{-R^2})$, which we multiply by $2\pi$).
\end{spoken_clean}

\begin{math_stroke}[Polar Coordinates and Fubini's Theorem]
Applying the polar substitution $x = r\cos\phi$, $y = r\sin\phi$ with the Jacobian volume scaling factor $|\det J| = r$:
\[
\int_{B_R(0)} e^{-x^2-y^2} \, dx \, dy = \int_0^{2\pi} \int_0^R \underbrace{e^{-r^2}}_{\text{``becomes } r^2 \text{''}} \underbrace{r}_{\text{``extra factor of } r \text{''}} \, dr \, d\phi
\]
Using Fubini's theorem to separate the variables and evaluate the iterated integral:
\[
= \underbrace{\left( \int_0^{2\pi} 1 \, d\phi \right)}_{\text{``just gives } 2\pi \text{''}} \left( \int_0^R r e^{-r^2} \, dr \right) = 2\pi \underbrace{\left[ -\frac{1}{2} e^{-r^2} \right]_0^R}_{\text{``evaluating this primitive from } 0 \text{ to } R \text{''}} 
\]
\[
= 2\pi \left( -\frac{1}{2} e^{-R^2} - \left(-\frac{1}{2}\right) \right) = \underbrace{\pi (1 - e^{-R^2})}_{\text{``we obtain exactly this''}}
\]
\end{math_stroke}

\textbf{Step 3: Taking the Limit to Infinity}
\begin{spoken_clean}[24:00 - 25:00]
Finally, we must take the limit as our radius $R$ goes to infinity. As $R$ becomes infinitely large, the term $e^{-R^2}$ decays to $0$, leaving us with exactly $\pi$. So, the improper integral over all of $\mathbb{R}^2$ for this Gaussian function is simply $\pi$.
\end{spoken_clean}

\begin{math_stroke}[Evaluating the Limit]
\[
\lim_{R \to \infty} \pi (1 - \underbrace{e^{-R^2}}_{\text{``decays to } 0 \text{''}}) = \pi (1 - 0) = \underbrace{\pi}_{\text{``integral over all of } \mathbb{R}^2 \text{''}}
\]
\end{math_stroke}

\textbf{Step 4: The 1D Gaussian Integral (Exhausting by Squares)}
\begin{spoken_clean}[25:00 - 29:00]
But wait! Because the value of an improper integral is completely independent of the sequence of compact sets we choose to exhaust the space, we can exhaust $\mathbb{R}^2$ in another way. Let us exhaust it by squares instead of balls! We define our sequence of expanding squares as $K_R = [-R, R] \times [-R, R]$. Because the sequence doesn't matter, the limit as $R$ goes to infinity over these expanding squares must also evaluate to exactly $\pi$. 

Notice that our function $e^{-x^2 - y^2}$ can be split into $e^{-x^2} e^{-y^2}$. Using Fubini's theorem, we can split this two-dimensional integral over the square into the product of two identical, independent one-dimensional integrals. This gives us the square of the integral from $-R$ to $R$ of $e^{-x^2}$. Taking the limit as $R$ goes to infinity, we find that the integral from $-\infty$ to $\infty$ of $e^{-x^2}$ squared is $\pi$. Therefore, taking the square root, the one-dimensional integral is exactly $\sqrt{\pi}$ (i.e., $\int_{-\infty}^\infty e^{-x^2} \, dx = \sqrt{\pi}$). This is the famous Gaussian integral, which you will see absolutely everywhere in probability!
\end{spoken_clean}

\begin{math_stroke}[Exhaustion by Squares and Algebraic Factoring]
By the topological independence of the exhausting sequence, we evaluate the same improper integral using squares $K_R = [-R, R]^2$:
\[
\lim_{R \to \infty} \int_{[-R,R]^2} e^{-x^2-y^2} \, dx \, dy = \pi
\]
Using Fubini's theorem to cleanly factor the separable integrand across the rectangular domain:
\[
\int_{-R}^R \int_{-R}^R \underbrace{e^{-x^2} e^{-y^2}}_{\text{``split into...''}} \, dx \, dy = \left( \int_{-R}^R e^{-x^2} \, dx \right) \left( \int_{-R}^R e^{-y^2} \, dy \right) = \underbrace{\left( \int_{-R}^R e^{-x^2} \, dx \right)^2}_{\text{``the square of the integral''}}
\]
Equating the limits to our known 2D volume yields the standard 1D Gaussian integral:
\[
\left( \int_{-\infty}^\infty e^{-x^2} \, dx \right)^2 = \pi \implies \underbrace{\int_{-\infty}^\infty e^{-x^2} \, dx = \sqrt{\pi}}_{\text{``the famous Gaussian integral''}}
\]
\end{math_stroke}
\end{proof}



%%%%%%%%%%%%%%%%%%%%%%%%%%%%%%%%%%%%%%%%%%%

%\chapter{Real Analysis: Measuring Volumes on Submanifolds}

\begin{didactic_insight}[The Gavel and the "Extra Circus"]
The professor opens the lecture holding a toy gavel, explicitly preparing the students for an "extra circus." This playful, theatrical prop serves a distinct pedagogical purpose: acknowledging the escalating difficulty of the material ("The Final Boss of Analysis II") while keeping the classroom atmosphere grounded and engaged. He then beautifully recalls his earlier, informal ``potato'' analogies to show how far their rigor has come.
\end{didactic_insight}

% ==========================================
% SECTION 1: THE FINAL BOSS (DIVERGENCE THEOREM)
% ==========================================

\section{Recall: The Final Boss of Analysis II}

\begin{spoken_clean}[00:00:00 - 00:01:28]
So I brought back the gavel because today I expect extra circus, okay? 

So, remember that I told you the first day the final goal, or the final boss of Analysis II, or one of the most representative theorems, is this Divergence Theorem. And the first days I needed to state it very informally using ``potatoes'' and stuff because you didn't know mathematically any of the elements. But now we saw the statement again some weeks ago, and we already knew some parts. And now we know almost all of it. Right? 

I'll remind you of the statement. Now we know many elements: $\Omega$ in $\mathbb{R}^n$ is an open, bounded set. The boundary $\partial \Omega$ will be a $C^1$ submanifold. This we know what it means. Now, $F$ will be a $C^1$ function up to the boundary, taking values in $\mathbb{R}^n$. This is a vector field. And the theorem says that the integral over $\Omega$ of this combination of partial derivatives—the derivative $i$ of the component $i$ of the vector field—and you integrate with respect to the Jordan measure, Riemann integral... this is equal to this integral over the boundary of the domain with respect to the surface measure of the boundary. 

Now we know, and this is $F$ multiplied by the normal vector, which we already saw what is a tangent vector and a normal vector to the boundary. So more or less everything you know what it is now, except for one thing: this integral over a surface, of the differential of volume in the surface. We know how to integrate on the full space, but not over a curved surface.
\end{spoken_clean}

\begin{orangeformula}[Divergence Theorem (The Final Boss)]
Let $\Omega \subset \mathbb{R}^n$ be an open, bounded set, and $\partial \Omega$ a $C^1$ submanifold. Let $F \in C^1(\overline{\Omega}, \mathbb{R}^n)$ be a vector field. The theorem states:
\begin{equation} \label{eq:divergence}
\int_\Omega \sum_{i=1}^n \partial_i F_i \, dx = \int_{\partial \Omega} F \cdot \underbrace{\nu}_{\text{Unit vector normal to } \partial \Omega} \, \underbrace{dS}_{\text{Surface measure (Unknown)}}
\end{equation}
\end{orangeformula}

% ==========================================
% SECTION 2: QUEST LOG - INTEGRATION OVER SUBMANIFOLDS
% ==========================================

\section{Defining Integrals over d-dimensional Submanifolds}

\begin{spoken_clean}[00:01:28 - 00:03:50]
And this is where we are going today. We are going to learn, to define what is the integral over the surface. Okay, so next step in our quest: define integral over $d$-dimensional submanifolds. 

And to define an integral, as it happened in the case of $\mathbb{R}^n$, first we need to define how to measure the volume of sets over a submanifold. So today, towards this goal, what we will do is measure volumes over submanifolds of dimension $d$. 

Towards this goal, we start by the easiest case. What is the easiest possible submanifold? It is the submanifold of dimension one. Right? So let's start by $d=1$. A submanifold of dimension one, we call it a curve. And we will see how to measure length. So we need to introduce length.
\end{spoken_clean}

\begin{math_stroke}[Quest Log: Volumes and Curves]
\textbf{Next Step:} Define $\int$ over $d$-dim submanifolds.

\textbf{Today's Objective:} Measure volumes over submanifolds of $\dim d$.

We build intuition starting with the simplest case:
\begin{itemize}
    \item $\dim d = 1 \implies$ \textbf{Curve} $\implies$ ``Volume'' is defined as \textbf{Length}.
\end{itemize}
\end{math_stroke}

% ==========================================
% SECTION 3: PARAMETRIZED CURVES AND ARC LENGTH
% ==========================================

\section{Formalizing Curves in $\mathbb{R}^n$}

\begin{spoken_clean}[03:50 - 05:20]
So, how do we rigorously characterize a one-dimensional submanifold? We use what we call a parametrized curve. We start with a simple interval, let us call it $[a, b]$, on the real line, and we map it into our higher-dimensional space $\mathbb{R}^n$ using a function that we will denote as $\gamma$. 

Now, we cannot just use any arbitrary function. We want our curve to be smooth, without any jagged edges or sudden jumps. Therefore, we require $\gamma$ to be continuously differentiable (i.e., $\gamma \in C^1([a, b], \mathbb{R}^n)$). Furthermore, to avoid pathological cases where the curve stops or creates a sharp cusp, we must enforce that the derivative vector—the velocity—is never the zero vector at any point $t$ (i.e., $\gamma'(t) \neq 0$ for all $t$). If these conditions hold, we say the curve is regular.
\end{spoken_clean}

\begin{math_stroke}[Definition of a Regular Parametrized Curve]
We define a \textbf{parametrized curve} as a map:
\[
\gamma: [a, b] \to \mathbb{R}^n
\]
We say that $\gamma$ is a \textbf{regular curve} if it satisfies two conditions:
\begin{enumerate}
    \item $\gamma \in C^1([a, b], \mathbb{R}^n)$ 
    \item $\gamma'(t) \neq 0$ for all $t \in [a, b]$ 
\end{enumerate}
\end{math_stroke}

\begin{spoken_clean}[05:20 - 07:15]
Once we have this regular curve, the immediate next question is: how do we compute its length? Think about physics for a second. If $\gamma(t)$ represents your position in space at time $t$, then the derivative $\gamma'(t)$ is your velocity vector. If you want to know how far you have traveled, you do not just look at the velocity vector; you look at your speed, which is the magnitude of that velocity vector (i.e., $\|\gamma'(t)\|$).

If you integrate your speed over the entire duration of your journey, from time $a$ to time $b$, you obtain the total distance traveled. This total distance is exactly what we define as the length of the curve.
\end{spoken_clean}

\begin{orangeformula}[Length of a Parametrized Curve]
\begin{definition}
Let $\gamma: [a, b] \to \mathbb{R}^n$ be a regular parametrized curve. The \textbf{length} of $\gamma$, denoted as $L(\gamma)$, is defined by the integral of the Euclidean norm of its derivative:
\begin{equation} \label{eq:arc_length}
L(\gamma) = \int_a^b \|\gamma'(t)\| \, dt
\end{equation}
\end{definition}
\end{orangeformula}

\begin{didactic_insight}[The Speedometer Analogy]
The professor relies on a classical kinematic analogy here. The formula for arc length is presented not just as an abstract analytical construct, but as a physical one: Length is the accumulation of speed over time. This effectively maps the abstract $dt$ integration back to familiar real-world intuition.
\end{didactic_insight}

% ==========================================
% SECTION 4: THE AXIOMS OF LENGTH
% ==========================================

\section{Geometric Independence and the Axioms of Length}

\begin{spoken_clean}[07:15 - 09:30]
Now, here is the conceptual leap we must make. The length $L(\gamma)$ that we just defined depends explicitly on the parametrization $\gamma$. It depends on the specific function and the specific time interval. But length, fundamentally, is a geometric property of the set of points in space, not a property of how fast you drive along them. 

If two different people drive along the exact same path, but one drives twice as fast, the geometric length of the path they covered is exactly the same, even though their parametrizations are completely different. Therefore, if our formula is to make any geometric sense, we need to prove that it is invariant under reparametrization.
\end{spoken_clean}

\begin{math_stroke}[Reparametrization Setup]
Let $\phi: [c, d] \to [a, b]$ be a bijective, $C^1$ function with $\phi'(s) \neq 0$ for all $s \in [c, d]$.
We define a reparametrization of $\gamma$ as:
\[
\tilde{\gamma}(s) = \gamma(\phi(s))
\]

\textbf{Our analytical goal:} We must show that $L(\tilde{\gamma}) = L(\gamma)$.
\end{math_stroke}

\begin{spoken_clean}[09:30 - 11:00]
Before we prove that our integral formula actually survives this reparametrization mathematically, we have to take a step back. We are constructing a measure of volume—specifically, a 1-dimensional volume, which is length. But what does it even mean to be a valid measure of length? 

We cannot just write down an integral, compute it, and blindly accept it as "length." There are fundamental, axiomatic properties that any concept of length must satisfy in Euclidean space. If our formula fails even one of these, it is not a valid measure. So we must set the ground rules.
\end{spoken_clean}
%\chapter{Real Analysis: Measuring Volumes on Submanifolds}

\begin{metanote}[Lecture Commencement]
The professor walks to the center of the blackboard, holding a small toy gavel. The blackboard is clean, and the camera is fixed on the main panel.
\end{metanote}

\begin{didacticinsight}[The Gavel and the ``Extra Circus'']
The professor opens the lecture holding a toy gavel, explicitly preparing the students for an ``extra circus''. This playful, theatrical prop serves a distinct pedagogical purpose: acknowledging the escalating difficulty of the material (``The Final Boss of Analysis II'') while keeping the classroom atmosphere grounded and engaged. He beautifully recalls his earlier, informal ``potato'' analogies to show how far the class's rigor has progressed.
\end{didacticinsight}

% ==========================================
% SECTION 1: THE FINAL BOSS (DIVERGENCE THEOREM)
% ==========================================

\section{Recall: The Final Boss of Analysis II}

\begin{spokenclean}[00:00:00 - 00:00:58]
So I brought back the gavel because today I expect extra circus, okay? 

So, remember that I told you on the first day that the final goal, or the final boss of Analysis II, or one of the most representative theorems, is this Divergence Theorem. And in the first days, I needed to state it very informally using ``potatoes'' and stuff because you did not know mathematically any of the elements. But now we saw the statement again some weeks ago, and we already knew some parts. And now we know almost all of it. Right? 
\end{spokenclean}

\begin{nicebox}[Board Action: Divergence Setup]
The professor places the gavel on the desk and begins writing the formal elements of the Divergence Theorem on the left side of the blackboard, underlining the conditions for the domain and the vector field.
\end{nicebox}

\begin{spokenclean}[00:00:58 - 00:01:55]
I will remind you of the statement. Now we know many of the elements: $\Omega$ in $\mathbb{R}^n$ is an open, bounded set. The boundary $\partial \Omega$ will be a $C^1$ submanifold. This we know what it means. Now, $F$ will be a $C^1$ function up to the boundary, taking values in $\mathbb{R}^n$. This is a vector field. 

And the theorem says that the integral over $\Omega$ of this combination of partial derivatives—the derivative $i$ of the component $i$ of the vector field—and you integrate with respect to the Jordan measure (i.e., the Riemann integral, $dx$)... this is equal to this integral over the boundary of the domain with respect to the surface measure of the boundary.
\end{spokenclean}

\begin{spokenclean}[00:01:55 - 00:02:40]
Now we know, and this is $F$ multiplied by the normal vector, which we already saw what is a tangent vector and a normal vector to the boundary. So more or less everything you know what it is now, except for one thing: this integral over a surface, of the differential of volume in the surface (i.e., $dS$). We know how to integrate on the full space, but not over a curved extra surface.
\end{spokenclean}

\begin{orangeformula}[Divergence Theorem (The Final Boss)]
\begin{theorem}[Divergence Theorem]
Let $\Omega \subset \mathbb{R}^n$ be an open, bounded set, and $\partial \Omega$ a $C^1$ submanifold. Let $F \in C^1(\overline{\Omega}, \mathbb{R}^n)$ be a vector field. The theorem states:
\begin{equation} \label{eq:divergence}
\int_\Omega \sum_{i=1}^n \partial_i F_i \, dx = \int_{\partial \Omega} F \cdot \underbrace{\nu}_{\text{Unit vector normal to } \partial \Omega} \, \underbrace{dS}_{\text{Surface measure (Unknown)}}
\end{equation}
\end{theorem}
\end{orangeformula}

\begin{redundantexplanation}[The Unknown Surface Measure]
The entire motivation for this lecture hinges on $dS$. While the left side of the equation uses the standard $n$-dimensional volume measure $dx$ that the students have mastered via the Riemann integral and Fubini's Theorem, the right side requires integrating over an $(n-1)$-dimensional curved manifold embedded in $\mathbb{R}^n$. We do not yet have a mathematical definition for how to rigorously quantify this ``curved volume''.
\end{redundantexplanation}

% ==========================================
% SECTION 2: QUEST LOG - INTEGRATION OVER SUBMANIFOLDS
% ==========================================

\section{Defining Integrals over d-dimensional Submanifolds}

\begin{spokenclean}[00:02:40 - 00:03:35]
And this is where we are going today. We are going to learn, to define what is the integral over the surface. Okay, so next step in our quest: define the integral over $d$-dimensional submanifolds. 

And to define an integral, as it happened in the case of $\mathbb{R}^n$, first we need to define how to measure the volume of sets over a submanifold. So today, towards this goal, what we will do is measure volumes over submanifolds of dimension $d$. 
\end{spokenclean}

\begin{nicebox}[Board Action: Quest Log]
The professor draws a distinct box on the middle panel to represent the current ``Quest Log'', a recurring motif in his lectures to keep the overarching goal visible.
\end{nicebox}

\begin{spokenclean}[00:03:35 - 00:04:20]
Towards this goal, we start by the easiest case. What is the easiest possible submanifold? It is the submanifold of dimension one. Right? So let us start by $d=1$. A submanifold of dimension one, we call it a curve. And we will see how to measure length. So we need to introduce length.
\end{spokenclean}

\begin{mathstroke}[Quest Log: Volumes and Curves]
\textbf{Next Step:} Define $\int$ over $d$-dim submanifolds.

\textbf{Today's Objective:} Measure volumes over submanifolds of $\dim d$.

We build intuition starting with the simplest case:
\begin{itemize}
    \item $\dim d = 1 \implies$ \textbf{Curve} $\implies$ ``Volume'' is defined as \textbf{Length}.
\end{itemize}
\end{mathstroke}

% ==========================================
% SECTION 3: PARAMETRIZED CURVES AND ARC LENGTH
% ==========================================

\section{Formalizing Curves in \texorpdfstring{$\mathbb{R}^n$}{Rn}}    

\begin{spokenclean}[00:04:20 - 00:05:18]
So, how do we rigorously characterize a one-dimensional submanifold? We use what we call a parametrized curve. We start with a simple interval, let us call it $[a, b]$, on the real line, and we map it into our higher-dimensional space $\mathbb{R}^n$ using a function that we will denote as $\gamma$. 

Now, we cannot just use any arbitrary function. We want our curve to be smooth, without any jagged edges or sudden jumps. Therefore, we require $\gamma$ to be continuously differentiable (i.e., $\gamma \in C^1([a, b], \mathbb{R}^n)$).
\end{spokenclean}

\begin{nicebox}[Board Action: Injectivity and Regularity]
He underlines the $C^1$ condition heavily. Then he draws a curve intersecting itself to visually demonstrate why injectivity might be considered, before shifting focus entirely to the non-zero derivative condition to prevent cusps.
\end{nicebox}

\begin{spokenclean}[00:05:18 - 00:06:10]
Furthermore, to avoid pathological cases where the curve stops or creates a sharp cusp, we must enforce that the derivative vector—the velocity—is never the zero vector at any point $t$ (i.e., $\gamma'(t) \neq 0$ for all $t$). If these conditions hold, we say the curve is regular.
\end{spokenclean}

\begin{mathstroke}[Definition of a Regular Parametrized Curve]
We define a \textbf{parametrized curve} as a map:
\[
\gamma: [a, b] \to \mathbb{R}^n
\]
We say that $\gamma$ is a \textbf{regular curve} if it satisfies two conditions:
\begin{enumerate}
    \item $\gamma \in C^1([a, b], \mathbb{R}^n)$ 
    \item $\gamma'(t) \neq 0$ for all $t \in [a, b]$ 
\end{enumerate}
\end{mathstroke}

\begin{spokenclean}[00:06:10 - 00:07:05]
Once we have this regular curve, the immediate next question is: how do we compute its length? Think about physics for a second. If $\gamma(t)$ represents your position in space at time $t$, then the derivative $\gamma'(t)$ is your velocity vector. If you want to know how far you have traveled, you do not just look at the velocity vector; you look at your speed, which is the magnitude of that velocity vector (i.e., $\|\gamma'(t)\|$).
\end{spokenclean}

\begin{spokenclean}[00:07:05 - 00:07:15]
If you integrate your speed over the entire duration of your journey, from time $a$ to time $b$, you obtain the total distance traveled. This total distance is exactly what we define as the length of the curve.
\end{spokenclean}

\begin{orangeformula}[Length of a Parametrized Curve]
\begin{definition}[Arc Length]
Let $\gamma: [a, b] \to \mathbb{R}^n$ be a regular parametrized curve. The \textbf{length} of $\gamma$, denoted as $L(\gamma)$, is defined by the integral of the Euclidean norm of its derivative:
\begin{equation} \label{eq:arc_length}
L(\gamma) = \int_a^b \|\gamma'(t)\| \, dt
\end{equation}
\end{definition}
\end{orangeformula}

% ==========================================
% SECTION 4: THE AXIOMS OF LENGTH
% ==========================================

\section{Geometric Independence and the Axioms of Length}

\begin{spokenclean}[00:07:15 - 00:08:30]
Now, here is the conceptual leap we must make. The length $L(\gamma)$ that we just defined depends explicitly on the parametrization $\gamma$. It depends on the specific function and the specific time interval. But length, fundamentally, is a geometric property of the set of points in space, not a property of how fast you drive along them. 
\end{spokenclean}

\begin{spokenclean}[00:08:30 - 00:09:30]
If two different people drive along the exact same path, but one drives twice as fast, the geometric length of the path they covered is exactly the same, even though their parametrizations are completely different. Therefore, if our formula is to make any geometric sense, we need to prove that it is invariant under reparametrization.
\end{spokenclean}

\begin{mathstroke}[Reparametrization Setup]
Let $\phi: [c, d] \to [a, b]$ be a bijective, $C^1$ function with $\phi'(s) \neq 0$ for all $s \in [c, d]$.
We define a reparametrization of $\gamma$ as:
\[
\tilde{\gamma}(s) = \gamma(\phi(s))
\]

\textbf{Our analytical goal:} We must show that $L(\tilde{\gamma}) = L(\gamma)$.
\end{mathstroke}

\begin{spokenclean}[00:09:30 - 00:10:30]
Before we prove that our integral formula actually survives this reparametrization mathematically, we have to take a step back. We are constructing a measure of volume—specifically, a 1-dimensional volume, which is length. But what does it even mean to be a valid measure of length?
\end{spokenclean}

\begin{spokenclean}[00:10:30 - 00:11:30]
We cannot just write down an integral, compute it, and blindly accept it as ``length''. There are fundamental, axiomatic properties that any concept of length must satisfy in Euclidean space. If our formula fails even one of these, it is not a valid measure. So we must set the ground rules.
\end{spokenclean}

% ==========================================
% SECTION 5: THE FOUR AXIOMS OF LENGTH
% ==========================================

\begin{spokenclean}[00:11:30 - 00:12:30]
If we are going to define a concept as fundamental as length, it cannot just be a random formula we pull from the air. It has to behave in the way our physical universe expects lengths to behave. If it does not, then we have defined something else entirely. 

So, before we even look at the integral of the norm of the velocity, we must ask ourselves: what properties \textit{must} length satisfy? What are the absolute baseline requirements? There will be four. So there will be four fundamental axioms that govern any valid measurement of length.
\end{spokenclean}

\begin{nicebox}[Classroom Interaction: The Catchbox]
The professor walks back to the main desk. He writes the numbers 1 through 4 on the far right blackboard, leaving the definitions blank. He then picks up a bright red throwable ``Catchbox'' microphone and tosses it to a student in the audience.
\end{nicebox}

\begin{spokenclean}[00:12:30 - 00:13:15]
Okay, so what properties must length satisfy? I have written four empty slots on the board. Maybe you can tell me two, and I will complete the other two. Who wants to catch this? Tell me what properties length must satisfy. Who has a guess for the most obvious property a length should have?
\end{spokenclean}

\begin{studentquestion}
It has to be positive? You can't have a negative length. And you just add them together if you have multiple pieces?
\end{studentquestion}

\begin{spokenclean}[00:13:15 - 00:14:15]
Exactly! First, we have positivity, but more rigorously, we need normalization. The length of the straight unit interval (i.e., $[0,1] \times \{0\}$ in $\mathbb{R}^n$) must exactly equal 1. And yes, it must be additive. If we split a curve into two pieces, the total length must be the sum of the lengths of the individual pieces. 
\end{spokenclean}

\begin{spokenclean}[00:14:15 - 00:15:00]
Now I will complete the other two, which are the geometric anchors. We already discussed one: it should be independent of the chosen parametrization. The geometric length does not care how fast you draw the curve. And finally, invariance under Euclidean isometry. If I rotate or translate the curve rigidly in space, its length cannot change.
\end{spokenclean}

\begin{orangeformula}[The Four Axioms of Length]
\begin{definition}[Axioms of Length]
For a definition of length $L(\gamma)$ to be geometrically valid for a curve $\gamma: [a, b] \to \mathbb{R}^n$, it \textit{must} satisfy:
\begin{enumerate}
    \item \textbf{Independence of Parametrization:} $L(\gamma \circ \phi) = L(\gamma)$ for any valid reparametrization $\phi$.
    \item \textbf{Additivity:} For any $c \in (a, b)$, $L(\gamma|_{[a, b]}) = L(\gamma|_{[a, c]}) + L(\gamma|_{[c, b]})$.
    \item \textbf{Euclidean Isometry Invariance:} Rigid rotations and translations do not alter $L(\gamma)$.
    \item \textbf{Normalization:} The standard unit interval $[0,1] \times \{0\}$ has a length of $1$.
\end{enumerate}
\end{definition}
\end{orangeformula}


% ==========================================
% SECTION 6: PROVING REPARAMETRIZATION INVARIANCE
% ==========================================

\section{Proving the Invariance}

\begin{spokenclean}[00:15:00 - 00:16:30]
Let us focus on the first property we established: Independence of Parametrization. We need to rigorously prove that our integral formula actually honors this axiom. 

Suppose we have our reparametrization $\phi$ mapping the interval $[c, d]$ to $[a, b]$. We assume $\phi$ is strictly increasing, so its derivative is strictly positive (i.e., $\phi'(s) > 0$). We want to compute the length of the new, reparametrized curve $\tilde{\gamma}(s) = \gamma(\phi(s))$. 
\end{spokenclean}

\begin{nicebox}[Board Action: The Chain Rule Setup]
The professor walks over to the center panel of the blackboard, pointing at the definition of the regular curve, and begins writing out the composition of the functions to apply the Chain Rule.
\end{nicebox}

\begin{spokenclean}[00:16:30 - 00:17:45]
How do we find its velocity vector? We just apply the Chain Rule. The derivative $\tilde{\gamma}'(s)$ is simply the derivative of the outer function evaluated at the inner function, multiplied by the derivative of the inner function (i.e., $\tilde{\gamma}'(s) = \gamma'(\phi(s)) \cdot \phi'(s)$). 
\end{spokenclean}

\begin{mathstroke}[Applying the Chain Rule]
Given the reparametrized curve $\tilde{\gamma}(s) = \gamma(\phi(s))$ for $s \in [c, d]$.

By the Chain Rule, the velocity vector is:
\[
\tilde{\gamma}'(s) = \gamma'(\phi(s)) \phi'(s)
\]
Taking the Euclidean norm of both sides:
\[
\|\tilde{\gamma}'(s)\| = \|\gamma'(\phi(s)) \phi'(s)\| = \|\gamma'(\phi(s))\| |\phi'(s)|
\]
Since we assume $\phi$ is strictly increasing, $\phi'(s) > 0$, yielding $|\phi'(s)| = \phi'(s)$.
\[
\implies \|\tilde{\gamma}'(s)\| = \|\gamma'(\phi(s))\| \phi'(s)
\]
\end{mathstroke}

\begin{spokenclean}[00:17:45 - 00:19:00]
Now we plug this norm directly into our length formula for $\tilde{\gamma}$. We get the integral from $c$ to $d$ of $\|\gamma'(\phi(s))\| \phi'(s) \, ds$. 

Does this look familiar? It is exactly the setup for the 1-dimensional Change of Variables formula—or as you likely learned it in Analysis I, Integration by Substitution. 
\end{spokenclean}

\begin{spokenclean}[00:19:00 - 00:20:00]
If we substitute $t = \phi(s)$, then our differential $dt$ is exactly $\phi'(s) ds$. The limits of integration change from $c$ and $d$ to $\phi(c)=a$ and $\phi(d)=b$. And as if by magic, the formula collapses back perfectly into the length of the original curve $\gamma$. The geometry is entirely preserved. We have proven the first axiom.
\end{spokenclean}

\begin{mathstroke}[Integration by Substitution]
Computing the length of the reparametrized curve $\tilde{\gamma}$:
\begin{align*}
L(\tilde{\gamma}) &= \int_c^d \|\tilde{\gamma}'(s)\| \, ds \\
&= \int_c^d \|\gamma'(\phi(s))\| \underbrace{\phi'(s) \, ds}_{\text{Substitution: } dt}
\end{align*}

Let $t = \phi(s)$. The differential becomes $dt = \phi'(s) ds$.
The boundaries transform as: $t(c) = \phi(c) = a$ and $t(d) = \phi(d) = b$.

\begin{align*}
L(\tilde{\gamma}) &= \int_a^b \|\gamma'(t)\| \, dt \\
&= L(\gamma)
\end{align*}

\begin{explanationofsteps}
We have successfully demonstrated that the total arc length is invariant under an orientation-preserving reparametrization, confirming the first fundamental axiom of length. The 1D substitution rule elegantly offsets the scaling of the velocity vectors.
\end{explanationofsteps}
\end{mathstroke}

%\chapter{Real Analysis: Cavalieri's Principle}

\begin{spoken-clean}[00:00:00 - 00:00:04]
Okay, so.
\end{spoken-clean}

\begin{spoken-clean}[00:00:04 - 00:00:10]
Let me continue with the statement. So, we were trying to state this theorem that I informally described as cutting the ``potato'', right? But now we are trying to be more serious.
\end{spoken-clean}

\begin{spoken-clean}[00:00:10 - 00:00:19]
And I was introducing this notation. So, this theorem, you will see, is a bit the opposite of the Change of Variables formula. The Change of Variables formula is a very simple statement and a very hard proof. This is a very complicated statement, but the proof is very easy, right?
\end{spoken-clean}

\section{The General Setup: Cavalieri's Principle for Sets}

\begin{spoken-clean}[00:00:19 - 00:00:39]
So, we have our measurable set. Let me do a drawing. So, we have our set $A$, right? That is contained in a box. Here, we will put the $y$-axis, and this is the $x$-axis. This is actually $\mathbb{R}^k$, and on the $y$-axis we have $\mathbb{R}^{n-k}$, right? And the whole space is $\mathbb{R}^n$, the product of the two (i.e., $\mathbb{R}^n = \mathbb{R}^k \times \mathbb{R}^{n-k}$).
\end{spoken-clean}

\begin{math-stroke}[Geometric Visualization Setup for Cavalieri's Principle]
\begin{center}
\begin{tikzpicture}[scale=1.5]
    % Axes
    \draw[->, thick, profgreen] (0,0) -- (4,0) node[below right] {$x \in \mathbb{R}^k$};
    \draw[->, thick, profgreen] (0,0) -- (0,3) node[above left] {$y \in \mathbb{R}^{n-k}$};
    \node[above right] at (4,3) {$\mathbb{R}^n$};

    % Set A
    \draw[thick] (1,1) .. controls (0.5,2.5) and (2.5,2.8) .. (3.5,1.5) .. controls (3,0.5) and (1.5,0.2) .. cycle;
    \node at (2,1.5) {$A$};

    % Bounding Box (dashed)
    \draw[dashed, gray] (0.5,0.5) rectangle (3.8,2.8);

    % Slice at x
    \draw[profred, thick] (2, 0.5) -- (2, 2.8); % The vertical slice
    \node[profred, below] at (2, 0.5) {$x$};
\end{tikzpicture}
\end{center}

\begin{explanation-of-steps}
We visualize an arbitrary Jordan measurable set $A$ within $\mathbb{R}^n$, which is decomposed into a Cartesian product of $\mathbb{R}^k$ (our $x$-axis) and $\mathbb{R}^{n-k}$ (our $y$-axis). This setup allows us to consider ``slices'' of the set $A$.
\end{explanation-of-steps}
\end{math-stroke}

\begin{spoken-clean}[00:00:39 - 00:01:23]
Okay, now we are saying that our set is contained in some box, but this is not very important. So, this is the box, a very large box. And now for $x$ inside of my box, so this will be some $x$ here, an $x$ point, I will consider the slice. So, the relevant piece is this slice over here. This is what I'm calling $A_x$. Well, the projection... So, $A_x$ will be the projection onto this axis, right? So, you take this set and you project it. So, you take the $y$-coordinate, so given the $x$, you take the $y$-coordinates of the pairs $(x,y)$, is over there. The $y$-coordinates of the pairs $(x,y)$ that belong to $A$.
\end{spoken-clean}

\begin{math-stroke}[Formalizing the Slices of Set A]
Let $A \subset \mathbb{R}^n$ be a Jordan measurable set. We can express $\mathbb{R}^n$ as a product space $\mathbb{R}^k \times \mathbb{R}^{n-k}$, where we denote points as $(x,y)$.
For any fixed $x \in \mathbb{R}^k$, we define the slice $A_x$ as:
\[
A_x = \{ y \in \mathbb{R}^{n-k} \mid (x,y) \in A \}
\]
For $A$ to be Jordan measurable, it must be bounded. We assume $A$ is contained within a large box, so $A_x$ will be a bounded subset of $\mathbb{R}^{n-k}$, which we denote as:
\[
A_x \subset (-c,c)^{n-k} \subset \mathbb{R}^{n-k}
\]
\end{math-stroke}

\begin{spoken-clean}[00:01:47 - 00:02:20]
And now this set, here is very nice, because in the drawing it's an interval, right? But this, we are in very high dimensions, maybe. So, this is just a subset of... $A_x$ is a subset of $\mathbb{R}^{n-k}$. It's a bounded subset, but we just know that it's just containing $(-c,c)^{n-k}$, right? In the box. And what we can do, because we know that $A$ is Jordan measurable, but maybe the slices are not Jordan measurable. But we can surely consider these marginals, the $\bar{\varphi}$ that was the... So, we can always measure the $\mu_{n-k}^{\text{out}}$ and the $\mu_{n-k}^{\text{in}}$. So, the measure from inside and outside of this set is always well-defined. And this is what defines the function $\overline{\varphi}$ and $\underline{\varphi}$ of $x$, right?
\end{spoken-clean}

\begin{math-stroke}[Upper and Lower Measures of Slices]
Given the slices $A_x$, even if $A_x$ itself is not Jordan measurable (meaning its inner and outer measures may not coincide), we can always define its $(n-k)$-dimensional outer measure and inner measure:
\[
\overline{\varphi}(x) = \mu_{n-k}^{\text{out}}(A_x) \quad \text{and} \quad \underline{\varphi}(x) = \mu_{n-k}^{\text{in}}(A_x)
\]
These functions, $\overline{\varphi}(x)$ and $\underline{\varphi}(x)$, map from $\mathbb{R}^k$ to $\mathbb{R}$.
\end{math-stroke}

\begin{orange-formula}
\begin{theorem}[Cavalieri's Principle for Jordan Measurable Sets]
Let $A \subset \mathbb{R}^n$ be a Jordan measurable set. We consider the decomposition $\mathbb{R}^n = \mathbb{R}^k \times \mathbb{R}^{n-k}$, with coordinates $(x,y)$. For each $x \in [-c,c]^k$, we define the slices $A_x = \{y \in \mathbb{R}^{n-k} \mid (x,y) \in A \}$. We then define the upper and lower marginal functions:
\[
\overline{\varphi}(x) = \mu_{n-k}^{\text{out}}(A_x) \quad \text{and} \quad \underline{\varphi}(x) = \mu_{n-k}^{\text{in}}(A_x)
\]
Then, both $\overline{\varphi}: [-c,c]^k \to \mathbb{R}$ and $\underline{\varphi}: [-c,c]^k \to \mathbb{R}$ are Riemann integrable functions over $[-c,c]^k$, and their integrals are equal to the total $n$-dimensional measure of $A$:
\[
\int_{[-c,c]^k} \overline{\varphi}(x) \, dx = \int_{[-c,c]^k} \underline{\varphi}(x) \, dx = \mu_n(A)
\]
\end{theorem}
\end{orange-formula}

\begin{spoken-clean}[00:02:20 - 00:03:20]
So, this is the setup, and the theorem says that under these assumptions, under these definitions, then both functions, the $\bar{\varphi}(x)$ and $\underline{\varphi}(x)$ are Riemann integrable. And the integral of them in the set $[-c,c]^k$ is equal to the integral of the other one, which is equal to the measure of $A$. So this integral is... You sum so the $\bar{\varphi}$ is the measure of each slice. You are taking the measure of these slices. Now you can tell me, maybe for some slices this measure is not well-defined, because the outer measure and the inner measure do not coincide. I tell you, don't worry, just put any, the outer or the inner. I don't care. Because when you integrate, the outer or the inner or anything in between, it will give you always the measure... It will recover the measure of your set $A$.
\end{spoken-clean}

\begin{proof}
\begin{spoken-clean}[00:03:20 - 00:04:31]
Okay, so let's prove that. So maybe to do the proof, I will do a drawing. So remember my notation: this is the $y$, this is the $x$-axis. This is my set $A$, right? And now, always when we want to prove something about Jordan measure, remember that we were approximating our set from inside and outside. We will take the approximation from inside, right? Something like this, some approximation from inside. And then there will be the approximation from outside. And now, since our set is measurable, it's Jordan measurable, we can take a very small pixel size, $2^{-p}$. This will be the pixel size. And we can take an inner and an outer approximation of our set by the dyadic... so by by pixels, right? By the dyadic sets. This is the inner, this is the outer. And now, the difference between these two is very, um, is is is less than $\epsilon$. We can take the volume of the red minus the volume of the yellow is less than $\epsilon$ for every $\epsilon$ we want.
\end{spoken-clean}

\begin{math-stroke}[Approximation with Dyadic Cubes ($2^{-p}$)]
\begin{center}
\begin{tikzpicture}[scale=1.5]
    % Axes
    \draw[->, thick, profgreen] (0,0) -- (4,0) node[below right] {$x$};
    \draw[->, thick, profgreen] (0,0) -- (0,3) node[above left] {$y$};
    \node[above right] at (4,3) {$A$}; % Label A near the set

    % Set A (irregular shape)
    \draw[thick] (1,1) .. controls (0.5,2.5) and (2.5,2.8) .. (3.5,1.5) .. controls (3,0.5) and (1.5,0.2) .. cycle;

    % Inner approximation (yellow/light green)
    \draw[fill=yellow!20, draw=profgreen, thick] (1,1) -- (1,2) -- (2,2) -- (2,1) -- (2.5,1) -- (2.5,0.5) -- (3,0.5) -- (3,1.5) -- (2.5,1.5) -- (2.5,2.5) -- (1.5,2.5) -- (1.5,1) -- cycle;

    % Outer approximation (red/light red)
    \draw[fill=red!10, draw=profred, thick, dashed] (0.5,0.5) -- (0.5,2.8) -- (3.8,2.8) -- (3.8,0.5) -- cycle;
    \draw[draw=profred, thick] (0.8,0.8) rectangle (3.2,2.7); % Outer boundary simplified for visual clarity

    % Labels for slices
    \node at (0.7, 2.7) {\small $A$};
    \node[profred, right] at (3.8, 2.7) {$2^{-p}$}; % Pixel size

    % Slice at x
    \draw[blue, thick] (2, 0) -- (2, 2.8); % Vertical slice
    \node[blue, below] at (2,0) {$x$};
\end{tikzpicture}
\end{center}
\begin{explanation-of-steps}
We approximate the Jordan measurable set $A$ from inside and outside using unions of dyadic cubes (``pixels'') of side length $2^{-p}$. Let $F_p$ be the union of dyadic cubes entirely contained within $A$ (inner approximation, shown in green), and $G_p$ be the union of dyadic cubes that intersect $A$ (outer approximation, shown in red). We can choose $p$ large enough such that the difference in their measures is arbitrarily small, $\mu_n(G_p) - \mu_n(F_p) < \epsilon$.
\end{explanation-of-steps}
\end{math-stroke}

\begin{spoken-clean}[00:04:31 - 00:05:20]
And now we will we will see what happens when we take a slice of this set. When we take a slice of this set like this, let's say we take $x$ and we slice, right? What we see is that the the slice of my... the slice of $A$ contains the slices of the yellow and is contained in the slices of the red. But when you when you start to move $x$ like this, when you start to move $x$, you will see that whenever you move $x$ less than the size of the pixel, you're not jumping here. So the the function, the function, um, so the marginal's idea: the marginal densities computed for the two, for the yellow set and for the, um, the dyadic, the dyadic inner and outer approximations. So what is this again? What is this, um, these marginal? The marginal is just take the the yellow set, for instance. Compute, given some $x$, what is the intersection, right? And you see, because this is a union of cubes, when you intersect with a hyperplane parallel to the to the, or with a plane parallel to to to the $y$-axis, this also gives you a union of cubes, just of lesser dimensions. And then the measure of this, then it's constant. When you move $x$ here in the dyadic cube, this you will always get the same. You will always get the same, all the time you will get the same. And then only when you jump to the next dyadic cube, maybe you jump here, because it's like a staircase. So the marginal densities computed in the inner and outer approximations are the dyadic step functions. This is the key observation. And now we will write this, so so they are dyadic step functions in $x$. So this will, they approximate the Riemann integral. The dyadic step functions are what we need, or what we use, to approximate the Riemann integral.
\end{spoken-clean}

\begin{math-stroke}[Dyadic Cube Definition and Product Space]
\textbf{Given a dyadic cube (pixel $2^{-p}$) in $\mathbb{R}^n$:}
A general dyadic cube $Q_p^n(z)$ of side length $2^{-p}$ at point $z \in \mathbb{R}^n$ is defined as:
\[
Q_p^n(z) = z + 2^{-p} [0,1)^n
\]
If $z$ is expressed in product coordinates as $z=(a,b)$ where $a \in \mathbb{R}^k$ and $b \in \mathbb{R}^{n-k}$, then the $n$-dimensional dyadic cube can be factored into lower-dimensional dyadic cubes:
\[
Q_p^n(z) = Q_p^k(a) \times Q_p^{n-k}(b)
\]
Similarly, for a dyadic cube of dimension $m$:
\[
Q_p^m(\bar{z}) = \bar{z} + 2^{-p} [0,1)^m \subset \mathbb{R}^m
\]
\end{math-stroke}

\begin{spoken-clean}[00:05:20 - 00:06:01]
So, this is how we are going to proceed. This is the picture, and now we will write it formally. So, the idea is that, given a dyadic cube of size $2^{-p}$ in $\mathbb{R}^n$, we have... so we will say, okay, $Q_p^n(z)$ is $z + 2^{-p}[0,1)^n$, right? And since now there will be several dimensions, let me put the dimension here as a sub-index $Q_p^n(z)$. This is $n$.
\end{spoken-clean}

\begin{spoken-clean}[00:06:01 - 00:06:57]
Now, if $z$ is of the form $(a,b)$ where $a \in \mathbb{R}^k$ and $b \in \mathbb{R}^{n-k}$, right? In this coordinate splitting, I take a $z$ that has coordinates $a$ and $b$, right? Then this cube $Q_p^n(z)$ is equal to $Q_p^k(a) \times Q_p^{n-k}(b)$. Right? Where I'm putting the $Q_p^m$ of some $\bar{z}$ is $\bar{z} + 2^{-p}[0,1)^m$, this is a subset of $\mathbb{R}^m$. Right?
\end{spoken-clean}

\begin{math-stroke}[Dyadic Cubes and Approximations]
Given dyadic sets $F$ (inner) and $G$ (outer) such that $F \subset A \subset G$, and their measures satisfy:
\[
\mu(G) - \mu(F) < \epsilon
\]
This holds for dyadic sets of pixel size $2^{-p}$ where $p \in \mathbb{N}$ is sufficiently large.
\end{math-stroke}

\begin{spoken-clean}[00:06:57 - 00:08:14]
Something, something very complicated to write, but very trivial on the picture. I'm saying that the cube over here, the dyadic cube over here, is always the product of a cube over here times a dyadic cube over here, right? A cube in the in the dyadic cube in the full space is the product of a dyadic cube of the same size here, times a dyadic cube of the same size here. And a dyadic set, a dyadic set, the yellow or the red, to both applies, is just a finite union of these two cubes, right? Okay. So, now, given $\epsilon > 0$, let $p \in \mathbb{N}$ and $F,G$ be dyadic sets such that $F \subset A \subset G$, and the measure of $G$ minus the measure of $F$ (so the measure of the larger minus the measure of the smaller) is less than $\epsilon$, right? So, this is the definition of a Jordan measurable set. We can trap it between dyadic sets, and the measure of the outer minus the measure of the inner is less than $\epsilon$.
\end{spoken-clean}

\begin{math-stroke}[Slicing Dyadic Approximations]
For $x \in [-c,c]^k$, define the slices of the outer and inner dyadic approximations:
\[
G_x = \{ y \in \mathbb{R}^{n-k} \mid (x,y) \in G \}
\]
\[
F_x = \{ y \in \mathbb{R}^{n-k} \mid (x,y) \in F \}
\]
\textbf{Observation:} If $Q_p^k(a) \subset \mathbb{R}^k$ is a dyadic cube of side length $2^{-p}$, then for any $x, x' \in Q_p^k(a)$, the slices of the dyadic sets are identical:
\[
F_x = F_{x'} \quad \text{and} \quad G_x = G_{x'}
\]
\end{math-stroke}

\begin{spoken-clean}[00:08:14 - 00:09:50]
And now we, for all $x \in [-c,c]^k$, define the slices $G_x = \{y \in \mathbb{R}^{n-k} \mid (x,y) \in G \}$, and similarly with $F_x$. Okay? So we do the slices of $G_x$. Now, observation, that is what I was trying to say with the drawing: if $Q_p(a) \subset \mathbb{R}^k$ is some dyadic cube of pixel size $2^{-p}$ (so of size of side length $2^{-p}$), then what we see is because every cube in $\mathbb{R}^n$ is a product of cubes. You see that $F_x = F_{x'}$ for all $x, x' \in Q_p(a)$.
\end{spoken-clean}

\begin{math-stroke}[Defining Step Functions from Slice Measures]
Let $f(x)$ be the measure of the inner dyadic slice and $g(x)$ be the measure of the outer dyadic slice:
\[
f(x) := \mu(F_x)
\]
\[
g(x) := \mu(G_x)
\]
These functions $f(x)$ and $g(x)$ are constant in each dyadic cube $Q_p^k(a)$, and are zero outside a bounding box $(-C,C)^k$. This implies $f(x)$ and $g(x)$ are dyadic step functions.
\end{math-stroke}

\begin{spoken-clean}[00:09:50 - 00:10:33]
As long as the base point is moving in this cube, the slices cannot change; the slice is always constant. And the same is with the $G$. What does it mean? It means that if we now define... So, since this function, so since these functions are constant in each dyadic cube, this implies that $F$ and $G$ are dyadic step functions. And notice something else. That for each slice... So, this $F$ is measuring...
\end{spoken-clean}

\begin{math-stroke}[Relating Measures to Riemann Integrals]
Since $F_x \subset A_x \subset G_x$, we have the following inequality for the measures of the slices:
\[
\underline{f(x)} \equiv \mu(F_x) \leq \mu_{n-k}^{\text{in}}(A_x) \leq \mu_{n-k}^{\text{out}}(A_x) \leq \mu(G_x) \equiv \overline{g(x)}
\]
Integrating over the entire domain $[-c,c]^k$:
\[
\int_{[-c,c]^k} f(x) \, dx \leq \int_{[-c,c]^k} \underline{\varphi}(x) \, dx \leq \int_{[-c,c]^k} \overline{\varphi}(x) \, dx \leq \int_{[-c,c]^k} g(x) \, dx
\]
For dyadic step functions, the integral is simply the sum of the volumes of the dyadic cubes:
\[
\int_{[-c,c]^k} f(x) \, dx = 2^{-n \cdot p} \cdot (\text{number of dyadic cubes in } F) = \mu_n(F)
\]
And similarly for $g(x)$:
\[
\int_{[-c,c]^k} g(x) \, dx = 2^{-n \cdot p} \cdot (\text{number of dyadic cubes in } G) = \mu_n(G)
\]
Substituting these into the inequality:
\[
\mu_n(F) \leq \int_{[-c,c]^k} \underline{\varphi}(x) \, dx \leq \int_{[-c,c]^k} \overline{\varphi}(x) \, dx \leq \mu_n(G)
\]
\end{math-stroke}

\begin{spoken-clean}[00:10:52 - 00:11:42]
This $F$ is measuring this, right? For each $x$, this $F$ is measuring the the number of pixels that are in this slice. And the $G$ is measuring the number of pixels that are in the bigger slice of the red. So, what we always have is that for every $x$, $F_x$ is contained in $A_x$, is contained in $G_x$, right? We have this inclusion, and this implies that the measure of $F_x$ is less than the... The measure of an inner approximation is always less than the inner measure, that is less than the outer measure of $A_x$, that is less than or equal to the measure of the outer approximation. This for every $x$.
\end{spoken-clean}

\begin{spoken-clean}[00:11:42 - 00:14:47]
And this is... by definition, this function over here, this is $\underline{\varphi}(x)$, right? Of $\underline{\varphi}(x)$. And this is less than the... This, the integral of this... I can put the upper integral of the $\bar{\varphi}(x)$ and this is less than the integral of $G$. Which is equal to $\mu_n(G)$. Okay? Does that make sense? So, this is the chain of inequalities. Again, I repeat. This equality is just combinatorial, counting cubes. This equality is because this is a lower approximation of my integral. So, if you have a function that is below your function, and this is the dyadic step function, because the lower integral is defined as the supremum of those... this is always less than or equal to this, right? If two functions are ordered, like this the the, obviously the lower integral is always less than the upper integral. And and the $\underline{\varphi}$, that is the $\mu_{n-k}^{\text{in}}$, is less than the $\mu_{n-k}^{\text{out}}$. This is the $\bar{\varphi}$. So, the $\mu_{n-k}^{\text{out}}$ is the outer marginal. So, this inequality is also trivial, right? It's just because of ordering. And this is again the same function, but from a different argument, but from outside now. So, we know that this $G$ is an approximation of the $\bar{\varphi}$; it is above the $\bar{\varphi}$, the $G$. So, the integral of the $G$ must be bigger than the upper integral of the $\bar{\varphi}$, right? And again, this is combinatorial, the integral of $G$ is the measure of $G$. Counting pixels.
\end{spoken-clean}

\begin{spoken-clean}[00:14:47 - 00:15:39]
But now we remember that this minus this was $\epsilon$. But $\mu(G_\epsilon) - \mu(F_\epsilon)$ is less than $\epsilon$. And actually, not only this, but the value... but we know that $G_\epsilon$ is contained in $A$, and sorry, $F_\epsilon$ is contained in $A$, and $A$ is contained in $G_\epsilon$. So, the measure of this is trapped in between, right? And then, putting this together with this, what does it what does it tell you? So, if I if I send $\epsilon$ to zero, all these integrals in between that do not depend on $\epsilon$, you see, this is the depends on $\epsilon$. These integrals need to be the same. When I send $\epsilon$ to zero, you have my squeeze. So, this integral that a priori is less than this needs to be the same. This proved that these functions, both functions, are Riemann integrable, because the upper integral and the lower integral will coincide, right? And actually, all the two functions have the same value of the integral that is $\mu_n(A)$, because these these two values, right? So, what I wanted to say is that $\mu_n(F_\epsilon)$ is less than $\mu_n(A)$, is less than $\mu_n(G_\epsilon)$, right? So, when you send $\epsilon$ to zero, this value and this value are converging to the common value $\mu_n(A)$, right? And everything that is in between needs to converge to the same value.
\end{spoken-clean}

\begin{spoken-clean}[00:15:39 - 00:15:42]
So, this is how we prove the theorem.
\end{spoken-clean}

\end{proof}

\begin{spoken-clean}[00:15:46 - 00:15:47]
Is there any question about this?
\end{spoken-clean}

\begin{spoken-clean}[00:16:00 - 00:16:21]
Okay, so, as I said, the main difficulty of this theorem is the notation. I need to introduce a lot of notation, and this is always confusing to anyone, right? So, when you introduce new notation, until you load it to your RAM memory, all this notation, you are confused by the notation, right? So, if you got lost at some point, you remember the "potato", that is what is important to understand the intuition, and then you go back with a cup of tea on the proof, and then you read step-by-step and you will see there's nothing mysterious. It's just a lot of notation.
\end{spoken-clean}

\begin{spoken-clean}[00:16:21 - 00:16:22]
Okay? And the key idea to remember is the counting argument.
\end{spoken-clean}

\section{Extension: Fubini's Theorem for Functions}

\begin{spoken-clean}[00:16:22 - 00:16:27]
So, now we will do an extension of this of this theorem.
\end{spoken-clean}

\begin{spoken-clean}[00:16:27 - 00:16:28]
That is for for functions.
\end{spoken-clean}

\begin{spoken-clean}[00:16:28 - 00:16:32]
It will actually be an immediate consequence of this theorem.
\end{spoken-clean}

\begin{math-stroke}[Fubini's Theorem: Erasing for Clarity]
\begin{center}
\vspace{1cm} % Placeholder for erased content
\end{center}
\begin{explanation-of-steps}
The professor erases the board to prepare for Fubini's Theorem. He emphasizes that the previous proof for Cavalieri's Principle for sets is foundational and that Fubini's Theorem for functions is a direct extension, leveraging the same underlying logic.
\end{explanation-of-steps}
\end{math-stroke}

\begin{spoken-clean}[00:16:58 - 00:17:02]
So, Fubini's Theorem is the following. So, now we have a function...
\end{spoken-clean}

\begin{orange-formula}
\begin{theorem}[Fubini's Theorem]
Let $A \subset \mathbb{R}^n$ be a bounded Jordan measurable set, and $f: A \to \mathbb{R}$ be a Riemann integrable function over $A$.
We decompose $\mathbb{R}^n = \mathbb{R}^k \times \mathbb{R}^{n-k}$, with coordinates $(x,y)$.
We extend $f$ to a function $\tilde{f}: \mathbb{R}^n \to \mathbb{R}$ by setting $\tilde{f}(x,y) = f(x,y)$ if $(x,y) \in A$ and $\tilde{f}(x,y) = 0$ if $(x,y) \notin A$.
For $x \in \mathbb{R}^k$, we define the marginal functions:
\[
\underline{\varphi}(x) = \underline{\int_{\mathbb{R}^{n-k}} \tilde{f}_x(y) \, dy}
\]
\[
\overline{\varphi}(x) = \overline{\int_{\mathbb{R}^{n-k}} \tilde{f}_x(y) \, dy}
\]
where $\tilde{f}_x(y) = \tilde{f}(x,y)$ (treating $x$ as a constant parameter).

Then, $\underline{\varphi}(x)$ and $\overline{\varphi}(x)$ are Riemann integrable functions over $\mathbb{R}^k$, and their integrals are equal to the total integral of $f$ over $A$:
\[
\int_A f(x,y) \, d(x,y) = \int_{\mathbb{R}^k} \underline{\varphi}(x) \, dx = \int_{\mathbb{R}^k} \overline{\varphi}(x) \, dx
\]
\end{theorem}
\end{orange-formula}

\begin{spoken-clean}[00:17:02 - 00:17:22]
So, we have some $A \subset \mathbb{R}^n$ (bounded). Jordan measurable. Right? Now, as before, I will take the $\mathbb{R}^n$ as $\mathbb{R}^k \times \mathbb{R}^{n-k}$, and we'll call the coordinates here $(x,y)$. Or do I need Jordan measurable? Maybe I don't even need Jordan measurable. I just need set bounded. And then I have a function $f$ from $A$ to $\mathbb{R}$ that is Riemann integrable over $A$.
\end{spoken-clean}

\begin{spoken-clean}[00:17:36 - 00:18:53]
And now I define... I will define $\tilde{f}$ just $f$ multiplied by the characteristic function of $A$, right? Or if you want, this is the... So, $\tilde{f}$ will be $\tilde{f}(x,y)$ will be $f(x,y)$ if $(x,y)$ belongs to $A$, and $0$ if $(x,y)$ does not belong to $A$, right? This is just a trivial way to extend the function to be defined everywhere, right? I just take zero where it is not defined. And then, um, I define the marginals. So for $x \in \mathbb{R}^k$, so for $x \in \mathbb{R}^k$, define the marginals. So, as before, $\underline{\varphi}$ will be the lower integral of... Now I take $\tilde{f}$ of $x$, so let me put it like this: $\tilde{f}_x(y) \, dy$. Where I define $\tilde{f}_x(y)$ simply as $\tilde{f}(x,y)$. I'm putting it like this to be entirely sure that this is a function of... So, that you understand that this is a function defined in... So, this is with a tilde, tilde tilde. This is a function defined for $y$ in $\mathbb{R}^{n-k}$.
\end{spoken-clean}

\begin{math-stroke}[Fubini's Theorem Statement]
\[
\int_A f \, d(x,y) = \int_{\mathbb{R}^k} \underline{\varphi}(x) \, dx = \int_{\mathbb{R}^k} \overline{\varphi}(x) \, dx
\]
\end{math-stroke}

\begin{spoken-clean}[00:18:53 - 00:19:14]
So, this given $x$, $x$ is fixed now, $x$ is fixed. So, this is a function of $y$, this is a function in $\mathbb{R}^{n-k}$. So, you can compute an integral or a lower integral is always well-defined, the the lower integral and the upper integral. You can compute the lower integral of this function with respect to the variables $y$, right? And we can do the same with the upper. These are definitions.
\end{spoken-clean}

\begin{spoken-clean}[00:19:14 - 00:19:53]
Now what the theorem says is that the integral in $A$ of $f(x,y)$, so this is the integral in $\mathbb{R}^n$, right? Is the same as the integral of, um, now I need to integrate with respect to $x$. $x$ lives in $\mathbb{R}^k$, right? Integral in $\mathbb{R}^k$ of $\underline{\varphi}(x) \, dx$. And this is the same as taking the upper.
\end{spoken-clean}

\begin{math-stroke}[Visualizing Fubini's Theorem]
\begin{center}
\begin{tikzpicture}[scale=1.5]
    % Axes (Muted to push them to the background, keeping focus on the surfaces)
    \draw[->, thick, gray!80!black] (0,0,0) -- (4,0,0) node[right] {$y$ axis};
    \draw[->, thick, gray!80!black] (0,0,0) -- (0,3,0) node[above] {$z = f(x,y)$};
    \draw[->, thick, gray!80!black] (0,0,0) -- (0,0,4) node[below left] {$x$ axis};

    % Domain in xy-plane
    \draw[dashed, profblue, thick] (1,0,1) -- (3,0,1) -- (3,0,3) -- (1,0,3) -- cycle;
    \node[profblue] at (2,-0.3,3.5) {Base Domain $A = X \times Y$};

    % The Slice at constant x_0 (Drawn FIRST so it is properly occluded by the surface)
    \draw[thick, profred, fill=profred!20, opacity=0.7] (1,0,2) -- (3,0,2) -- (3,2.0,2) to[out=150,in=-10] (1,1.5,2) -- cycle;
    
    % Slice Annotations
    \draw[dashed, profred, thick] (0,0,2) -- (1,0,2);
    \node[left, profred] at (0,0,2) {$x_0$};
    % Moved the Area label to the right to avoid overlapping the surface
    \draw[<-, profred, thick] (2.5, 0.8, 2) -- (3.8, 1.5, 2) node[right, profred, fill=white, inner sep=1pt] {Area $= \int f(x_0, y) \, dy$};

    % Surface (Drawn SECOND so it is in front, opacity adjusted)
    \draw[thick, proforange, fill=proforange!20, opacity=0.8] (1,2,1) to[out=20,in=160] (3,2.5,1) to[out=-20,in=70] (3,1.5,3) to[out=160,in=-20] (1,1,3) to[out=70,in=-20] cycle;
    \node[proforange] at (2,2.8,1) {Surface $z = f(x,y)$};
\end{tikzpicture}
\end{center}
\begin{explanation-of-steps}
The visual clarifies the core concept: the inner integral calculates the area of the 2D cross-sectional slice at a constant $x_0$. Fubini's Theorem simply states that integrating this varying 1D area function across the $x$-axis accumulates the full 3D volume (the hypograph).
\end{explanation-of-steps}
\end{math-stroke}

\begin{spoken-clean}[00:19:53 - 00:20:16]
Which again... So, this is similar to the theorem. And again, going to the picture, what what this tells me is that... This is the $y$, this is the $x$. Is that if you want to integrate a function over a body, you can integrate first for each $x$, you consider the slice of your body, and you can integrate over the slices, and then sum over all the slices.
\end{spoken-clean}

\begin{spoken-clean}[00:20:16 - 00:21:01]
This thing is integrating over the slices. And you will tell me, okay, maybe for some slices, if the integral is not well-defined, then just take the upper or the lower integral. You don't care. So, for most of the slices the integral will be defined, a posteriori. But if for some it's not defined, just take any, the upper or the lower integral. When you integrate over the slices, you sum over all the slices, you get the total integral of your function. And this will be very useful, and we will use next time to compute that integral over there that we left. Okay? You will see that with this observation, and we will prove a corollary that is very very nice to use, um, we can compute in one line that integral over there. Okay? But let me say one thing about the proof. The proof of this is just the same, you repeat almost line by line the proof that I gave for the for the Cavalieri Principle. Or if you want to be even smarter, what you can do is you can use the... So, you can take your function, you can decompose it into negative part and positive part. And remember that any integral is just a volume. It's just measuring a volume in one dimension more. So, what you can apply if you apply the Cavalieri Principle to the hypograph, so to the hypograph of this function $F$, you get exactly that. So, remember that any integral is just a volume of a hypograph. So, this is just the Cavalieri Principle applied to the hypograph of the function $F$ when $F$ is positive, or when $F$ is negative. And then you decompose. Okay. So, we are going to use this next time.
\end{spoken-clean}


%\begin{spoken-clean}[00:00:00 - 00:00:12]
Let me continue with the statement. We were trying to state this theorem that I informally described as \qt{cutting the potato}. But now we are trying to be more serious, and I will introduce the proper notations.
\end{spoken-clean}

\begin{math-stroke}[Transition from Spherical Coordinates to Slicing Ideas]
The previous discussion led to the volume integral for the unit ball:
\[
\mu_3(B_1) = \mu_3(\Psi(D)) = \int_{(0,1)\times(0,2\pi)\times(-\pi/2,\pi/2)} y_1^2 \cos y_3 \, dy = \, \textcolor{red}{?}
\]
This naturally introduces the concept of ``Slicing ideas'':
\begin{center}
%\includegraphics[width=0.4\textwidth]{path/to/cavalieri_principle_concept.png} % Placeholder for visible board image
\end{center}
\begin{explanation-of-steps}
The board visually summarizes the core idea of Cavalieri's principle: for an $n$-dimensional body, its total volume can be computed by summing the "areas" (measures) of its $(n-1)$-dimensional slices, multiplied by an infinitesimal thickness (i.e., $\text{volume}(A) \approx \sum A_x \, dx$).
\end{explanation-of-steps}
\end{math-stroke}

\begin{spoken-clean}[00:00:12 - 00:02:20]
This theorem, you will see, is a bit the opposite of the Change of Variables formula. The Change of Variables formula has a very simple statement but a very hard proof. This has a very complicated statement, but the proof is very easy. We have our measurable set. Let me do a drawing of it. We have our set $A$ that is contained in a bounding box. Here, we will put the $y$-axis, and this is the $x$-axis. This is actually $\mathbb{R}^k$, and on the $y$-axis, we have $\mathbb{R}^{n-k}$. The whole space is $\mathbb{R}^n$, the product of the two (i.e., $\mathbb{R}^n = \mathbb{R}^k \times \mathbb{R}^{n-k}$).
\end{spoken-clean}

\begin{math-stroke}[Geometric Visualization Setup]
\begin{center}
\begin{tikzpicture}[scale=1.0]
    % Axes
    \draw[->, thick, profgreen] (-0.5,0) -- (4,0) node[right] {$x \in \mathbb{R}^k$};
    \draw[->, thick, profblue] (0,-0.5) -- (0,3) node[above] {$y \in \mathbb{R}^{n-k}$};
    \node[above right] at (3.5, 2.5) {$\mathbb{R}^n$};

    % Bounding Box (dashed)
    \draw[dashed] (0.5,0.5) rectangle (3.5,2.5);

    % Set A
    \draw[thick] (1.5,1) to[out=45,in=135] (2.5,1.5) to[out=-45,in=45] (2,2) to[out=-135,in=-45] (1,1.5) to[out=135,in=-135] cycle;
    \node at (2, 1.5) {$A$};

    % Slice at x
    \draw[red, very thick] (1.75, 0.5) -- (1.75, 2.5);
    \node[below, xshift=-0.1cm, red] at (1.75, 0) {$x$};
    \node[right, red] at (1.75, 1.5) {$A_x$};
\end{tikzpicture}
\end{center}
\begin{explanation-of-steps}
The professor draws a set $A$ within a bounding box in $\mathbb{R}^n$, split into $\mathbb{R}^k$ ($x$-axis) and $\mathbb{R}^{n-k}$ ($y$-axis). A vertical slice at a specific $x$ value is indicated in red and labeled $A_x$. This $A_x$ represents the $y$-coordinates corresponding to a fixed $x$ within $A$.
\end{explanation-of-steps}
\end{math-stroke}

\begin{spoken-clean}[00:02:20 - 00:03:37]
We are saying that our set is contained in some very large bounding box. Now, for an $x$ inside of my box, I will consider the vertical slice. The relevant piece is this slice over here. This is what I am calling $A_x$. The set $A_x$ will be the projection onto this vertical axis. So, you take this slice and you project it.
\end{spoken-clean}

\begin{nice-box}[Cavalieri's Principle for Sets]
\begin{theorem}[Cavalieri's Principle for Sets]
Let $\mathbb{R}^n = \mathbb{R}^k \times \mathbb{R}^{n-k}$, with coordinates $(x,y) \in \mathbb{R}^k \times \mathbb{R}^{n-k}$.
Let $A \subset \mathbb{R}^n$ be a Jordan measurable and bounded set, such that $A \subset (-C,C)^n$ for some $C > 0$.

We define the slices of $A$ for each $x \in (-C,C)^k$ as:
\[
A_x = \{ y \in \mathbb{R}^{n-k} \mid (x,y) \in A \} \subset (-C,C)^{n-k}
\]
We then define the upper and lower marginal densities (inner and outer measures of the slices):
\[
\overline{\varphi}(x) = \mu_{n-k}^{\text{out}}(A_x) \quad \text{and} \quad \underline{\varphi}(x) = \mu_{n-k}^{\text{in}}(A_x)
\]
Then, both functions $\overline{\varphi}, \underline{\varphi}: [-C,C]^k \to \mathbb{R}$ are Riemann integrable.
And their integrals over $[-C,C]^k$ are equal to the $n$-dimensional measure of $A$:
\[
\int_{[-C,C]^k} \overline{\varphi}(x) \, dx = \int_{[-C,C]^k} \underline{\varphi}(x) \, dx = \mu_n(A)
\]
\end{theorem}
\end{nice-box}

\begin{spoken-clean}[00:03:37 - 00:04:08]
Let us prove that. Remember that whenever we want to prove something about Jordan measures, we approximate our set from the inside and the outside.
\end{spoken-clean}

\begin{proof}[Proof of the Cavalieri's Principle for Sets]
\begin{math-stroke}[Approximation by Dyadic Cubes (Upgraded Visual)]
\begin{center}
\begin{tikzpicture}[scale=1.8]
    % Faint Background Grid to emphasize the "pixels"
    \draw[step=0.25, gray!15, very thin] (0.4, 0.4) grid (3.6, 2.1);

    % Axes
    \draw[->, thick, profgreen] (-0.2,0) -- (4,0) node[right] {$x \in \mathbb{R}^k$};
    \draw[->, thick, profgreen] (0,-0.2) -- (0,2.5) node[above] {$y \in \mathbb{R}^{n-k}$};

    % 1. The True Smooth Set A (Drawn behind the cubes)
    \draw[thick, gray, dashed, fill=gray!10] (1.1,0.6) to[out=45,in=180] (1.8,1.6) to[out=0,in=135] (3.1,1.1) to[out=-45,in=0] (2.5,0.6) to[out=180,in=-135] cycle;
    \node[gray, font=\bfseries] at (2.9, 1.2) {$A$};

    % 2. Inner Approximation F (Yellow, perfectly blocky)
    \draw[thick, profyellow!80!black, fill=profyellow!30] 
        (1.25, 0.75) -- (1.25, 1.25) -- (1.5, 1.25) -- (1.5, 1.5) -- (2.25, 1.5) -- (2.25, 1.25) -- (2.5, 1.25) -- (2.5, 1.0) -- (2.75, 1.0) -- (2.75, 0.75) -- cycle;
    \node[profyellow!80!black] at (2.0, 1.1) {$F$};

    % 3. Outer Approximation G (Red Outline, perfectly blocky)
    \draw[very thick, profred] 
        (1.0, 0.5) -- (1.0, 1.5) -- (1.25, 1.5) -- (1.25, 1.75) -- (2.5, 1.75) -- (2.5, 1.5) -- (3.0, 1.5) -- (3.0, 1.25) -- (3.25, 1.25) -- (3.25, 0.5) -- cycle;
    \node[profred] at (3.1, 1.6) {$G$};

    % 4. The vertical slice at x
    \draw[thick, profblue] (1.625, 0) -- (1.625, 2.2);
    \node[below, profblue] at (1.625, 0) {$x$};
    
    % Highlighting the trapped 1D slice bounds on the slice line
    \draw[line width=2.5pt, profred] (1.625, 0.5) -- (1.625, 1.75); % Outer slice G_x
    \draw[line width=2.5pt, gray] (1.625, 0.68) -- (1.625, 1.55); % True slice A_x (approx bound)
    \draw[line width=2.5pt, profyellow] (1.625, 0.75) -- (1.625, 1.5); % Inner slice F_x
    
    % Slice Labels
    \node[profred, right] at (1.625, 1.85) {$G_x$};
    \node[profyellow!80!black, right] at (1.625, 0.85) {$F_x$};
\end{tikzpicture}
\end{center}
\begin{explanation-of-steps}
To rigorously prove Cavalieri's Principle, we exploit the definition of Jordan measurability. We trap the smooth set $A$ between two ``pixelated'' approximations made entirely of dyadic cubes: an inner approximation $F$ (yellow) and an outer approximation $G$ (red). Because these sets are built from orthogonal blocks, taking a 1D vertical slice at a fixed $x$ yields perfectly vertical piecewise-constant line segments. The true slice $A_x$ is rigorously trapped: $F_x \subset A_x \subset G_x$.
\end{explanation-of-steps}
\end{math-stroke}

\begin{math-stroke}[Approximation by Dyadic Cubes (Original Visual)]
\begin{center}
\begin{tikzpicture}[scale=1.5]
    % Axes
    \draw[->, thick, profgreen] (-0.5,0) -- (4,0) node[right] {$x$};
    \draw[->, thick, profblue] (0,-0.5) -- (0,3) node[above] {$y$};
    \node[above right] at (3.5, 2.5) {$A$};

    % Outer approximation (red outline)
    \draw[red, thick] (0.5,0.5) rectangle (3.5,2.5); % Bounding box

    % Inner approximation (green inside) - simplified for visual clarity
    \draw[profgreen, thick] (1.0, 0.7) -- (1.0, 1.3) -- (1.5, 1.3) -- (1.5, 0.7) -- cycle;
    \draw[profgreen, thick] (1.0, 1.7) -- (1.0, 2.3) -- (1.5, 2.3) -- (1.5, 1.7) -- cycle;
    \draw[profgreen, thick] (2.0, 0.7) -- (2.0, 1.3) -- (2.5, 1.3) -- (2.5, 0.7) -- cycle;
    \draw[profgreen, thick] (2.0, 1.7) -- (2.0, 2.3) -- (2.5, 2.3) -- (2.5, 1.7) -- cycle;

    % Outer approximation of A by union of dyadic cubes (red outline)
    \draw[red, thick, smooth] plot coordinates { (0.7,0.7) (0.7,1.5) (1.2,1.5) (1.2,2.3) (1.8,2.3) (1.8,2.7) (2.7,2.7) (2.7,1.8) (3.3,1.8) (3.3,0.7) (2.5,0.7) (2.5,0.2) (1.0,0.2) (1.0,0.7) (0.7,0.7) };

    % Inner approximation of A by union of dyadic cubes (green outline)
    \draw[profgreen, thick, smooth] plot coordinates { (0.9,0.9) (0.9,1.3) (1.3,1.3) (1.3,2.1) (1.7,2.1) (1.7,2.5) (2.3,2.5) (2.3,1.7) (2.9,1.7) (2.9,0.9) (2.0,0.9) (2.0,0.5) (1.0,0.5) (1.0,0.9) (0.9,0.9) };

    % A specific slice for x
    \draw[blue, very thick] (1.7, 0.5) -- (1.7, 2.5);
    \node[below, xshift=-0.1cm, blue] at (1.7, 0) {$x$};

    % Slices of the approximations
    \draw[densely dashed, red, thin] (1.7, 0.7) -- (1.7, 2.5);
    \draw[densely dashed, profgreen, thin] (1.7, 0.9) -- (1.7, 2.3);
\end{tikzpicture}
\end{center}
\begin{explanation-of-steps}
To prove the theorem, we rely on approximating the Jordan measurable set $A$ from both the inside and the outside using unions of small dyadic cubes. The outer approximation is shown with a red outline, and the inner approximation with a green outline. When we take a slice ($x$) of these approximations, the slices of the actual set $A$ are bounded by the slices of the approximations.
\end{explanation-of-steps}
\end{math-stroke}

\begin{spoken-clean}[00:04:08 - 00:05:00]
We will take the approximation from the inside, and then there will be the approximation from the outside. Let me describe with a picture what we are going to do in the proof, and later we will write the symbols. Since our set is Jordan measurable, we can take a very small pixel size (i.e., $2^{-p}$). We can take an inner and an outer approximation of our set by dyadic cubes, so by \qt{pixels}. The difference between these two is very small. We can make the volume of the red outer approximation minus the volume of the yellow inner approximation less than $\epsilon$ for any $\epsilon > 0$ we choose (i.e., $\mu_n(G) - \mu_n(F) < \epsilon$). Now we will see what happens when we take a vertical slice of this set. If we fix $x$ and take a slice, we see that the slice of our set $A$ contains the slice of the yellow inner approximation, and is contained within the slice of the red outer approximation (i.e., $F_x \subset A_x \subset G_x$).
\end{spoken-clean}

\begin{math-stroke}[Idea: Marginal Densities for Dyadic Approximations]
The marginal densities computed for dyadic inner and outer approximations are simply dyadic step functions.
\end{math-stroke}

\begin{spoken-clean}[00:05:00 - 00:06:00]
As you start to move $x$, you will see that whenever you move $x$ by a distance less than the size of a single pixel, you do not jump to a new cube. Therefore, the marginal densities computed for the dyadic approximations are simply dyadic step functions.
\end{spoken-clean}

\begin{math-stroke}[Definition of Dyadic Cubes in $\mathbb{R}^n$]
Given a dyadic cube (``pixel'' of size $2^{-p}$) in $\mathbb{R}^n$:
\[
Q_p^n(z) = z + 2^{-p} [0,1)^n, \quad \text{where } z = (a,b), \quad a \in \mathbb{R}^k, b \in \mathbb{R}^{n-k}
\]
The full $n$-dimensional dyadic cube can be decomposed into a product of lower-dimensional dyadic cubes:
\[
Q_p^n(z) = Q_p^k(a) \times Q_p^{n-k}(b)
\]
\end{math-stroke}

\begin{spoken-clean}[00:06:00 - 00:07:10]
Let us take the size of a pixel to be $2^{-p}$ in $\mathbb{R}^n$. We will define the dyadic cube $Q_p^n(z)$ as $z + 2^{-p} [0,1)^n$. Since there will be cubes of several dimensions, I will explicitly indicate the dimension $n$ as a sub-index. Now, if $z$ is of the form $(a,b)$ where $a \in \mathbb{R}^k$ and $b \in \mathbb{R}^{n-k}$, then in this coordinate splitting, the $n$-dimensional cube $Q_p^n(z)$ is exactly equal to the Cartesian product $Q_p^k(a) \times Q_p^{n-k}(b)$.
\end{spoken-clean}

\begin{math-stroke}[Definition of Dyadic Cubes in $\mathbb{R}^m$]
Where, in general, for $m$-dimensions:
\[
Q_p^m(\bar{z}) = \bar{z} + 2^{-p} [0,1)^m \subset \mathbb{R}^m
\]
\end{math-stroke}

\begin{spoken-clean}[00:07:10 - 00:07:37]
This is something very complicated to write, but geometrically very trivial on the picture. I am simply saying that a dyadic cube in the full space is always the product of a dyadic cube of the same size on the base axis times a dyadic cube of the same size on the vertical axis.
\end{spoken-clean}

\begin{spoken-clean}[00:07:37 - 00:08:15]
Now, given $\epsilon > 0$, let $p \in \mathbb{N}$ be such that $F$ and $G$ are dyadic sets (i.e., unions of dyadic cubes of pixel size $2^{-p}$) such that $F \subset A \subset G$, and the measure of $G$ minus the measure of $F$ (i.e., $\mu_n(G) - \mu_n(F)$) is less than $\epsilon$. This is the very definition of a Jordan measurable set. We can trap it between dyadic sets, such that the outer measure minus the inner measure is less than $\epsilon$.
\end{spoken-clean}

\begin{math-stroke}[Defining Dyadic Slices]
For all $x \in [-C,C]^k$, we define the vertical slices:
\[
G_x = \{ y \in \mathbb{R}^{n-k} \mid (x,y) \in G \} \quad \text{and} \quad F_x = \{ y \in \mathbb{R}^{n-k} \mid (x,y) \in F \}
\]
\textbf{Observation:} If $Q_p^k(a) \subset \mathbb{R}^k$ is a dyadic cube of side length $2^{-p}$, then for all $x, x' \in Q_p^k(a)$, we have:
\[
F_x = F_{x'} \quad \text{and} \quad G_x = G_{x'}
\]
\end{math-stroke}

\begin{spoken-clean}[00:08:15 - 00:09:47]
Now, for all $x \in [-C,C]^k$, we define the slices $G_x$ (i.e., the set of all $y \in \mathbb{R}^{n-k}$ such that $(x,y) \in G$), and we do similarly for $F_x$. The key observation, which I was trying to convey with the drawing, is this: if $Q_p^k(a) \subset \mathbb{R}^k$ is a dyadic base cube of pixel size $2^{-p}$, then because every cube in $\mathbb{R}^n$ is a product of cubes, $F_x = F_{x'}$ for all $x, x'$ within this base cube $Q_p^k(a)$. As long as the base point remains inside this specific cube, the vertical slices cannot change; they are perfectly constant. The same logic applies to the outer approximation $G$.
\end{spoken-clean}

\begin{math-stroke}[Defining Dyadic Step Functions]
What does it mean? It means that if we define:
\[
f(x) := \mu(F_x) \quad \text{and} \quad g(x) := \mu(G_x)
\]
Then these functions are constant in each dyadic cube $Q_p(a)$, and are zero outside of $(-C,C)^k$.
\end{math-stroke}
\end{proof}


%\chapter{Real Analysis: Divergence Theorem \& Measure Theory}

\section{Recall: The Divergence Theorem}

\begin{spoken-clean}[00:00:00 - 00:02:19]
I brought back the \qt{gavel} because today I expect an \qt{extra circus}. Remember that I told you the first day, the final goal, or the \qt{final boss} of Analysis II, or one of the most representative theorems, is this Divergence Theorem. On the first day, I needed to state it very informally, using \qt{potatoes} and stuff, because you didn't know any of the mathematical elements. But now, we saw the statement again some weeks ago, and we already knew some parts. And now we know almost all of it.

Let me remind you of the statement. We know many elements now. $\Omega$ in $\mathbb{R}^n$ is an open, bounded set. Its boundary, $\partial\Omega$, will be a $C^1$ submanifold. This we know what it means. $F$ will be a $C^1$ function up to the boundary, taking values in $\mathbb{R}^n$. This is a vector field. You can imagine this function as a vector field. And the theorem says that the integral over $\Omega$ of this combination of partial derivatives (i.e., the partial derivative with respect to $x_i$ of the $i$-th component of the vector field, summed over $i$) integrated with respect to the Jordan measure (i.e., the Riemann integral) is equal to this integral over the boundary of the domain with respect to the surface measure of the boundary. We know that $F \cdot \nu$ is $F$ multiplied by the normal vector, which we already saw what a tangent vector and a normal vector to the boundary are. So, more or less everything you know what it is now, except for one thing. This (i.e., the integral over the boundary) goes with this (i.e., the surface measure $dS$). The integral over a surface of the differential of volume in the surface. We know how to integrate on the full space, but not over a curved surface. And this is where we are going today. We are going to learn, we are going to define what is the integral over the surface.
\end{spoken-clean}

\begin{math-stroke}[The Divergence Theorem]
\begin{nice-box}[Divergence Theorem]
\begin{theorem}[Divergence Theorem]
Let $\Omega \subset \mathbb{R}^n$ be an open, bounded set, and let $\partial\Omega$ be a $C^1$ submanifold.
Let $F \in C^1(\overline{\Omega}, \mathbb{R}^n)$ be a vector field. Then:
\[
\int_{\Omega} \sum_{i=1}^n \partial_i F_i \, dx = \int_{\partial\Omega} F \cdot \underbrace{\nu}_{\text{unit vector normal to } \partial\Omega} \, \underbrace{dS}_{\mathclap{\substack{\text{Surface measure} \\ \text{(Goal of today's lecture)}}}}
\]
\end{theorem}
\end{nice-box}
\end{math-stroke}

\section{Measuring Volumes over Submanifolds}

\begin{spoken-clean}[02:19 - 04:30]
Okay, so the next step in our quest is to define integrals over d-dimensional submanifolds. Today, we will measure volumes over submanifolds of dimension $d$. Towards this goal, we start with the easiest case. What is the easiest possible submanifold? It's the submanifold of dimension one. So, let's start by $d=1$. A submanifold of dimension one, we call it a curve. And we will see how to measure length.
\end{spoken-clean}

\begin{math-stroke}[Quest Log: From Volumes to Lengths]
The lecturer outlines the plan for the day.
\begin{description}
    \item[Next Step:] Define $\int$ over $d$-dim submanifolds.
    \item[Today's Goal:] Define how to measure volumes over submanifolds.
    \item[Starting Case:] Begin with the simplest case, $d=1$, where a submanifold is a curve and its "volume" is its length.
\end{description}
\end{math-stroke}

\begin{spoken-clean}[04:30 - 06:25]
I will work in the set of parameterized submanifolds so far, and I will give you the definition. Let me add here, over parameterized submanifolds. For today, then we will extend, but for today, let me add \emph{parameterized}. Okay.

So, a parameterized 1-manifold is a curve. Remember that we gave the definition some weeks ago. We have some $\phi$ (i.e., a map $\phi$) from $V$ to $\mathbb{R}^n$, where $V \subset \mathbb{R}$ is open. That is the parametrization. And this has to satisfy some conditions. The most important was that it is injective, and also that the differential (which in this case is just the velocity vector, $\phi'(t)$) is injective. In the case of curves, it means that this number, the velocity vector, is never zero.
\end{spoken-clean}

\begin{math-stroke}[Definition: Parameterized 1-Manifold (Curve)]
A \textbf{parameterized 1-manifold}, or curve, is defined by a map $\phi: V \to \mathbb{R}^n$, where $V \subset \mathbb{R}$ is an open set. The map must satisfy several conditions:
\begin{itemize}
    \item It must be \textbf{injective}.
    \item Its derivative must never be zero: $\phi'(t) \neq 0$.
    \item It must satisfy a topological property to prevent self-intersections where the curve loops back on itself, which is a more advanced condition from differential geometry.
\end{itemize}
\end{math-stroke}

\begin{spoken-clean}[06:25 - 08:38]
And now we will give the definition of length. So, if in this setup, if $[a,b]$ is a closed interval contained in the domain of the parametrization, then I define the length of $\phi([a,b])$ as the integral between $a$ and $b$ of the speed. This is the velocity vector, I take the modulus, and then I integrate with respect to $t$. This is the definition of length.

So let me do maybe a bit of a picture of what this situation is. So we have our $V$ here. So this will be $\mathbb{R}$, and we have our $V$. It could be an open interval or some open set. And then we have our map, that is $\phi$. And then there is the $\phi(V)$, that is some curve. And then we are taking a subset that is a closed subset here, that is $[a,b]$. And this piece is defining a closed piece of the curve. And then we are saying that by definition, the length of this red piece is this integral over there. Now I ask you, is this definition sound?
\end{spoken-clean}

\begin{math-stroke}[Definition of Arc Length]
\begin{definition}[Length of a Curve]
If $[a,b] \subset V$ is a closed interval within the domain of the parametrization $\phi$, the length of the curve segment $\phi([a,b])$ is defined as:
\[
L(\phi([a,b])) = \int_a^b |\phi'(t)| \, dt
\]
\end{definition}
\begin{center}
\begin{tikzpicture}
    % Domain V in R
    \draw[->, thick] (-1,0) -- (3,0) node[right] {$\mathbb{R}$};
    \draw[thick, blue] (0,0) -- (2,0);
    \node[below] at (1,0) {$V$};
    \draw[thick, BrickRed] (0.5,0) -- (1.5,0);
    \node[below] at (0.5,0) {$a$};
    \node[below] at (1.5,0) {$b$};

    % Mapping phi
    \draw[->, thick] (1, 0.5) to[bend left] node[midway, above] {$\phi$} (4, 1.5);

    % Image curve in R^n
    \draw[thick, blue] (3,1) to[out=30, in=150] (6,2);
    \node[above] at (5.5, 2) {$\phi(V)$};
    \draw[thick, BrickRed, line width=1.5pt] (3.5, 1.18) to[out=30, in=160] (5, 1.8);
    \node[above] at (4.25, 2) {$\phi([a,b])$};
\end{tikzpicture}
\end{center}
\end{math-stroke}

\begin{spoken-clean}[08:38 - 11:46]
What properties must length satisfy? And there will be four. But maybe you can tell me two, and I will complete the other two. So what should I check to see that this definition is sound? Two properties. So think that I'm defining the length of a curve, like this wire over here. What I'm saying is that to define the length of the curve, I am choosing a parametrization of the curve, I am running the curve through some map, and then I'm defining the length as this integral. Is this sound? Is this correct? What would you like to ask this definition to satisfy so that it is correct?
\end{spoken-clean}

\begin{student-question}[Student Answer]
Well, I think it should be positive, first of all. I wouldn't want to have a negative length or something like that. And I mean, this obviously satisfies that.
\end{student-question}

\begin{spoken-clean}[10:20 - 10:22]
It's positive. Okay.
\end{spoken-clean}

\begin{student-question}[Student Answer]
I also think it should be independent of $\phi$, because it should depend only on the object you measure, not on how we measure it, so to say.
\end{student-question}

\begin{spoken-clean}[10:22 - 10:23]
Louder.
\end{spoken-clean}

\begin{student-question}[Student Answer]
You should, like, the length should have additivity as one of its properties, meaning you could add the distance of two paths.
\end{student-question}

\begin{spoken-clean}[10:23 - 11:46]
So, additivity. Like the measure, right? Yes, additivity. Anything that is measuring something, if you... So additivity means that if I cut into this curve, the length of this piece plus the length of this piece is the length of the total curve. If this is not happening, the definition of length must be ill-posed, right? So we need to check that. But the additivity, since it is given by an integral, is usually very easy, right? Because integrals are additive. Okay. Some other property? I asked two, so you already did your work, but do you want to say some other property? There.
\end{spoken-clean}

\begin{student-question}[Student Answer]
It would be nice if if we have like the curve just in $\mathbb{R}$, if then it corresponds to just what we would intuitively guess.
\end{student-question}

\begin{spoken-clean}[11:46 - 13:55]
Yeah. So we can call it normalization. So he's saying that if the curve is in $\mathbb{R}$, meaning like, I guess if the curve is for instance $[0,1] \times \{0\}$, right? And this is the zero of $\mathbb{R}^{n-1}$. So the curve is just a straight segment along the $x_1$ axis, then the length of this curve should be one, right? Okay. And the last one, I will give you a hint. My curve is put in the space in some position, but this is a rigid curve, that's why I made it of thick wire. So if I put this curve in this position, or in this other position, should the length be the same? Yeah. So this is the invariance under rigid motion.
\end{spoken-clean}

\begin{math-stroke}[Axioms of Length]
For a definition of length to be sound, it must satisfy four fundamental properties:
\begin{enumerate}
    \item \textbf{Independence of Parametrization:} The length must depend only on the geometric path, not the function used to trace it.
    \item \textbf{Additivity:} If a curve is split into two segments, the total length is the sum of the segment lengths.
    \item \textbf{Invariance under Euclidean Isometry:} Rigid motions (translations and rotations) of the curve in space should not change its length.
    \item \textbf{Normalization:} The length of the standard unit interval $[0,1]$ in $\mathbb{R}$ must be 1.
\end{enumerate}
\end{math-stroke}

\begin{spoken-clean}[13:55 - 15:07]
So what is Euclidean isometry? What types of maps preserve rigid motions of the space? The rigid motions of space are translations and rotations, or compositions, right? So $F$ is a rigid motion, or a Euclidean isometry, if it is of the form $F(x) = Rx+b$, where $b$ is a vector (i.e., the translation vector) and $R$ is orthogonal. So if you want, a rotation or a symmetry or a composition of them.
\end{spoken-clean}

\begin{math-stroke}[Euclidean Isometry]
A map $F: \mathbb{R}^n \to \mathbb{R}^n$ is a \textbf{Euclidean isometry} (or rigid motion) if it has the form:
\[
F(x) = Rx+b
\]
where $R$ is an orthogonal matrix ($R^T R = I$) and $b \in \mathbb{R}^n$ is a translation vector.
\end{math-stroke}

\begin{spoken-clean}[15:07 - 17:52]
So if you take $\Psi$ from $[c,d]$ to $[a,b]$ that is bijective, $C^1$, $C^1$ inverse. Then a reparametrization is this: you would take... Let me do the drawing here. So here, you are taking some $\Psi$, and the new interval now is $[c,d]$, that maybe is longer, or maybe shorter, or maybe the speed here is irregular. Here maybe the speed is one and here the speed is irregular, right? But this map is going from here to here. Then the image of the composition of the maps is the same. So the length should be the same. And this is translated in my definition of length that I would like that when I apply this integral between $c$ and $d$ of the composition (i.e., $(\phi \circ \Psi)'(s)$) $ds$, it should be equal to the original integral (i.e., $\int_a^b |\phi'(t)| \, dt$).
This we will prove in a second, a more general property. So I think this is a good... I put this property over here as an exercise. It will follow from the more general property I will give for general $d$, right? But it's very good that you try to do this exercise by yourself. It's one line of computation if you do it right. Okay. So you need to apply the chain rule. Okay, so the additivity...
\end{spoken-clean}

\begin{math-stroke}[Invariance under Reparametrization]
Let $\Psi: [c,d] \to [a,b]$ be a bijective $C^1$ map with a $C^1$ inverse (a diffeomorphism). The reparametrized curve is $\tilde{\phi} = \phi \circ \Psi$. We must prove that $L(\tilde{\phi}) = L(\phi)$.
\[
\int_c^d |(\phi \circ \Psi)'(s)| \, ds = \int_a^b |\phi'(t)| \, dt
\]
This is left as an exercise and follows directly from the chain rule and 1D integration by substitution.
\end{math-stroke}

\begin{spoken-clean}[17:52 - 18:10]
The additivity is what I said, right? So this is property number one, property number two, the additivity is that when you have the length of $\phi([a,b])$ should be equal to the length of $\phi([a,c])$ plus the length of $\phi([c,b])$ if $c$ is any point between $a$ and $b$, for example. This is the additivity. But this follows from the additivity of the integral. So here there is nothing to prove.
\end{spoken-clean}

\end{document}